% Options for packages loaded elsewhere
\PassOptionsToPackage{unicode}{hyperref}
\PassOptionsToPackage{hyphens}{url}
\PassOptionsToPackage{dvipsnames,svgnames,x11names}{xcolor}
%
\documentclass[
  11pt,
]{article}
\usepackage{amsmath,amssymb}
\usepackage{iftex}
\ifPDFTeX
  \usepackage[T1]{fontenc}
  \usepackage[utf8]{inputenc}
  \usepackage{textcomp} % provide euro and other symbols
\else % if luatex or xetex
  \usepackage{unicode-math} % this also loads fontspec
  \defaultfontfeatures{Scale=MatchLowercase}
  \defaultfontfeatures[\rmfamily]{Ligatures=TeX,Scale=1}
\fi
\usepackage{lmodern}
\ifPDFTeX\else
  % xetex/luatex font selection
  \setmainfont[]{Liberation Serif}
  \setsansfont[]{Liberation Sans}
  \setmonofont[]{DejaVu Sans Mono}
\fi
% Use upquote if available, for straight quotes in verbatim environments
\IfFileExists{upquote.sty}{\usepackage{upquote}}{}
\IfFileExists{microtype.sty}{% use microtype if available
  \usepackage[]{microtype}
  \UseMicrotypeSet[protrusion]{basicmath} % disable protrusion for tt fonts
}{}
\makeatletter
\@ifundefined{KOMAClassName}{% if non-KOMA class
  \IfFileExists{parskip.sty}{%
    \usepackage{parskip}
  }{% else
    \setlength{\parindent}{0pt}
    \setlength{\parskip}{6pt plus 2pt minus 1pt}}
}{% if KOMA class
  \KOMAoptions{parskip=half}}
\makeatother
\usepackage{xcolor}
\usepackage[margin=0.8in]{geometry}
\usepackage{longtable,booktabs,array}
\usepackage{calc} % for calculating minipage widths
% Correct order of tables after \paragraph or \subparagraph
\usepackage{etoolbox}
\makeatletter
\patchcmd\longtable{\par}{\if@noskipsec\mbox{}\fi\par}{}{}
\makeatother
% Allow footnotes in longtable head/foot
\IfFileExists{footnotehyper.sty}{\usepackage{footnotehyper}}{\usepackage{footnote}}
\makesavenoteenv{longtable}
\setlength{\emergencystretch}{3em} % prevent overfull lines
\providecommand{\tightlist}{%
  \setlength{\itemsep}{0pt}\setlength{\parskip}{0pt}}
\setcounter{secnumdepth}{-\maxdimen} % remove section numbering
\ifLuaTeX
\usepackage[bidi=basic]{babel}
\else
\usepackage[bidi=default]{babel}
\fi
\babelprovide[main,import]{american}
\ifPDFTeX
\else
\babelfont{rm}[]{Liberation Serif}
\fi
% get rid of language-specific shorthands (see #6817):
\let\LanguageShortHands\languageshorthands
\def\languageshorthands#1{}
\usepackage{adjustbox}
\emergencystretch=3em
\AtBeginEnvironment{longtable}{\small\sloppy}
\usepackage{stmaryrd}
% Printing shims for run-on conditional-bar tokens in the supplied Biology source.
\providecommand{\midL}{\mid L}
\providecommand{\midN}{\mid N}
\sloppy
\ifLuaTeX
  \usepackage{selnolig}  % disable illegal ligatures
\fi
\usepackage{bookmark}
\IfFileExists{xurl.sty}{\usepackage{xurl}}{} % add URL line breaks if available
\urlstyle{same}
\hypersetup{
  pdflang={en-US},
  colorlinks=true,
  linkcolor={blue},
  filecolor={Maroon},
  citecolor={Blue},
  urlcolor={blue},
  pdfcreator={LaTeX via pandoc}}

\author{}
\date{}

\begin{document}

\section{COHERENCE UNIVERSALISM:
BIOLOGY}\label{coherence-universalism-biology}

\subsection{Life as Active Self-Maintenance: Viability, Individuation,
and Multiscale
Regulation}\label{life-as-active-self-maintenance-viability-individuation-and-multiscale-regulation}

\textbf{Gaura Kiśora Dās Rader}\\
\textbf{Heaven≡Earth Foundation}\\
\textbf{White Paper Series}

\textbf{Version 0.10}\\
\textbf{October 2026}

\textbf{Release status.} This manuscript is part of the complete Version
0.10 proposal. Numbered principles identify the framework's current
formulations; their claim status, external validation, and authorization
for consequential application remain separate questions.

\textbf{Freeze status.} This manuscript is synchronized to the frozen
Version 0.10 release registry, evidence package, and review gates.
Future changes require a new versioned release.

\begin{center}\rule{0.5\linewidth}{0.5pt}\end{center}

\section{Abstract}\label{abstract}

Biology concerns systems that remain organized while their materials,
environments, and internal states continually change. Coherence
Universalism describes this domain as \textbf{active self-maintenance
under constraint}, but that description becomes scientifically useful
only when the system boundary, protected variables, regulatory
organization, perturbation class, and temporal horizon are specified.

This paper develops a non-scalar formal architecture for doing so. A
candidate biological organization is represented by a state space,
environment, dynamics, admissible regulatory actions, relational
organization, viability constraints, and history. Viability is treated
as a trajectory property rather than a point score. The central
mathematical objects are a viability region and its finite- or
indefinite-horizon viability kernel; a robust viability kernel when
disturbances are included; interventionally estimated maintenance
contributions; and a \textbf{biological coherence profile} recording
viability margin, resilience, causal integration, organizational
closure, adaptive reach, historical continuity, and energetic or
material cost. These dimensions need not move together and are not
assumed to admit one universal scalar aggregation.

The paper proposes that a physical organization becomes a candidate
biological individual when a relatively bounded network of processes
causally contributes to preserving and regenerating the conditions of
its own continuation. This yields a naturalized sense of
organism-relative stake: some states are better or worse \emph{for the
organization} because crossing particular boundaries destroys or
seriously degrades the organization whose regulatory processes are
arranged around those boundaries. Such stake is teleonomic and
functional. It does not by itself establish consciousness, felt valence,
moral patiency, or reflective agency.

On this basis, the origin of life is analyzed as a family of coupled
autonomy transitions rather than one universal phase transition.
Autocatalysis, compartmentalization, energy transduction, constraint
closure, regulation, heredity, and open-ended evolution are separable
achievements. Contemporary work on autocatalytic sets, artificial cells,
synthetic protocells, and prebiotic regulatory models provides candidate
systems in which these dimensions can be manipulated independently. A
recurrence-dominance criterion is shown to be insufficient because
passive stable structures can satisfy it without producing, repairing,
or regulating the conditions of their own persistence.

Development and regeneration are treated as multiscale control problems.
Gene regulation, biochemical signaling, mechanics, bioelectric coupling,
and tissue history jointly constrain trajectories toward viable forms.
Attractor and target-morphology language is warranted when it generates
interventionally discriminating predictions; it is not a license to
posit a hidden blueprint or a universal organism-level optimizer.
Planarian bioelectric experiments, engineered multicellular constructs,
and recent multiscale models support a research program in which
long-range coordination can be perturbed independently of local cellular
competence.

Biological agency is presented as graded. Homeostatic regulation,
sensorimotor control, adaptive reconfiguration, learning, and
model-based anticipation are distinct capacities. Single cells can
instantiate some of them without implying human-like cognition or
consciousness. Evolution is treated as a stochastic and selective
exploration of \textbf{structured evolutionary possibility}: mutation
biases, genotype--phenotype mappings, developmental dynamics, physical
feasibility, ecology, inheritance, drift, and prior history make some
transformations far more reachable than others. This structure can
contribute to convergent evolution without imposing one endpoint.
Biological coherence is neither identical with reproductive fitness nor
guaranteed to increase over evolutionary time. Major evolutionary
transitions are analyzed as changes in individuality involving
cooperation, dependence, division of labor, conflict mediation, and the
relocation of heritable fitness.

The resulting framework makes a bounded claim. CU Biology does not
derive life from physics, define all life by one number, or show that
the universe aims at life. It offers a governed set of questions and
formalisms for comparing how candidate living systems maintain
themselves, where their boundaries lie, what disturbances they can
survive, how they reorganize, and which costs their persistence exports.
Its scientific value will depend on whether these variables improve
explanation and prediction beyond established models. Its interpretive
value is that life can be understood---without violating
thermodynamics---as matter organized so that the continuation of an
organization becomes causally relevant to what the organization does
next. Its commitment register receives living beings as sacred
historical achievements and the biosphere as the living field within
which every organismal life becomes possible.

\begin{center}\rule{0.5\linewidth}{0.5pt}\end{center}

\textbf{Keywords:} active self-maintenance; viability theory; biological
individuality; autonomy; morphogenesis; regeneration; basal cognition;
evolution; organism-relative stake; multiscale regulation

\begin{center}\rule{0.5\linewidth}{0.5pt}\end{center}

\subsection{Relation to the series}\label{relation-to-the-series}

This paper can be read independently. It defines the biological concepts
needed for its own argument. \emph{Metaphysics and Epistemology}
develops the framework's account of coherence, truth, bridge principles,
and interpretive commitment; \emph{Foundations} records the shared
principle architecture; \emph{Physics} supplies the physical and
thermodynamic constraints within which biological organization exists.
The later papers examine consciousness, psychology, social dynamics,
ethics, and artificial intelligence. Formal proofs, full claim ledgers,
illustrative-model details, audits, and reproducibility materials are
preserved in the
\href{../supplements/04_Biology_Technical_Supplement_v0.10.md}{technical
supplement}.

\subsection{Note on Version 0.10}\label{note-on-version-0.10}

Version 0.10 preserves the original integrative question---what connects
origin, development, maintenance, evolution, and biological
failure---while replacing a universal scalar, a universal biological
gradient, one sharp origin threshold, and validation claims based on toy
models with a non-scalar, interventionally governed architecture. The
main text presents the current theory directly; Appendix B and the
technical archive preserve the detailed revision history.\footnote{The
  March 2026 paper introduced the scalar \(C=A I\), a universal
  biological gradient, a single origin threshold, and stronger claims
  about entropy, evolution, cancer, aging, and death. The technical
  archive preserves that manuscript and the complete revision record.}

\textbf{Release-candidate synchronization.} This manuscript uses the
shared Version 0.10 definitions, principle identifiers, Mandala
cross-references, and interpretive and commitment language published in
\emph{Coherence Universalism: Metaphysics and Epistemology} and
\emph{Coherence Universalism: Foundations}. Its audited technical core,
formal results, numerical results, and claim-status boundaries remain
unchanged except where this manuscript explicitly states otherwise.

\subsection{How claims are marked}\label{how-claims-are-marked}

Definitions, formal results, hypotheses, models, interpretations, and
open questions are labeled where their status is load-bearing. Imported
empirical results are presented through their citations and scope
conditions. Biological bridges do not follow from Physics alone, and
spiritual interpretations are constrained by the scientific account
without being entailed by it.

\section{Contents}\label{contents}

\begin{enumerate}
\def\labelenumi{\arabic{enumi}.}
\tightlist
\item
  Biology as Active Self-Maintenance\\
\item
  Conceptual Discipline: Organization, Constraint, and Function\\
\item
  A Formal Architecture for Biological Coherence\\
\item
  Biological Individuality and Organism-Relative Stake\\
\item
  Thermodynamic Grounding\\
\item
  Origins of Life as an Autonomy Transition\\
\item
  Development, Morphogenesis, and Regeneration\\
\item
  Biological Agency, Learning, and Minimal Epistemology\\
\item
  Evolution, Constraint, and Major Transitions\\
\item
  Pathology, Cancer, Aging, and Death\\
\item
  Synthetic and Artificial Contrast Cases\\
\item
  Empirical Program, Discriminating Predictions, and Failure
  Conditions\\
\item
  What Commitments This Account Licenses
\end{enumerate}

Appendix A. Biology Principles and Claim-Governance Summary\\
Appendix B. Major Revisions from v0.09\\
Appendix C. Technical Supplement and Reproducibility\\
References

\section{Principal contributions}\label{principal-contributions}

\begin{enumerate}
\def\labelenumi{\arabic{enumi}.}
\tightlist
\item
  It defines biological coherence as a domain-specific assessment of
  active self-maintenance rather than as a universal substance or score.
\item
  It replaces the universal scalar \(C=A I\) with a structured profile
  and explicit scalarization rules.
\item
  It makes viability regions, viability kernels, disturbance classes,
  control sets, boundaries, and horizons explicit.
\item
  It distinguishes passive persistence, externally maintained stability,
  self-maintenance, reproduction, adaptation, and open-ended evolution.
\item
  It gives an intervention-based criterion for endogenous maintenance.
\item
  It reconstructs biological individuality as multidimensional and
  potentially nested.
\item
  It defines organism-relative stake without equating it with phenomenal
  valence or moral status.
\item
  It replaces a single abiogenesis threshold with a family of autonomy
  transitions.
\item
  It narrows morphogenetic global-constraint claims to testable
  multiscale-control hypotheses.
\item
  It presents agency and cognition as graded capacities and separates
  them from consciousness.
\item
  It distinguishes biological organization from fitness and states when
  selection can stabilize regulatory architectures.
\item
  It converts broad predictions into comparative, interventionally
  discriminating research programs.
\item
  It shows through two deterministic constructions that the viability
  effect of one diffusive coupling law can reverse when floors, state
  distribution, and disturbances change.
\item
  It requires every biological coherence claim to declare its unit,
  horizon, disturbance class, scalarization choices, and register
  boundary.
\item
  It distinguishes designer optimization, as-if control analysis, and
  evolutionary selection.
\item
  It reconstructs evolution as exploration of structured evolutionary
  possibility rather than as either uniform random search or universal
  coherence ascent.
\item
  It places the spiritual meaning of life in a designated ``What
  Commitments This Account Licenses'' section rather than presenting it
  as a biological derivation.
\end{enumerate}

\begin{center}\rule{0.5\linewidth}{0.5pt}\end{center}

\section{Domain Analytic Frame and CAP
Record}\label{domain-analytic-frame-and-cap-record}

\textbf{Target:} Coherence Universalism Version 0.10 release\\
\textbf{Status:} Release declaration; scope-limited and revisable under
the declared review gates\\
\textbf{Use rule:} Unresolved fields are left explicit rather than
completed by analogy.

\subsection{Paper and model identity}\label{paper-and-model-identity}

\begin{longtable}[]{@{}
  >{\raggedright\arraybackslash}p{(\columnwidth - 4\tabcolsep) * \real{0.3333}}
  >{\raggedright\arraybackslash}p{(\columnwidth - 4\tabcolsep) * \real{0.3333}}
  >{\raggedright\arraybackslash}p{(\columnwidth - 4\tabcolsep) * \real{0.3333}}@{}}
\toprule\noalign{}
\begin{minipage}[b]{\linewidth}\raggedright
Field
\end{minipage} & \begin{minipage}[b]{\linewidth}\raggedright
Declaration
\end{minipage} & \begin{minipage}[b]{\linewidth}\raggedright
Status
\end{minipage} \\
\midrule\noalign{}
\endhead
\bottomrule\noalign{}
\endlastfoot
Paper & Coherence Universalism: Biology & Specified \\
Domain owner & Biology & Specified by the Version 0.10 release corpus \\
Primary model or analysis & Candidate biological system
\(\mathcal B=(X,E,U,F,K,G,R,H,\pi,T)\) & Paper-relative \\
Primary claim status & Definitions, imported viability and control
formalisms, empirical generalizations, bridge principles, hypotheses,
and commitments & Mixed statuses remain explicit \\
Review gate & Biology, origin-of-life, systems biology, control and
viability theory, philosophy of biology, and relevant ethical review &
Proportionate to claim and consequential reach \\
Methodological authority & CU-FP2, CU-M2, CU-M11, CU-M12, CU-M13,
CU-M14, and CU-M15 & Controlling Version 0.10 methodological
framework \\
\end{longtable}

\subsection{Shared Analytic Frame}\label{shared-analytic-frame}

\begin{longtable}[]{@{}
  >{\raggedright\arraybackslash}p{(\columnwidth - 4\tabcolsep) * \real{0.3333}}
  >{\raggedright\arraybackslash}p{(\columnwidth - 4\tabcolsep) * \real{0.3333}}
  >{\raggedright\arraybackslash}p{(\columnwidth - 4\tabcolsep) * \real{0.3333}}@{}}
\toprule\noalign{}
\begin{minipage}[b]{\linewidth}\raggedright
\#
\end{minipage} & \begin{minipage}[b]{\linewidth}\raggedright
Slot and domain declaration
\end{minipage} & \begin{minipage}[b]{\linewidth}\raggedright
Status and evidence
\end{minipage} \\
\midrule\noalign{}
\endhead
\bottomrule\noalign{}
\endlastfoot
1 & \textbf{Unit and boundary.} A declared candidate biological
organization, which may be a cell, organism, tissue, lineage,
consortium, or other bounded unit. & Boundary-sensitive \\
2 & \textbf{State or trajectory space.} \(X\) with history \(H\);
viability is primarily a trajectory property rather than a point score.
& Specified by model \\
3 & \textbf{Environment.} \(E\) contains material, energetic,
ecological, developmental, and social conditions outside the focal
boundary. & Specified per application \\
4 & \textbf{Constraints and viable region.} \(K\) contains physical,
viability, regulatory, generative, historical, and observer-imposed
constraints; \(V\), viability kernels, and robust kernels are central
objects. & Formal architecture specified \\
5 & \textbf{Actions, options, and transitions.} \(U\) contains
regulatory actions and interventions; identity-preserving, adaptive,
reproductive, and failure transitions are distinguished. &
Model-specific \\
6 & \textbf{Dynamics or update rule.} \(F\) may be stochastic, discrete,
hybrid, reaction-network, field, or agent-based; no universal biological
gradient is asserted. & Plural \\
7 & \textbf{Observation and access.} \(\pi\), sensors, assays,
interventions, and causal comparisons connect latent organization to
evidence. & Observation model must be declared \\
8 & \textbf{Perturbation and intervention class.} Disturbance class
\(\mathcal D\), control set, counterfactual disabling, and repair
interventions define robustness and maintenance claims. & Required \\
9 & \textbf{Equilibria, basins, or regimes.} Attractors and target sets
are permitted only in models with demonstrated dynamics; viability need
not be an attractor and an attractor need not be viable. &
Conditional \\
10 & \textbf{Local--global status.} Primarily coordinate-wise and
qualitative across cell, tissue, organism, lineage, and ecosystem
scales; exact relations belong only to declared submodels. & No
whole-profile decomposition \\
11 & \textbf{Scalar status.} The biological coherence object is a
profile; coordinate measures and decision scalarizations are permitted
after floor screening. & No canonical scalar \\
12 & \textbf{Statistical-image and allocation status.} Population
frequencies, stochastic state distributions, occupancy measures, and
some resource budgets may qualify. The biological coherence profile is
not a simplex. & Candidate subobjects only \\
13 & \textbf{Translation status.} Viability, control, organizational
closure, information, and thermodynamics enter through separate licenses
rather than one hidden equation. & Levels 1--4 by claim \\
14 & \textbf{Open gaps.} Boundary individuation, operational closure,
causal integration, cost accounting, origin transitions, and comparative
validation remain incomplete. & Research program \\
\end{longtable}

\subsection{Compact Coherence Analysis Protocol
record}\label{compact-coherence-analysis-protocol-record}

\begin{longtable}[]{@{}
  >{\raggedright\arraybackslash}p{(\columnwidth - 4\tabcolsep) * \real{0.3333}}
  >{\raggedright\arraybackslash}p{(\columnwidth - 4\tabcolsep) * \real{0.3333}}
  >{\raggedright\arraybackslash}p{(\columnwidth - 4\tabcolsep) * \real{0.3333}}@{}}
\toprule\noalign{}
\begin{minipage}[b]{\linewidth}\raggedright
CAP module
\end{minipage} & \begin{minipage}[b]{\linewidth}\raggedright
Declaration
\end{minipage} & \begin{minipage}[b]{\linewidth}\raggedright
Boundary
\end{minipage} \\
\midrule\noalign{}
\endhead
\bottomrule\noalign{}
\endlastfoot
System declaration & Name \(B,\pi,T,\mathcal D,\mathcal U\), identity
criterion, state variables, environment, and regulatory relations. &
Already required in Sections 2--3 \\
Observation declaration & State assays, proxies, causal interventions,
missing scales, and observer-imposed categories. & Correlation is
insufficient for integration \\
Evaluation declaration & Declare viability floors, profile coordinates,
costs, disturbances, scalarization, and uncertainty. & Viability-first
screening is often appropriate \\
Power and authorship declaration & Name who fixes the organism boundary,
function, intervention set, cost boundary, and ethical use of biological
classifications. & Especially important in synthetic and human
applications \\
Revision declaration & Boundary, horizon, disturbance class, identity
criterion, or function changes constitute model revision. & Must be
versioned \\
\end{longtable}

\subsection{License summary}\label{license-summary}

\begin{longtable}[]{@{}
  >{\raggedright\arraybackslash}p{(\columnwidth - 4\tabcolsep) * \real{0.3333}}
  >{\raggedright\arraybackslash}p{(\columnwidth - 4\tabcolsep) * \real{0.3333}}
  >{\raggedright\arraybackslash}p{(\columnwidth - 4\tabcolsep) * \real{0.3333}}@{}}
\toprule\noalign{}
\begin{minipage}[b]{\linewidth}\raggedright
License
\end{minipage} & \begin{minipage}[b]{\linewidth}\raggedright
Declaration
\end{minipage} & \begin{minipage}[b]{\linewidth}\raggedright
Current status
\end{minipage} \\
\midrule\noalign{}
\endhead
\bottomrule\noalign{}
\endlastfoot
Scalar license & Coordinate measures and purpose-bound decision
scalarizations only; the former product AI is a permitted local model,
not the definition. & Conditional \\
Local--global license & Coordinate-wise or qualitative unless an exact
relation is proved for a declared quantity and partition. & Default \\
Statistical Image License & Possible for type frequencies, transition
distributions, stochastic state models, and occupancy measures. &
Object-specific \\
Allocation License & Possible for defended metabolic, developmental,
energetic, or material budgets with a real total and no double counting.
& Empirical hypothesis \\
Refinement and Coarse-Graining License & Permitted only when biological
categories retain meaning under aggregation and support changes are
handled. & Conditional \\
Fisher-Rao License & Possible for licensed frequency or probability
models; not for the multidimensional biological profile as a whole. &
Object-specific only \\
Translation license & Imported viability and control results can reach
Level 3 inside their mathematical class; biological interpretation and
prediction require further evidence. & Mixed \\
\end{longtable}

\subsection{Open gaps and review
disposition}\label{open-gaps-and-review-disposition}

\begin{longtable}[]{@{}
  >{\raggedright\arraybackslash}p{(\columnwidth - 4\tabcolsep) * \real{0.3333}}
  >{\raggedright\arraybackslash}p{(\columnwidth - 4\tabcolsep) * \real{0.3333}}
  >{\raggedright\arraybackslash}p{(\columnwidth - 4\tabcolsep) * \real{0.3333}}@{}}
\toprule\noalign{}
\begin{minipage}[b]{\linewidth}\raggedright
Gap
\end{minipage} & \begin{minipage}[b]{\linewidth}\raggedright
Consequence
\end{minipage} & \begin{minipage}[b]{\linewidth}\raggedright
Required work
\end{minipage} \\
\midrule\noalign{}
\endhead
\bottomrule\noalign{}
\endlastfoot
Biological individuality & Different boundaries can reverse the
assessment. & Interventionally test boundary and causal closure \\
Profile measurement & Coordinates lack universal units and validation. &
Develop construct-specific protocols \\
Cross-scale cost & Local maintenance can export ecological or
organism-level damage. & Disaggregate receivers and horizons \\
\end{longtable}

\begin{quote}
\textbf{Propagation rule.} This card makes the paper's existing objects
and limits visible. It does not strengthen any empirical, mathematical,
normative, metaphysical, or spiritual claim beyond the wording supported
in the manuscript.
\end{quote}

\begin{center}\rule{0.5\linewidth}{0.5pt}\end{center}

\section{Section 1. Biology as active
self-maintenance}\label{section-1.-biology-as-active-self-maintenance}

\subsection{1.1 The integrative problem}\label{the-integrative-problem}

Biology studies entities that persist by changing. A cell replaces
constituents while maintaining a membrane, metabolism, and regulatory
organization. An organism develops through states that differ radically
in material composition and form. A lineage persists by producing
non-identical descendants. A tissue can recover from damage by
rebuilding relations that no individual cell contains as a complete
instruction.

These phenomena are investigated by mature disciplines: molecular
biology, physiology, developmental biology, ecology, evolutionary
biology, systems biology, control theory, and origin-of-life research.
The problem is not that biology lacks mechanisms. It is that the
language of \emph{organization} appears across these mechanisms without
always receiving a common operational grammar.

Coherence Universalism uses \textbf{coherence} in this domain to name
integrated relational organization under declared constraints. In
biology, the relevant organization is not static order, material
identity, or one-time assembly. Living systems maintain changing sets of
mutually dependent conditions across scales and time. The scientific
task is therefore not to assign every organism one score, but to specify
the boundary, viability conditions, maintenance processes,
perturbations, costs, and horizon that make a claim of biological
organization testable.

\subsection{1.2 The bounded central
claim}\label{the-bounded-central-claim}

\begin{quote}
\textbf{Bridge principle.} Living systems are materially and
energetically open organizations in which internally generated processes
make a counterfactual difference to the continued maintenance, repair,
reproduction, or adaptive transformation of the organization that
generates them.
\end{quote}

This is a proposal about the characteristic organization of life, not a
universally accepted necessary-and-sufficient definition. It joins four
ideas:

\begin{enumerate}
\def\labelenumi{\arabic{enumi}.}
\tightlist
\item
  \textbf{Openness:} matter and energy cross the boundary.
\item
  \textbf{Relative closure:} a network of constraints or regulatory
  processes depends on and contributes to maintaining other members of
  the network.
\item
  \textbf{Viability:} some trajectories preserve the organization while
  others cross failure boundaries.
\item
  \textbf{Endogenous contribution:} the system's own organization
  causally enlarges the set or duration of viable trajectories.
\end{enumerate}

Autopoiesis made self-production and boundary generation central to
living organization. Later organizational accounts developed the idea of
closure of constraints: relatively stable constraints act on underlying
processes while mutually depending on and maintaining one another
(Varela, Maturana, \& Uribe, 1974; Mossio, Saborido, \& Moreno, 2009;
Montévil \& Mossio, 2015). This paper treats these as major sources
rather than as settled definitions of life.

\subsection{1.3 The bridge from physics to
biology}\label{the-bridge-from-physics-to-biology}

Physics supplies state spaces, dynamics, thermodynamic constraints, and
material processes. It does not by itself derive life. The biological
bridge is crossed when a physically realized organization contains
processes whose effects contribute to maintaining the conditions under
which those processes and the larger organization continue.

The difference is not that living matter violates physics or contains a
new force. It is that certain physical processes enter a recursively
organized control relation. A flame, hurricane, and convection cell can
persist through throughput. Whether any such system counts as
self-maintaining depends on more than duration or dissipation. We must
ask whether the organization is differentially sensitive to departures
from its viability conditions, acts through internally generated
pathways, repairs or regenerates relevant constraints, and preserves a
boundary or causal organization across perturbation.

The bridge is therefore empirical and theoretical. The word
\emph{coherence} cannot cross it by itself. Foundations records this
transition through the Biology-owned principle CU-B1 and treats the
dynamics principle CU-D3 as a gateway rather than an independent law.

\subsection{1.4 Scope boundaries}\label{scope-boundaries}

This paper does not claim that biology has one essence captured by a
universal checklist. Viruses, dormant spores, sterile organisms,
obligate symbionts, colonial organisms, microbiomes, and artificial
living systems place pressure on any sharp demarcation.

It does not claim that every successful biological process maximizes
coherence. Selection can favor fragility, exploitation, short-horizon
reproduction, or specialization. An organism may preserve itself by
damaging its environment; a cancer clone may thrive while the organism
fails. Every assessment must identify its boundary and horizon.

It also does not derive moral value from survival. That a process
supports an organism's continuation grounds a biological function and an
organism-relative norm. It does not establish that the continuation is
morally good, that every living entity has equal moral standing, or that
death is an ethical failure. Those are later normative questions.

\subsection{1.5 Abiogenesis, morphogenesis, and
evolution}\label{abiogenesis-morphogenesis-and-evolution}

Abiogenesis, morphogenesis, and evolution do not share one mechanism.
Their legitimate unity is more abstract:

\begin{itemize}
\tightlist
\item
  \textbf{abiogenesis} concerns how chemical processes became organized
  so that some of their own products and constraints contributed to
  continuation and reproduction;
\item
  \textbf{morphogenesis} concerns how an already living system
  constructs and maintains viable form across changing local states;
\item
  \textbf{evolution} concerns how heritable variation and differential
  reproduction alter the architectures by which organisms and
  collectives maintain, reproduce, and transform themselves.
\end{itemize}

Calling all three ``coherence dynamics'' is useful only when the
domain-specific mechanisms remain visible.

\subsection{1.6 What would count as
progress}\label{what-would-count-as-progress}

The proposal succeeds scientifically only if it improves at least one of
the following:

\begin{itemize}
\tightlist
\item
  \textbf{operationalization:} turning organization, integration, or
  stake into measurable variables;
\item
  \textbf{comparison:} explaining why systems with similar outputs
  differ in robustness, autonomy, or individuality;
\item
  \textbf{prediction:} identifying interventions whose outcomes are not
  already implied by local mechanism alone;
\item
  \textbf{compression:} relating apparently separate biological
  phenomena without erasing their mechanisms;
\item
  \textbf{failure analysis:} locating where maintenance breaks and which
  scale bears the cost.
\end{itemize}

A vocabulary that merely renames homeostasis, autonomy, fitness, or
control as coherence would add little.

\subsection{1.7 Relation to the series and
roadmap}\label{relation-to-the-series-and-roadmap}

\emph{Metaphysics and Epistemology} defines coherence and the
framework's claim registers. \emph{Foundations} records the shared
principle architecture and assigns CU-B1 through CU-B5 to Biology.
\emph{Physics} supplies the thermodynamic and dynamical constraints
within which living systems exist. The later papers examine candidate
perspective, psychological regulation, social organization, ethics, and
artificial systems.

Sections 2--5 establish the conceptual and physical basis. Sections
6--10 apply it to origins, development, agency, evolution, and
pathology. Section 11 considers synthetic and artificial cases. Section
12 states the experimental program and failure conditions. Section 13
separates scientific conclusions from the framework's spiritual
interpretation. Formal propositions, illustrative models, the complete
claim ledger, and the audit record are preserved in the technical
supplement.

\begin{center}\rule{0.5\linewidth}{0.5pt}\end{center}

\section{Section 2. Conceptual Discipline: organization, constraint, and
function}\label{section-2.-conceptual-discipline-organization-constraint-and-function}

\subsection{2.1 Coherence in the biological
domain}\label{coherence-in-the-biological-domain}

\textbf{{[}Definition{]}} For a declared candidate system \(B\),
biological coherence is the integrated organization by which \(B\)
maintains, repairs, reproduces, or adaptively transforms its
viability-bearing relations under a stated perturbation class and
horizon.

In the canonical dependency notation,

\[\adjustbox{max width=\linewidth}{$\displaystyle 
C_B=C_B(x;K,G,\pi,T),
$}\]

but in this paper \(C_B\) is normally a \textbf{profile}, not a scalar.
The hidden arguments are:

\begin{itemize}
\tightlist
\item
  \(x\): the state, trajectory, or organization being assessed;
\item
  \(K\): viability-bearing and functional constraints;
\item
  \(G\): causal, material, regulatory, and informational relations;
\item
  \(\pi\): system boundary, scale, measurement map, or partition;
\item
  \(T\): the relevant temporal horizon.
\end{itemize}

The phrase ``the organism became more coherent'' is incomplete until
these arguments and the comparison rule are given.

\subsection{2.2 Constraint types}\label{constraint-types}

Biology uses \emph{constraint} in several senses. This paper
distinguishes at least six.

\subsubsection{2.2.1 Physical constraints}\label{physical-constraints}

Conservation laws, stoichiometry, diffusion, reaction kinetics,
mechanics, and thermodynamic boundary conditions restrict possible
trajectories whether or not a living organization exists.

\subsubsection{2.2.2 Viability constraints}\label{viability-constraints}

These mark boundaries beyond which the organization cannot continue,
cannot recover within the chosen intervention set, or loses an
identity-relevant capacity. Examples may include membrane integrity,
redox balance, temperature range, energy charge, tissue perfusion, or
reproductive continuity. Their exact form is system-relative.

\subsubsection{2.2.3 Regulatory
constraints}\label{regulatory-constraints}

These are relatively stable structures or processes that channel faster
processes: enzymes, membranes, gene-regulatory networks, transport
systems, signaling circuits, tissue architectures, and behavioral
policies. Their stability is always timescale-relative.

\subsubsection{2.2.4 Generative
constraints}\label{generative-constraints}

Some restrictions create capacities. A membrane excludes many exchanges
while making concentration gradients, selective transport, and a
persistent reaction environment possible. Developmental canalization
narrows trajectories while enabling reliable form.

\subsubsection{2.2.5 Historical
constraints}\label{historical-constraints}

Development, learning, epigenetic modification, symbiosis, and evolution
alter which future trajectories are reachable. The current organization
cannot always be explained from the instantaneous state alone.

\subsubsection{2.2.6 Observer-imposed
constraints}\label{observer-imposed-constraints}

Experimental protocols, model boundaries, engineering goals, and
classification schemes may be useful without belonging to the organism.
Confusing these with organism-relative constraints is a central failure
mode.

\subsection{2.3 Persistence is not yet
self-maintenance}\label{persistence-is-not-yet-self-maintenance}

A state can persist because:

\begin{enumerate}
\def\labelenumi{\arabic{enumi}.}
\tightlist
\item
  it is near equilibrium;
\item
  an environment continually recreates it;
\item
  an external controller holds it in place;
\item
  it lies in a dynamical attractor;
\item
  it reproduces through a template;
\item
  its own organization regulates and repairs the conditions of its
  continuation.
\end{enumerate}

Only the sixth is the strongest form of biological self-maintenance
proposed here. The categories can overlap. A cell uses environmental
free energy, inherits evolved structure, occupies dynamical attractors,
and may depend on a host. The point is not to demand independence from
everything external. No organism is independent. The question is whether
a relatively bounded organization makes an endogenous, counterfactual
contribution to its own continued organization.

\subsection{2.4 Function without
foresight}\label{function-without-foresight}

Biological function can be described historically, dispositionally, or
organizationally. This paper uses an organizational criterion where
appropriate:

\begin{quote}
\textbf{{[}Definition{]}} A component or process has an organizational
function relative to \(B\) when its causal activity contributes to
maintaining one or more conditions required for the persistence or
reproduction of \(B\), and the organization of \(B\) contributes to
producing, repairing, regulating, or retaining that component or
process.
\end{quote}

This reciprocal relation grounds a limited form of teleology without
future causation or conscious purpose. The heart circulates blood
because that effect contributes to the organism whose development and
maintenance produce a functioning heart. A stone that happens to block
wind may benefit an organism but does not thereby have a biological
function within it.

The criterion is not exhaustive. Selected-effect accounts remain
important for evolutionary explanation, and artifacts have functions
assigned by designers. This paper does not force all function talk into
one theory.

\subsection{2.5 Organism-relative
normativity}\label{organism-relative-normativity}

Once a candidate organization and its viability boundaries are
specified, states can be evaluated relative to its continuation. A
membrane rupture, loss of electrochemical gradient, or failure of repair
can be \emph{bad for the organism} in a functional sense. This is not
merely the researcher's preference if the system's own regulatory
processes are organized around preventing or correcting those states.

The word \emph{normativity} must be scoped. This paper distinguishes:

\begin{itemize}
\tightlist
\item
  \textbf{functional normativity:} correct or defective relative to a
  biological function;
\item
  \textbf{viability normativity:} better or worse relative to continued
  organization;
\item
  \textbf{phenomenal valence:} felt better or worse for a subject;
\item
  \textbf{ethical normativity:} morally permitted, required, or
  forbidden.
\end{itemize}

The first two do not entail the latter two.

\subsection{2.6 Why one scalar is not
canonical}\label{why-one-scalar-is-not-canonical}

The former Biology paper used

\[\adjustbox{max width=\linewidth}{$\displaystyle 
C=A I,
$}\]

where \(A\) aggregated constraint satisfaction and \(I\) aggregated
integration. This construction had heuristic virtues: it prevented
excellent local performance from compensating automatically for zero
integration, and it forced attention to both parts and whole. But it was
not derived, its dimensions and normalizations were not stable across
systems, and its integration term combined coupling, synergy, and
robustness without proving that they formed one quantity.

More importantly, biological dimensions genuinely conflict. Strong
coupling can improve robustness to a common shock while reducing
flexibility under asymmetric environmental change. Redundancy can raise
resilience while increasing energetic cost. Rapid growth can improve
lineage fitness while reducing repair reserve. A scalar can still be
built for a particular decision, but the weights then express a model or
value choice.

This paper therefore uses a structured profile and preserves Pareto
incomparability where warranted.

\subsection{2.7 Boundary and horizon
discipline}\label{boundary-and-horizon-discipline}

A coherence claim can reverse when the boundary changes.

\begin{itemize}
\tightlist
\item
  A cancer cell may be robust at the cell scale and destructive at the
  organism scale.
\item
  An obligate symbiont may appear non-autonomous in isolation and
  functionally integrated within a host consortium.
\item
  A short-lived reproductive strategy may succeed over one generation
  and undermine lineage persistence under environmental change.
\item
  A dormant spore may show little active regulation over a short
  observation window while preserving a capacity for reactivation over a
  longer one.
\end{itemize}

Every substantive assessment must state:

\[\adjustbox{max width=\linewidth}{$\displaystyle 
(B,\pi,T,\mathcal D,\mathcal A),
$}\]

where \(B\) is the candidate system, \(\pi\) the chosen boundary and
scale, \(T\) the horizon, \(\mathcal D\) the disturbance class, and
\(\mathcal A\) the admissible internal and external intervention class.

\subsubsection{2.7.1 Claim-governance
protocol}\label{claim-governance-protocol}

The standardized Domain Analytic Frame above is this paper's compact
Coherence Analysis Protocol record under proposed CU-M11--CU-M14. The
detailed rules below remain the domain-specific implementation.

The following rules govern every biological coherence claim in this
manuscript and in downstream CU papers:

\begin{enumerate}
\def\labelenumi{\arabic{enumi}.}
\tightlist
\item
  \textbf{Boundary rule:} name the candidate biological unit or units. A
  cell-, tissue-, organism-, lineage-, and ecosystem-level assessment
  may order the same process differently.
\item
  \textbf{Horizon rule:} name the temporal horizon. Maintenance,
  resilience, adaptation, and failure are incomplete claims without a
  timescale.
\item
  \textbf{Disturbance rule:} name the perturbation or disturbance class.
  Robustness under one challenge family does not imply robustness under
  another.
\item
  \textbf{Profile rule:} declare viability floors, weights, aggregation
  rules, and protected dimensions whenever a profile is scalarized. A
  scalar ranking is a model or decision rule, not a framework-level
  discovery.
\item
  \textbf{Register rule:} do not promote organism-relative function or
  stake directly into phenomenal valence, moral patiency, or ethical
  worth.
\end{enumerate}

A claim that omits these declarations should be treated as
underspecified rather than silently completed by the reader. The
technical supplement records the corresponding registry and downstream
crosswalk.

\subsection{2.8 Relation to established
approaches}\label{relation-to-established-approaches}

CU Biology is closest to a family of approaches rather than one school:

\begin{itemize}
\tightlist
\item
  \textbf{autopoiesis} emphasizes self-production and boundary
  constitution;
\item
  \textbf{organizational closure} analyzes mutually sustaining
  constraints;
\item
  \textbf{viability theory} asks which states can remain inside
  constraint sets under admissible controls;
\item
  \textbf{control theory} specifies feedback, setpoints, disturbances,
  observability, and controllability;
\item
  \textbf{systems biology} reconstructs networks and multiscale
  interactions;
\item
  \textbf{developmental systems and evo-devo} examine how inherited
  resources and developmental constraints shape form;
\item
  \textbf{major-transition theory} studies the emergence of new levels
  of individuality;
\item
  \textbf{active-inference approaches} model self-organization and
  action through probabilistic generative dynamics;
\item
  \textbf{basal-cognition research} investigates adaptive capacities in
  aneural organisms and tissues.
\end{itemize}

The contribution proposed here is not that these fields unknowingly
describe one hidden equation. It is that their different objects can be
located within one declared schema: boundaries, constraints, causal
organization, viable trajectories, endogenous control, history, and
cost.

\begin{center}\rule{0.5\linewidth}{0.5pt}\end{center}

\section{Section 3. A formal architecture for biological
coherence}\label{section-3.-a-formal-architecture-for-biological-coherence}

\subsection{3.1 Candidate-system
specification}\label{candidate-system-specification}

Let a candidate biological system be represented by

\[\adjustbox{max width=\linewidth}{$\displaystyle 
\mathcal B
=
(X,E,U,F,K,G,R,H,\pi,T).
$}\]

Here:

\begin{itemize}
\tightlist
\item
  \(X\) is a state space;
\item
  \(E\) is an environment or family of environmental processes;
\item
  \(U\) is the set of admissible regulatory actions available to the
  candidate system;
\item
  \(F\) specifies the dynamics;
\item
  \(K\) is a set of viability and functional constraints;
\item
  \(G\) is a relational graph, hypergraph, field, or causal structure
  among components and scales;
\item
  \(R\) is the set of regulatory processes and their realization;
\item
  \(H\) records history or path-dependent state;
\item
  \(\pi\) specifies the system boundary and observation map; and
\item
  \(T\) is the evaluation horizon.
\end{itemize}

A continuous-time stochastic representation is

\[\adjustbox{max width=\linewidth}{$\displaystyle 
dx_t
=
f(x_t,e_t,u_t;\theta,H_t)\,dt
+
\Sigma(x_t,e_t)\,dW_t,
$}\]

with policy

\[\adjustbox{max width=\linewidth}{$\displaystyle 
u_t=\rho(y_{0:t},H_t),
$}\]

where \(y_t\) is what the system's regulatory organization can sense or
transduce. Discrete, hybrid, agent-based, reaction-network, and field
models are equally admissible. The tuple is a specification discipline,
not a claim that every organism is optimally represented by a stochastic
differential equation.

\subsection{3.2 Viability region}\label{viability-region}

Let the relevant constraints be functions

\[\adjustbox{max width=\linewidth}{$\displaystyle 
k_j(x,e,H)\ge 0,
\qquad j=1,\ldots,m.
$}\]

The instantaneous viability region is

\[\adjustbox{max width=\linewidth}{$\displaystyle 
V(e,H)
=
\{x\in X:k_j(x,e,H)\ge0\ \text{for all }j\}.
$}\]

This region is model-relative. It can contain hard boundaries, graded
margins, and context-dependent conditions. Not every biologically
relevant variable should be treated as a hard viability constraint.

The normalized margin of constraint \(j\) may be written

\[\adjustbox{max width=\linewidth}{$\displaystyle 
m_j(x,e,H)
=
\frac{k_j(x,e,H)}{s_j},
$}\]

where \(s_j\) is a declared scale. A conservative instantaneous margin
is

\[\adjustbox{max width=\linewidth}{$\displaystyle 
m_{\min}(x)=\min_j m_j(x).
$}\]

This minimum does not measure the whole organism. It records proximity
to the most threatened modeled boundary.

\subsection{3.3 Viability kernel}\label{viability-kernel}

A point can lie in \(V\) while being unable to remain there. The
trajectory-level object is the finite-horizon viability kernel developed
in viability theory (Aubin, 1991):

\[\adjustbox{max width=\linewidth}{$\displaystyle 
\operatorname{Viab}_{T}(V;F,U)
=
\left\{
 x_0\in V:
 \exists u(\cdot)\in\mathcal A
 \text{ such that }
 x_t\in V
 \text{ for all }t\in[0,T]
\right\}.
$}\]

For indefinite maintenance, one considers \(T\to\infty\) when the model
permits it. In biological applications, finite horizons are often more
meaningful because organisms develop, reproduce, senesce, and inhabit
changing environments.

The viability kernel distinguishes present health from maintainability.
Two organisms can have the same measured state while differing in
reserves, redundancy, or control architecture, and therefore in the
disturbances from which they can recover.

\subsection{3.4 Robust viability}\label{robust-viability}

Let \(\mathcal D\) be an admissible disturbance class. The robust
viability kernel is

\[\adjustbox{max width=\linewidth}{$\displaystyle 
\operatorname{RViab}_{T}(V;F,U,\mathcal D)
=
\left\{
 x_0\in V:
 \exists \rho
 \ \forall d\in\mathcal D,
 \ x_t^{\rho,d}\in V
 \text{ for all }t\in[0,T]
\right\}.
$}\]

The quantifier order matters. The policy must cope with every
disturbance in the declared class rather than being chosen after the
entire disturbance trajectory is known.

A system with a larger robust kernel under the same boundary,
constraints, action set, and disturbance class is more robust in one
precise sense. This does not imply greater efficiency, evolvability, or
moral value.

\subsection{3.5 Endogenous maintenance
contribution}\label{endogenous-maintenance-contribution}

The phrase ``the system maintains itself'' should be tested
counterfactually.

Let \(\tau_V\) be the first exit time from the viability region. For a
regulatory subsystem \(R_i\), define an intervention-based maintenance
contribution

\[\adjustbox{max width=\linewidth}{$\displaystyle 
\Delta_i(T)
=
\Pr(\tau_V>T\mid \operatorname{do}(R_i=\text{intact}))
-
\Pr(\tau_V>T\mid \operatorname{do}(R_i=\text{disabled})),
$}\]

under matched environments and appropriate controls.

A positive \(\Delta_i\) does not by itself establish self-maintenance.
The subsystem may be externally installed or may support one function
while being irrelevant to organizational closure. But a network of
positive contributions, together with reciprocal production or repair
relations, makes the claim empirically tractable.

A stronger system-level comparison contrasts the intact candidate
organization with a boundary-matched passive or externally controlled
surrogate:

\[\adjustbox{max width=\linewidth}{$\displaystyle 
\Delta_{\mathrm{endo}}(T)
=
\Pr(\tau_V>T\mid \mathcal B_{\mathrm{intact}})
-
\Pr(\tau_V>T\mid \mathcal B_{\mathrm{passive/control}}).
$}\]

\subsection{3.6 Biological coherence
profile}\label{biological-coherence-profile}

This paper proposes the profile

\[\adjustbox{max width=\linewidth}{$\displaystyle 
\mathbf C_B
=
\bigl(
M_T,
R_T,
I_{G,T},
O_T,
A_T,
H_T;
J_T
\bigr),
$}\]

where:

\begin{itemize}
\tightlist
\item
  \(M_T\): viability margin or survival probability over horizon \(T\);
\item
  \(R_T\): resilience and recovery performance under a declared
  perturbation distribution;
\item
  \(I_{G,T}\): causal integration relevant to maintenance across the
  declared graph and scales;
\item
  \(O_T\): organizational closure or endogeneity of the maintenance
  network;
\item
  \(A_T\): adaptive reach---the disturbance or environmental set across
  which viable reconfiguration is possible;
\item
  \(H_T\): historical continuity of identity-relevant organization; and
\item
  \(J_T\): cost and externalization profile, including matter, work,
  entropy production, ecological extraction, or dependence on external
  support.
\end{itemize}

The semicolon before \(J_T\) is deliberate. Cost is not simply
``negative coherence.'' It is a separate accounting dimension that
constrains whether apparent maintenance is sustainable or merely
subsidized by another system.

No component is assumed to be universally measurable in the same units.
Each requires a model and protocol.

\subsection{3.7 Resilience}\label{resilience}

For perturbations \(p\) drawn from a declared distribution
\(\mathcal P\), let \(\tau_{\mathrm{return}}(p)\) be the time required
to return to a core viable set \(V_{\mathrm{core}}\), with failure
assigned when no return occurs within \(T\). Candidate resilience
summaries include

\[\adjustbox{max width=\linewidth}{$\displaystyle 
R_T^{(1)}
=
\Pr_{p\sim\mathcal P}
\bigl(
\tau_{\mathrm{return}}(p)\le T
\bigr),
$}\]

and

\[\adjustbox{max width=\linewidth}{$\displaystyle 
R_T^{(2)}
=
\mathbb E_{p\sim\mathcal P}
\left[
\frac{1}{1+\tau_{\mathrm{return}}(p)}
\right].
$}\]

Neither is universal. The perturbation distribution, core set, censoring
rule, and intervention class determine the result.

Resilience is not identical with robustness. A robust system changes
little under perturbation. A resilient system may change substantially
and recover or reorganize. Plasticity can contribute to resilience while
reducing short-run stability.

\subsection{3.8 Causal integration}\label{causal-integration}

Integration should not be inferred from correlation alone. This paper
uses \textbf{maintenance-relevant causal integration} as a family of
measures.

Let the candidate system be decomposed into modules
\(X^{(1)},\ldots,X^{(n)}\). Let \(Q\) be a maintenance outcome such as
survival time, recovery probability, or preserved function. One
intervention-based measure is the performance loss caused by severing or
scrambling cross-module relations while preserving local dynamics as
closely as possible:

\[\adjustbox{max width=\linewidth}{$\displaystyle 
I_{G,T}
=
Q(\mathcal B_{\text{intact}})
-
Q(\mathcal B_{\text{cut}(G)}).
$}\]

Alternative measures may use causal influence, transfer entropy,
integrated control capacity, synergy decompositions, or network
controllability. Each answers a different question. Transfer entropy,
for example, measures directed predictive information under assumptions;
it does not by itself establish function, agency, or organizational
closure (Schreiber, 2000).

The phrase ``global constraint'' is acceptable only when the relevant
macrovariable or cross-scale relation changes local trajectories in an
interventionally identifiable way.

\subsection{3.9 Organizational closure}\label{organizational-closure}

Represent regulatory constraints or maintenance processes as nodes in a
directed graph \(G_R\). Draw an edge \(R_i\to R_j\) when the activity of
\(R_i\) produces, repairs, energizes, or stabilizes a condition required
for \(R_j\) to operate over the relevant timescale.

A candidate closure core is a strongly connected subgraph in which every
member receives at least one maintenance contribution from another
member and contributes to at least one other. This graph criterion is
necessary only within the chosen representation; it is not sufficient
for life. A formal cycle can be externally powered, externally
assembled, or dynamically irrelevant.

Accordingly, the closure profile should include:

\begin{itemize}
\tightlist
\item
  causal contribution under intervention;
\item
  material and energetic realization;
\item
  timescale compatibility;
\item
  regeneration or replacement of constraint bearers;
\item
  boundary relation;
\item
  and dependence on external scaffolding.
\end{itemize}

Autonomy is graded because closure is always relative. Cells depend on
environments, organisms depend on ecosystems, and symbionts depend on
hosts. The question is which dependencies are internalized into the
candidate unit and which remain external.

\subsection{3.10 Adaptive reach}\label{adaptive-reach}

Let \(\mathcal E\) be a family of environments or disturbance regimes
with measure \(\mu\). Define the viable adaptation set

\[\adjustbox{max width=\linewidth}{$\displaystyle 
\mathcal A_T(\mathcal B)
=
\{e\in\mathcal E:
\exists \rho_e
\text{ such that the organization remains viable or reaches an admissible transformed organization through }T
\}.
$}\]

A candidate adaptive-reach measure is

\[\adjustbox{max width=\linewidth}{$\displaystyle 
A_T=\mu(\mathcal A_T).
$}\]

This quantity depends on the environment family and measure. It should
not be compared across studies without a common challenge space.
Adaptive reach is also not unbounded flexibility: a system that can
become anything may preserve no identity.

\subsection{3.11 Historical continuity and biological
identity}\label{historical-continuity-and-biological-identity}

The former paper represented identity by small distance between
successive constraint manifolds,

\[\adjustbox{max width=\linewidth}{$\displaystyle 
D(K_{t+1},K_t)\approx0.
$}\]

That condition is neither necessary nor sufficient. Development can
transform the relevant constraints dramatically while preserving an
organism. A static counterfeit can remain close in state space without
sharing causal continuity.

This paper therefore treats token identity as a history-sensitive
relation. Let

\[\adjustbox{max width=\linewidth}{$\displaystyle 
\mathcal R_{\mathrm{cont}}(B_t,B_{t+\Delta})
$}\]

hold when there exists an admissible causal process connecting the two
stages that preserves enough of the organization responsible for
boundary maintenance, regulation, repair, and developmental or lineage
continuity. The phrase ``enough'' is domain-specific and sometimes
contested.

Identity is not material sameness. It is also not arbitrary pattern
similarity. It is continuity through an organization-preserving process.

\subsection{3.12 Scalarization rules}\label{scalarization-rules}

A scalar score may be constructed for a specific experiment or decision:

\[\adjustbox{max width=\linewidth}{$\displaystyle 
C_B^{(w)}
=
S_w
\left(
\phi_M(M_T),
\phi_R(R_T),
\phi_I(I_{G,T}),
\phi_O(O_T),
\phi_A(A_T),
\phi_H(H_T),
\phi_J(J_T)
\right).
$}\]

To be admissible, the paper or experiment must state:

\begin{enumerate}
\def\labelenumi{\arabic{enumi}.}
\tightlist
\item
  the normalization maps \(\phi\);
\item
  the aggregator \(S_w\);
\item
  the weights or ordering;
\item
  protected floors that cannot be compensated away;
\item
  how uncertainty is propagated;
\item
  the boundary and horizon; and
\item
  sensitivity of conclusions to these choices.
\end{enumerate}

A viability-first lexicographic ordering will often be more defensible
than unrestricted weighted addition:

\begin{enumerate}
\def\labelenumi{\arabic{enumi}.}
\tightlist
\item
  reject trajectories that breach protected viability floors;
\item
  among the remainder, compare recovery and closure;
\item
  only then compare secondary performance and cost.
\end{enumerate}

The old product \(A I\) is permitted as one member of this model family,
not as the biological definition of coherence.

\subsection{3.13 Attractors, gradients, and target
sets}\label{attractors-gradients-and-target-sets}

Attractor language is useful when a model supplies dynamics. A
developmental or homeostatic regime may be represented by an invariant
set \(\mathcal A\), a basin of attraction, or a controlled target set.
But three distinctions are required.

First, an attractor need not be viable. A disease state can be
dynamically stable.

Second, a viable state need not be an attractor. Active control may be
required to keep the system there.

Third, an organism need not follow the gradient of one scalar. Feedback
controllers, hybrid switching, stochastic search, predictive policies,
mechanical constraints, and evolutionary inheritance can generate
directed behavior without global gradient ascent.

A coherence gradient is therefore used only when a declared scalar and
metric exist:

\[\adjustbox{max width=\linewidth}{$\displaystyle 
\operatorname{grad}_{g}C_B^{(w)}.
$}\]

No framework-level \(\nabla C\) governs all living systems.

\subsection{3.14 Failure and collapse}\label{failure-and-collapse}

Biological failure is not one universal bifurcation. This paper
distinguishes:

\begin{itemize}
\tightlist
\item
  \textbf{constraint breach:} one modeled viability boundary is crossed;
\item
  \textbf{functional failure:} a capacity is lost while the organism
  persists;
\item
  \textbf{loss of resilience:} current function remains but the
  perturbation-recovery set contracts;
\item
  \textbf{organizational decoupling:} cross-scale relations cease to
  coordinate maintenance;
\item
  \textbf{identity transition:} the candidate unit changes, as in some
  symbioses, metamorphoses, or major evolutionary transitions;
\item
  \textbf{irrecoverable organismal failure:} no admissible internal or
  external intervention in the declared class restores organism-level
  self-maintenance.
\end{itemize}

Threshold and cascade language is justified only where a model or
experiment identifies a control parameter, order parameter, and
discontinuity, bifurcation, tipping probability, or critical slowing
signature.

\subsection{3.15 Memory and path
dependence}\label{memory-and-path-dependence}

Biological memory includes any history-dependent modification that
changes present dynamics or future possibilities. Examples include
molecular modification, gene-regulatory state, immune memory, tissue
architecture, developmental commitment, neural plasticity, ecological
inheritance, and lineage-level genetic change.

A generic representation is

\[\adjustbox{max width=\linewidth}{$\displaystyle 
H_{t+1}=\mathcal F_H(H_t,x_t,e_t,u_t),
\qquad
x_{t+1}=F(x_t,e_t,u_t;H_t).
$}\]

Memory is not always stored in a localized code. It may be distributed
across concentrations, network topology, mechanical structure, or
environmental scaffolding. Nor does every memory bind a persisting
conscious subject.

\subsection{3.16 What the formal architecture
establishes}\label{what-the-formal-architecture-establishes}

The architecture establishes a disciplined space of models. It does not
prove that every organism possesses a unique boundary, that every
relevant quantity is measurable, or that one coherence profile will
outperform established domain-specific models.

Its immediate value is negative and constructive:

\begin{itemize}
\tightlist
\item
  it blocks circular definitions of life as whatever has high coherence;
\item
  it blocks universal gradient language;
\item
  it makes disturbances, controls, costs, and horizons explicit;
\item
  it distinguishes present performance from maintainability;
\item
  and it identifies intervention points by which ``self-maintenance''
  can be tested rather than admired.
\end{itemize}

The technical supplement proves several bounded propositions about these
objects and presents two illustrative configurations in which the same
diffusive coupling law produces opposite viability trends, demonstrating
why neither coupling nor a scalarized profile can be interpreted without
declared floors, disturbances, control, and horizon.

\begin{center}\rule{0.5\linewidth}{0.5pt}\end{center}

\section{Section 4. Biological individuality and organism-relative
stake}\label{section-4.-biological-individuality-and-organism-relative-stake}

\subsection{4.1 Why the biological individual cannot be
assumed}\label{why-the-biological-individual-cannot-be-assumed}

Biology routinely speaks about \emph{the organism}, but the unit whose
persistence is being explained is not always given in advance. A mammal
appears comparatively straightforward until one considers its
microbiota, inherited symbionts, immune tolerance, grafts, pregnancies,
cancers, and continuously replaced matter. Colonial organisms, lichens,
slime molds, eusocial colonies, and endosymbiotic lineages make the
problem explicit: physical contact, genetic identity, reproductive
alignment, metabolic dependence, and causal integration can identify
different boundaries.

This paper therefore treats biological individuality as an empirical
model-selection problem. A proposed individual is a candidate causal
unit whose boundary and internal organization explain maintenance,
reproduction, adaptation, or evolutionary change better than relevant
alternatives. There may be several legitimate units at different scales.

This position avoids two opposite errors. The first is atomism: assuming
that cells or organisms are fundamental individuals regardless of their
dependencies. The second is indiscriminate holism: declaring an
ecosystem or biosphere to be one organism because its parts interact.
Interaction is ubiquitous. Individuation requires a stronger pattern of
integration, common fate, regulatory dependence, or reproduction.

\subsection{4.2 Dimensions of
individuality}\label{dimensions-of-individuality}

No one criterion is sufficient across all cases. This paper
distinguishes at least six dimensions.

\subsubsection{4.2.1 Boundary
individuation}\label{boundary-individuation}

A candidate individual controls exchanges across a relatively persistent
boundary. The boundary can be a membrane, epithelium, behavioral
territory, chemical interface, or dynamical separation. It need not be
impermeable. What matters is selective regulation of flows relevant to
continued organization.

A boundary is not enough. A lipid vesicle can have a membrane without
metabolism, regulation, or reproduction. Conversely, a social insect
colony can display organism-like integration without one continuous
membrane.

\subsubsection{4.2.2 Metabolic
individuation}\label{metabolic-individuation}

A candidate unit coordinates material and energetic transformations so
that the resulting products support the unit's persistence. Metabolic
coupling can draw components into a common physiological fate. Yet
parasites can exploit a host metabolism without becoming part of the
same individual, and organelles can retain metabolic specialization
while being deeply integrated into a cell.

\subsubsection{4.2.3 Regulatory and causal
individuation}\label{regulatory-and-causal-individuation}

A candidate unit acts as a control system when disturbances to one part
elicit coordinated changes elsewhere that restore organism-level
variables or functions. This is the dimension most directly captured by
the maintenance contribution and causal-integration measures of Section
3.

A useful intervention test compares explanatory boundaries. Let
candidate partitions \(\mathcal P_1,\mathcal P_2,\dots\) represent
alternative units. For a maintenance outcome \(Q\), estimate how well
interventions and responses are predicted when cross-boundary influences
are removed or aggregated. The preferred boundary for that question is
the one that preserves the most maintenance-relevant causal structure
without importing unnecessary variables.

There may be no unique winner. Different partitions may best explain
metabolism, development, behavior, or reproduction.

\subsubsection{4.2.4 Reproductive
individuation}\label{reproductive-individuation}

A unit can be individuated by participation in a reproducing life cycle.
Bottlenecks often align component fates and reduce conflict because
descendants begin from a small propagule. But reproduction is not
required for every token organism: sterile workers, infertile adults,
and somatic tissues remain biological. Reproduction is more central to
lineage individuality than to moment-by-moment physiological
individuality.

\subsubsection{4.2.5 Evolutionary
individuation}\label{evolutionary-individuation}

A collective becomes an evolutionary individual when selection
increasingly acts on collective-level variation, heredity, and
differential reproduction, while within-collective conflict is
suppressed or redirected. Major-transition theory studies precisely
these reorganizations of the unit of selection (Maynard Smith \&
Szathmáry, 1995; West et al., 2015).

Evolutionary individuality is not identical with physiological autonomy.
A unit can be a strong physiological individual but a weak reproducer,
or a reproductive lineage can depend heavily on another organism.

\subsubsection{4.2.6 Historical
individuation}\label{historical-individuation}

Some biological identity is carried by causal history. A regenerated
planarian, metamorphosed insect, or transplanted tissue may differ
greatly from its earlier state while remaining historically continuous.
A perfect structural copy, by contrast, would not thereby become the
same token organism.

Historical continuity explains why identity cannot be reduced to state
similarity. It must also be constrained: any later state is causally
descended from earlier states, but not every causal descendant preserves
the same organismal organization.

\subsection{4.3 A multidimensional individuality
profile}\label{a-multidimensional-individuality-profile}

For a candidate unit \(B\), define an individuality profile

\[\adjustbox{max width=\linewidth}{$\displaystyle 
\mathbf I_B
=
\bigl(
I_{\partial},
I_{M},
I_{R},
I_{C},
I_{\mathrm{rep}},
I_{E},
I_{H}
\bigr),
$}\]

where the components summarize boundary control, metabolic integration,
regulatory integration, conflict mediation, reproductive unity,
evolutionary individuality, and historical continuity.

This profile is question-relative. A microbiome may score high in
metabolic contribution but low in reproductive unity. A eusocial colony
may score high in conflict mediation and collective reproduction while
individual workers retain behavioral and physiological individuality. A
tumor lineage can gain evolutionary individuality as it loses
integration with the host organism.

Calling one of these entities \emph{the} individual requires an
explanatory purpose. This paper does not force all biological questions
to use the same boundary.

\subsection{4.4 Nested and overlapping
organizations}\label{nested-and-overlapping-organizations}

Biological systems are nested but not perfectly hierarchical. Organelles
exist within cells, cells within tissues, tissues within organisms,
organisms within collectives and ecosystems. Yet horizontal gene
transfer, mobile immune cells, maternal-fetal exchange, symbiosis, and
ecological engineering produce overlapping causal organizations.

This paper therefore replaces the image of sealed boxes with a graph or
hypergraph of maintenance dependencies. A component can participate in
several organizations with different strengths and timescales. The
relevant question is not simply ``inside or outside?'' but:

\begin{itemize}
\tightlist
\item
  Which processes maintain this component?
\item
  Which larger organization does the component help maintain?
\item
  What happens to both when the relation is severed?
\item
  Over what horizon does dependence appear?
\item
  Where are conflicts generated and how are they controlled?
\end{itemize}

This view preserves organismal reality without pretending that
biological boundaries are metaphysically absolute.

\subsection{4.5 Organism-relative stake}\label{organism-relative-stake}

The language of \emph{stake} is useful but dangerous because it can
slide from function to feeling.

\textbf{{[}Definition{]}} A condition has \textbf{organism-relative
stake} for a candidate biological organization when variation in that
condition changes the probability or quality of the organization's
continued self-maintenance, and endogenous regulatory processes are
causally organized to prevent, detect, compensate for, or repair adverse
variation.

Stake therefore requires more than an observer labeling one outcome good
and another bad. It requires an asymmetry internal to the organization:

\[\adjustbox{max width=\linewidth}{$\displaystyle 
x\in V_B
\quad\text{supports continued organization, whereas}\quad
x\notin V_B
\quad\text{degrades or terminates it},
$}\]

and the organization's own processes make a counterfactual difference to
which side of that boundary its trajectory occupies.

This is a naturalized normativity of function. Temperature, pH, membrane
integrity, nutrient availability, genomic stability, and tissue
architecture can be better or worse \emph{for} an organism because the
organism exists only within bounded ranges and actively regulates them.

\subsection{4.6 Stake is not yet experience or moral
status}\label{stake-is-not-yet-experience-or-moral-status}

Organism-relative stake does not entail that the organism feels its
state. A thermostat also has a target range in an engineered
description, but its setpoint is supplied by designers and the device
does not regenerate the wider organization that gives the setpoint
purpose. A bacterium has a richer endogenously maintained organization,
yet the existence of biological norms does not settle whether there is
anything it is like to be that bacterium.

This paper therefore separates:

\begin{enumerate}
\def\labelenumi{\arabic{enumi}.}
\tightlist
\item
  \textbf{functional normativity:} a state supports or impairs an
  organized capacity;
\item
  \textbf{organism-relative stake:} the state supports or impairs the
  continued organization whose own processes regulate it;
\item
  \textbf{phenomenal valence:} the state is experienced as good, bad,
  attractive, or aversive;
\item
  \textbf{moral patiency:} the entity's states matter morally in their
  own right; and
\item
  \textbf{moral agency:} the entity can understand and answer for
  reasons.
\end{enumerate}

The transitions among these categories belong partly to the
Consciousness and Ethics papers. Biology supplies necessary evidence
about organization and stake but cannot close the phenomenological or
normative questions by itself.

\subsection{4.7 Functional attribution without
circularity}\label{functional-attribution-without-circularity}

A biological function is often described historically---what selection
favored---or organizationally---what contributes to present
self-maintenance. These accounts can diverge. The heart's pumping
function has both evolutionary history and current organizational
contribution. A novel compensatory pathway may acquire an organizational
role before selection has acted on it. A vestigial structure may retain
history without substantial present contribution.

This paper uses both registers explicitly.

\textbf{Historical function:} trait \(z\) has function \(f\) when
variants producing \(f\) contributed to the reproduction or persistence
of ancestors in the relevant lineage.

\textbf{Organizational function:} component \(z\) has function \(f\) for
\(B\) over horizon \(T\) when performing \(f\) makes an interventionally
detectable contribution to maintaining \(B\), while \(B\) contributes to
producing, repairing, or sustaining \(z\).

The organizational definition blocks one form of observer projection: a
rock can serve as a paperweight, but the paper stack does not maintain
the rock and the rock's paperweight effect does not participate in
closure of the stack.

It does not eliminate judgment. Selecting the system boundary, outcome,
and timescale remains a modeling decision. That is why the claim ledger
must state them.

\subsection{4.8 Conflict and the limits of local
maintenance}\label{conflict-and-the-limits-of-local-maintenance}

A component can maintain itself while undermining the larger unit.
Cancer is the clearest case, but analogous conflicts occur in selfish
genetic elements, meiotic drive, microbial cheating, and social
parasites.

Let \(Q_i\) measure maintenance of component \(i\) and \(Q_B\)
maintenance of the candidate organism. Then an intervention can satisfy

\[\adjustbox{max width=\linewidth}{$\displaystyle 
\Delta Q_i>0,
\qquad
\Delta Q_B<0.
$}\]

There is no contradiction. ``Coherence increased'' is incomplete until
the boundary is named.

Multicellular individuality depends partly on mechanisms that align
local and global interests: developmental bottlenecks, policing,
apoptosis, immune surveillance, resource regulation, and restrictions on
reproduction. These mechanisms are not moral governance. They are
evolved conflict-management architectures.

\subsection{4.9 Edge cases}\label{edge-cases}

\subsubsection{4.9.1 Viruses}\label{viruses}

Viruses reproduce and evolve but rely on host cellular machinery for
metabolism and many regulatory functions. They are strong evolutionary
individuals and weak autonomous physiological individuals. The profile
avoids forcing a binary verdict.

\subsubsection{4.9.2 Dormancy and spores}\label{dormancy-and-spores}

Dormant forms can suspend much active regulation while preserving
structures capable of restarting it. Their current metabolic
self-maintenance may be minimal, but their historical and prospective
organization remains. A definition requiring continuous high activity
would falsely exclude them.

\subsubsection{4.9.3 Sterile organisms}\label{sterile-organisms}

Sterility removes token reproductive capacity without erasing
physiological individuality. Reproduction is therefore not a necessary
condition for each living token, though it remains central to lineage
persistence.

\subsubsection{4.9.4 Obligate symbioses}\label{obligate-symbioses}

When partners cannot maintain or reproduce without one another, the
relevant individual may lie partly at the consortium level. Yet conflict
and separate inheritance may preserve lower-level individuality. The
correct description can be plural.

\subsubsection{4.9.5 Colonies and
superorganisms}\label{colonies-and-superorganisms}

A colony becomes more organism-like as division of labor, reproductive
specialization, communication, conflict control, and common fate
increase. The term \emph{superorganism} should mark a supported profile,
not rhetorical admiration.

\subsubsection{4.9.6 Microbiomes}\label{microbiomes}

A host-associated microbiome can be causally important without forming
one stable evolutionary individual with the host. Composition may be
reacquired, membership may turn over, and interests may diverge.
Specific host-microbe systems should be tested rather than subsumed
under a universal holobiont claim.

\subsubsection{4.9.7 Tumors}\label{tumors}

A tumor is simultaneously part of an organism, a breakdown of
organism-level regulation, and an evolving cell population with partial
autonomy. Its individuality profile changes during progression.

\subsection{4.10 Biology Principle CU-B5}\label{biology-principle-cu-b5}

\begin{quote}
\textbf{CU-B5 --- Biological Individuation and Stake}
\textbf{{[}Domain-authority bridge principle{]}}\\
A biological individual is a historically continuous, relatively bounded
organization whose components are causally integrated into maintaining,
repairing, reproducing, or adaptively transforming the conditions of
that organization. The organization thereby establishes
organism-relative stakes: conditions can be better or worse for it in
virtue of their effects on its continued self-maintenance. Individuality
is graded, multidimensional, and potentially nested; organism-relative
stake does not by itself establish consciousness or moral status.
\end{quote}

Foundations records CU-B5 as the Biology-owned bridge principle
governing the relation between biological individuation,
organism-relative stake, and later questions of experience and moral
standing.

\begin{center}\rule{0.5\linewidth}{0.5pt}\end{center}

\section{Section 5. Thermodynamic
grounding}\label{section-5.-thermodynamic-grounding}

\subsection{5.1 Life does not oppose the second
law}\label{life-does-not-oppose-the-second-law}

Living systems maintain gradients, compartments, macromolecules, and
coordinated dynamics while exchanging matter and energy with their
surroundings. Their local organization does not violate thermodynamics.
For a declared system boundary,

\[\adjustbox{max width=\linewidth}{$\displaystyle 
\frac{dS_B}{dt}
=
\dot S_{e}
+
\dot S_{i},
\qquad
\dot S_i\ge 0,
$}\]

where \(\dot S_e\) is entropy flow across the boundary and \(\dot S_i\)
is internal entropy production. \(\dot S_e\) may be negative when the
system exports entropy, while total entropy production for the closed
system-plus-environment remains nonnegative.

The popular phrase ``life fights entropy'' is therefore acceptable only
as shorthand for active maintenance of selected low-entropy or
far-from-equilibrium structures through throughput and dissipation. It
is not a new thermodynamic law and should not imply that all biological
organization is captured by one entropy value.

\subsection{5.2 Coherence is not low
entropy}\label{coherence-is-not-low-entropy}

Entropy and biological organization answer different questions. A
homogeneous equilibrium state can have high thermodynamic entropy and
little functional differentiation. A frozen crystal can have low entropy
but no metabolism, repair, or regulation. An organism can increase
internal diversity and informational entropy while improving functional
organization.

Nor does Shannon entropy of a biological data distribution equal
thermodynamic entropy without a physical encoding and ensemble. Physics
v0.10 established a restricted accessibility--entropy relation for block
averaging; it did not identify entropy with hidden substantive
coherence.

Consequently:

\[\adjustbox{max width=\linewidth}{$\displaystyle 
C_B\ne -S,
\qquad
C_B\ne S_{\max}-S,
$}\]

as framework-level identities.

A particular study may find that an entropy measure predicts a component
of \(\mathbf C_B\). The physical interpretation must be stated rather
than imported from the word \emph{entropy}.

\subsection{5.3 Free-energy throughput and maintained
constraints}\label{free-energy-throughput-and-maintained-constraints}

A cell maintains concentrations, membrane potentials, polymer
distributions, and reaction networks away from equilibrium by coupling
unfavorable processes to favorable ones. Nutrients, photons, redox
gradients, or host-provided metabolites supply free energy; waste and
heat leave the system.

The important biological fact is not simply dissipation. It is
\textbf{organized dissipation}: flows are channelled through constraints
that are themselves produced, repaired, and regulated by the system.
Enzymes, membranes, transporters, cytoskeletal structures, and
regulatory networks shape process rates while depending on those
processes for their own replacement.

This is where thermodynamics and organizational closure meet.
Thermodynamics constrains what work is possible and what costs must be
paid. Closure explains how a network of maintained constraints directs
that work toward continued organization.

\subsection{5.4 No universal maintenance-cost
law}\label{no-universal-maintenance-cost-law}

No framework-level law presently links a scalar amount of biological
coherence to a minimum maintenance energy, and no universal ratio of
biological coherence to throughput is defined.

There is no established cross-system conversion factor between energetic
work and the heterogeneous components of biological coherence. The
minimum work needed for a task depends on mechanism, kinetics, error
tolerance, boundary conditions, driving protocol, and information
processing. Real organisms also operate far from reversible limits.

Thermodynamic analyses can constrain particular processes---copying,
sensing, erasure, transport, replication, or growth---when the physical
model is specified. They do not support a generic law saying that twice
as much biological coherence costs twice as much energy.

Recent work has also clarified that proposed universal relations between
replicator growth, decay, and dissipation require careful kinetic
assumptions; dissipation alone does not supply a general
organism-independent bound. That correction reinforces the
methodological point: thermodynamic inequalities must travel with their
model assumptions, not with a broad narrative about life (Kolchinsky,
2024).

\subsection{5.5 Efficiency is a Pareto
problem}\label{efficiency-is-a-pareto-problem}

Biological performance involves multiple objectives:

\begin{itemize}
\tightlist
\item
  maintain viability;
\item
  recover from disturbance;
\item
  respond quickly;
\item
  preserve reserves;
\item
  minimize energetic and material costs;
\item
  avoid damaging other scales;
\item
  reproduce or support kin;
\item
  remain adaptable to future conditions.
\end{itemize}

These objectives conflict. Redundancy increases robustness but can raise
maintenance cost. Strong coupling can stabilize common variables but
reduce local flexibility. Rapid growth can reduce stress tolerance. A
highly efficient system near nominal conditions can be fragile under
rare perturbations.

The appropriate mathematical object is often a Pareto frontier. For two
strategies \(\rho_1\) and \(\rho_2\), \(\rho_1\) dominates \(\rho_2\)
only if it is no worse on every protected dimension and better on at
least one. Otherwise the strategies are incomparable until a decision
context supplies priorities.

A scalar ``coherence efficiency'' hides these tradeoffs unless it makes
the aggregation rule explicit.

\subsection{5.6 Robustness has a thermodynamic price, but not a fixed
one}\label{robustness-has-a-thermodynamic-price-but-not-a-fixed-one}

Feedback, proofreading, repair, and redundancy ordinarily require
resources. Yet the relation is not monotonic across all designs. Better
architecture can improve both reliability and energy use; passive
physical structure can reduce control effort; environmental scaffolding
can externalize costs.

A useful accounting identity is

\[\adjustbox{max width=\linewidth}{$\displaystyle 
J_T
=
J_{\mathrm{internal}}
+J_{\mathrm{imported}}
+J_{\mathrm{externalized}},
$}\]

where the terms record costs borne within the system, supplied by
another system, or imposed on the environment. Apparent organismal
efficiency can result from unpaid ecological or host costs.

This point becomes ethically important at higher levels, but it is first
a biological measurement requirement.

\subsection{5.7 Dissipative adaptation is not a complete theory of
life}\label{dissipative-adaptation-is-not-a-complete-theory-of-life}

Driven matter can form ordered structures and can, under some
conditions, preferentially occupy states that absorb and dissipate work.
Such results are relevant to prebiotic organization. They do not
establish that dissipation maximization defines life, that all
adaptation is thermodynamic optimization, or that biological functions
follow from dissipation alone.

A dissipative structure can lack heredity, regulation, boundary repair,
or organizational closure. Conversely, an organism may sacrifice
immediate throughput to survive stress or preserve future options.

This paper therefore treats far-from-equilibrium thermodynamics as
necessary background and a source of mechanisms---not as a replacement
for biological explanation.

\subsection{5.8 Markov blankets and statistical
boundaries}\label{markov-blankets-and-statistical-boundaries}

The free-energy principle and active-inference literature models living
systems through conditional-independence structures often called Markov
blankets. These provide a mathematically useful way to separate
internal, external, sensory, and active states in a stochastic model
(Friston, 2013).

A statistical blanket is not automatically a membrane, organism, or
autonomous agent. Its biological significance depends on whether the
variables correspond to a maintained causal boundary, whether the
partition is robust across interventions and timescales, and whether the
induced dynamics contribute to viability.

This paper therefore permits blanket-based models as one
operationalization of \(\pi\) and \(G\). It does not define life as
whatever admits a Markov blanket; many nonliving stochastic systems do.

\subsection{5.9 The thermodynamic bridge
principle}\label{the-thermodynamic-bridge-principle}

\begin{quote}
\textbf{{[}Bridge principle{]}} Biological self-maintenance requires
physically realized processes that draw on free-energy differences and
export entropy, while constraining those flows so that the organization
responsible for constraint and regulation is preserved or regenerated.
\end{quote}

This statement is strong enough to exclude vitalism and weak enough to
respect the diversity of biological mechanisms. It does not say that
life is a thermodynamic scalar or that physics alone selects the
biological organization.

\subsection{5.10 Failure conditions for the thermodynamic
account}\label{failure-conditions-for-the-thermodynamic-account}

The thermodynamic part of this paper would require revision if:

\begin{enumerate}
\def\labelenumi{\arabic{enumi}.}
\tightlist
\item
  a purported maintenance process could operate indefinitely without
  compensating matter, energy, or entropy exchange in a closed system;
\item
  the proposed system boundary concealed the major costs that sustain
  organization;
\item
  a claimed thermodynamic measure failed to correspond to a physical
  ensemble or encoding;
\item
  an efficiency comparison changed sign under reasonable accounting of
  externalized costs; or
\item
  the biological conclusion depended only on the word \emph{entropy}
  rather than on a stated physical model.
\end{enumerate}

\begin{center}\rule{0.5\linewidth}{0.5pt}\end{center}

\section{Section 6. Origins of life as an autonomy
transition}\label{section-6.-origins-of-life-as-an-autonomy-transition}

\subsection{6.1 Why one threshold is unlikely to be
enough}\label{why-one-threshold-is-unlikely-to-be-enough}

The origin of life is often presented as a boundary between nonlife and
life. Experimentally and theoretically, however, the relevant capacities
are separable. A compartment can persist without metabolism. An
autocatalytic network can reproduce chemical organization without a
membrane. A template can copy without regulating its environment. A
protocell can grow and divide while lacking open-ended heredity. A
reaction network can adapt transiently without forming a lineage.

This paper therefore treats abiogenesis as a \textbf{family of coupled
autonomy transitions}. This separation of autonomy from open-ended
evolution follows a well-established distinction in origin-of-life
theory (Ruiz-Mirazo, Peretó, \& Moreno, 2004). A historical origin may
have crossed these transitions in a particular order, perhaps repeatedly
and with reversals. The task is to explain how capacities became
mutually reinforcing enough to produce a unit that maintained and
reproduced its own organization under variable conditions.

\subsection{6.2 The autonomy-transition
profile}\label{the-autonomy-transition-profile}

For a candidate prebiotic system \(P\), define

\[\adjustbox{max width=\linewidth}{$\displaystyle 
\mathbf A_P
=
\bigl(
a_{\mathrm{boundary}},
a_{\mathrm{energy}},
a_{\mathrm{closure}},
a_{\mathrm{repair}},
a_{\mathrm{heredity}},
a_{\mathrm{viability}},
a_{\mathrm{evolution}}
\bigr),
$}\]

where:

\begin{itemize}
\tightlist
\item
  \(a_{\mathrm{boundary}}\): compartment or boundary control;
\item
  \(a_{\mathrm{energy}}\): endogenous or coupled energy transduction;
\item
  \(a_{\mathrm{closure}}\): catalytic and organizational closure;
\item
  \(a_{\mathrm{repair}}\): regulation and repair;
\item
  \(a_{\mathrm{heredity}}\): heredity with sufficient fidelity and
  variation;
\item
  \(a_{\mathrm{viability}}\): viability across a declared environmental
  range; and
\item
  \(a_{\mathrm{evolution}}\): capacity for cumulative or open-ended
  evolutionary change.
\end{itemize}

These are profiles, not universally normalized numbers. A system can
move along one dimension while regressing on another.

\subsection{6.3 Compartmentalization}\label{compartmentalization}

Compartments concentrate reactants, maintain chemical differences, and
make lineages possible by associating products with the structures that
produced them. Lipid vesicles, coacervates, mineral pores, droplets, and
surface-bound networks provide different mechanisms.

A compartment becomes biologically significant when its growth,
permeability, repair, and division become coupled to the internal
processes it houses. A vesicle supplied and repaired entirely by an
experimenter demonstrates useful chemistry but limited autonomy. If
internal metabolism produces membrane components and membrane properties
preserve that metabolism, a closure relation begins to form.

\subsection{6.4 Autocatalysis and
closure}\label{autocatalysis-and-closure}

Autocatalytic-set theory studies reaction systems in which each reaction
is catalyzed by molecules generated within the set and the required
reactants can be built from an available food set. Reflexively
autocatalytic and food-generated sets show how collective chemical
closure can arise without one self-copying molecule (Hordijk, Kauffman,
\& Steel, 2011).

This is an important transition but not a sufficient definition of life.
Formal autocatalysis can occur in networks that lack compartments,
heritable individuality, regulation, or robust persistence. The food set
and catalysts may also be externally maintained. This paper therefore
records autocatalytic closure as one axis of \(\mathbf A_P\).

\subsection{6.5 Energy transduction}\label{energy-transduction}

Persistent prebiotic organization requires usable free-energy
differences. Candidate sources include redox gradients, light, activated
chemical species, thermal gradients, and geochemical interfaces. The key
transition is not merely having energy available but coupling it to work
that maintains the candidate organization.

A protocell whose internal reactions generate osmotic rupture is not
improved merely because throughput increased. Energy transduction must
become regulated relative to boundary integrity, resource use, and
product removal.

Current experiments on light-driven synthetic protocells demonstrate
increasingly sophisticated coupling among energy input, phase
transition, growth, and structural change (Shi et al., 2025). Other
active-droplet systems link genetically encoded ribozyme catalysis to
cycles of dissolution, segregation, and regrowth (Giessler et al.,
2026). These systems are valuable precisely because they allow the axes
of \(\mathbf A_P\) to be manipulated separately. They do not by
themselves establish that a fully autonomous life cycle has been
achieved.

\subsection{6.6 Regulation before genes}\label{regulation-before-genes}

Regulation need not begin with modern genetic control. Chemical
feedback, product inhibition, catalytic competition, compartment
permeability, phase transitions, and resource coupling can stabilize
some trajectories and suppress others.

Recent protocell models explicitly investigate how primitive regulation
or agency-like behavior can arise before contemporary cellular
machinery. Such models should be evaluated by intervention: does the
regulatory relation enlarge robust viability, restore a disturbed
variable, or enable survival across environments unavailable to an
unregulated control?

The term \emph{agency} should remain graded. A feedback-stabilized
protocell may instantiate minimal control without representing goals or
choosing among articulated policies.

\subsection{6.7 Heredity and the transition to
evolution}\label{heredity-and-the-transition-to-evolution}

Reproduction alone produces more units. Heredity makes differences
persist across generations. Evolvability requires variation that is
neither erased by low fidelity nor frozen by perfect copying, and whose
effects interact with selection and ecology.

Early heredity may have been compositional, catalytic, spatial, or
environmental before sequence-based genomes dominated. But a putative
inheritance system must be distinguished from mere persistence of the
same material. The relevant test is whether offspring systematically
resemble parents in organizationally consequential ways under controlled
perturbations.

Open-ended evolution is stronger still. It requires an expandable space
of heritable innovations, not merely convergence among a small set of
attractors.

\subsection{6.8 Why recurrence dominance is
insufficient}\label{why-recurrence-dominance-is-insufficient}

The former paper proposed a criterion of the form

\[\adjustbox{max width=\linewidth}{$\displaystyle 
P(L_{t+1}\midL_t)
>
P(L_{t+1}\mid\neg L_t)
$}\]

as the abiogenesis transition: once a life-like state exists, another
life-like state becomes more likely.

This inequality measures persistence or recurrence, not
self-maintenance. A crystal, chemical equilibrium, or passive stable
attractor can satisfy it. Consider a two-state Markov chain with states
\(L\) and \(N\) and transition matrix

\[\adjustbox{max width=\linewidth}{$\displaystyle 
P=
\begin{pmatrix}
0.99&0.01\\
0.10&0.90
\end{pmatrix},
$}\]

where the first row is conditioned on \(L\). Then
\(P(L_{t+1}\midL_t)=0.99>0.10=P(L_{t+1}\midN_t)\). Nothing in the
inequality states that \(L\) constructs its boundary, acquires energy,
repairs damage, or reproduces with heredity.

The criterion can be retained as a persistence statistic inside a larger
assay. It cannot carry the definition of life.

\subsection{6.9 A counterfactual autonomy
assay}\label{a-counterfactual-autonomy-assay}

A more discriminating origin-of-life experiment compares four
conditions:

\begin{enumerate}
\def\labelenumi{\arabic{enumi}.}
\tightlist
\item
  intact candidate system;
\item
  boundary-matched passive system;
\item
  externally controlled system receiving the same corrective actions;
\item
  candidate system with one putative closure link disabled.
\end{enumerate}

Expose all four to a declared perturbation distribution and measure:

\begin{itemize}
\tightlist
\item
  time within the viability region;
\item
  probability of recovery;
\item
  maintenance of boundary and reaction network;
\item
  reproduction rate and offspring similarity;
\item
  energy and material costs;
\item
  dependence on experimenter intervention.
\end{itemize}

The candidate shows increased autonomy when its internal organization
enlarges the robust viability kernel and when disabling reciprocal
maintenance links selectively removes that advantage.

This assay does not require a universal life score. It asks whether the
claimed organization does causal work.

\subsection{6.10 Intermediate stages rather than one
miracle}\label{intermediate-stages-rather-than-one-miracle}

Geochemical and prebiotic research increasingly studies sequences of
intermediate stages: catalytic mineral environments, carbon fixation or
proto-metabolic networks, compartments, activated monomers, template
systems, and coupled cycles. This is methodologically preferable to
requiring one experiment to jump from simple chemistry to a modern cell.

An autonomy-transition account predicts that early systems were mosaics
of internally generated and environmentally scaffolded functions.
Hydrothermal settings, mineral catalysts, wet-dry cycles, or ice phases
may have supplied constraints later internalized by cells. Evolution can
be understood partly as progressive internalization, coordination, and
replacement of such scaffolds---but not as a guaranteed march toward
independence. Modern organisms remain environmentally dependent.

\subsection{6.11 Phase-transition language: when it is
warranted}\label{phase-transition-language-when-it-is-warranted}

A specific model may exhibit a phase transition. To justify that term,
it should identify:

\begin{itemize}
\tightlist
\item
  a control parameter, such as catalytic density, membrane permeability,
  resource flux, or copying fidelity;
\item
  an order parameter, such as closure probability, sustained
  reproduction, or robust-viability volume;
\item
  system-size behavior;
\item
  discontinuity, critical scaling, bifurcation, hysteresis, or another
  defined transition signature; and
\item
  the relation between the model transition and a measurable
  experimental system.
\end{itemize}

The broad historical emergence of life need not map to one universal
critical point. Different capacities may have different transitions, and
ecological feedback can produce hysteresis or path dependence.

\subsection{6.12 Relation to algorithmic and informational
accounts}\label{relation-to-algorithmic-and-informational-accounts}

Some origin-of-life theories emphasize information becoming causally
efficacious, algorithmic control, or a shift from chemistry governed
only by local kinetics to systems in which stored constraints influence
future transformations. This paper is compatible with this family of
proposals when ``information'' is physically instantiated and
interventionally linked to maintenance (Walker \& Davies, 2013).

It adds two cautions.

First, causal influence alone is not enough; a random contaminant can
alter a reaction. The influence must participate in maintaining or
reproducing an organization.

Second, information does not float free of thermodynamics. Copying,
expression, and error correction are physical processes with costs and
failure modes.

\subsection{6.13 Biology Principle CU-B1}\label{biology-principle-cu-b1}

\begin{quote}
\textbf{CU-B1 --- Active Biological Self-Maintenance}
\textbf{{[}Domain-authority bridge principle{]}}\\
A candidate living organization is not distinguished by static order or
persistence alone, but by a relatively closed network of materially
realized processes and constraints whose endogenous activity contributes
to maintaining, repairing, reproducing, or adaptively transforming the
conditions of that network's continuation. The degree and kind of
autonomy must be assessed relative to a boundary, environment,
perturbation class, and horizon.
\end{quote}

\subsection{6.14 What would falsify the autonomy-transition
account}\label{what-would-falsify-the-autonomy-transition-account}

The account would lose value if:

\begin{itemize}
\tightlist
\item
  passive persistence predicted the target phenomena as well as the
  proposed maintenance variables;
\item
  disabling alleged closure links had no selective effect on viability
  or reproduction;
\item
  the dimensions of \(\mathbf A_P\) could not be manipulated
  independently even in principle;
\item
  claims of autonomy depended entirely on hidden experimenter support;
\item
  no operational boundary separated self-produced regulation from
  external control; or
\item
  the framework classified every sufficiently complicated chemical
  system as partially alive without yielding discriminating predictions.
\end{itemize}

\begin{center}\rule{0.5\linewidth}{0.5pt}\end{center}

\section{Section 7. Development, morphogenesis, and
regeneration}\label{section-7.-development-morphogenesis-and-regeneration}

\subsection{7.1 The problem of reliable
form}\label{the-problem-of-reliable-form}

A multicellular organism develops from local interactions among cells,
extracellular matrices, chemical signals, mechanical forces, electrical
states, and inherited structures. No cell contains a literal miniature
of the whole organism. Yet development often produces reproducible
anatomy despite noise, variable growth, cell turnover, and injury.

The explanatory challenge is not to choose between genes and ``the
whole.'' Genes, proteins, mechanics, bioelectric signaling, metabolism,
and tissue geometry participate in coupled control systems. The question
is how these variables jointly restrict the space of developmental
trajectories and how their relative causal contributions can be
distinguished.

This paper describes morphogenesis as \textbf{multiscale regulation of
form under viability and developmental constraints}. This description is
compatible with gene-regulatory networks, reaction--diffusion systems,
positional information, mechanical feedback, tissue-level
electrophysiology, and hybrid models. It does not postulate a separate
morphogenetic substance.

\subsection{7.2 Developmental state and target
sets}\label{developmental-state-and-target-sets}

Let the developing system have state

\[\adjustbox{max width=\linewidth}{$\displaystyle 
z_t=(g_t,s_t,m_t,v_t,h_t),
$}\]

where the components may represent gene-regulatory activity, biochemical
signaling, mechanical state, bioelectric state, and history-dependent
variables. A developmental model specifies

\[\adjustbox{max width=\linewidth}{$\displaystyle 
\dot z = F_{\mathrm{dev}}(z,e,u;\theta),
$}\]

with parameters \(\theta\), environmental inputs \(e\), and possible
interventions \(u\).

The mature form need not be one point. Define an admissible target set

\[\adjustbox{max width=\linewidth}{$\displaystyle 
\mathcal T_{\mathrm{form}}
=
\{z:\text{anatomy and function fall within the declared viable class}\}.
$}\]

Developmental success over horizon \(T\) can then be evaluated by
probability of reaching \(\mathcal T_{\mathrm{form}}\), time to reach
it, path cost, residual defects, and robustness to perturbation.

Target-set language is useful because organisms tolerate variation. Two
animals can differ in size or microscopic detail while realizing the
same broad anatomy and function.

\subsection{7.3 A target is not a hidden
blueprint}\label{a-target-is-not-a-hidden-blueprint}

A target set can be implemented by distributed feedback without being
represented in one location. The term \emph{target morphology} is
warranted when a model predicts error-sensitive, many-to-one convergence
toward a class of forms and when interventions can move, erase, or
replace the apparent setpoint.

It is not warranted merely because normal development is reproducible. A
feed-forward program can also yield a reliable result. To distinguish
control from fixed execution, experiments should test:

\begin{itemize}
\tightlist
\item
  perturbations not encountered during normal development;
\item
  different starting configurations converging on similar outcomes;
\item
  correction whose direction depends on the sign of the error;
\item
  reallocation of cells or tissue after damage;
\item
  persistent change in the recovered form after a transient
  intervention; and
\item
  restoration of the original behavior when the candidate control
  variable is reset.
\end{itemize}

A ``blueprint'' metaphor becomes misleading if it discourages analysis
of the actual distributed mechanism.

\subsection{7.4 Local mechanisms and global
constraints}\label{local-mechanisms-and-global-constraints}

A macrovariable is causally relevant when intervening on it---through
the local variables that realize it---changes local trajectories in ways
not captured by a model that omits the macrorelation.

For example, tissue curvature, stress, electric potential patterns, or
organ-scale geometry can constrain cell migration and differentiation.
This is not an extra force acting from outside physics. It is a compact
description of boundary conditions and coupled relations instantiated by
local interactions.

This paper uses the phrase \textbf{global constraint} only in this
interventionally disciplined sense. It does not infer causal power from
scale alone.

\subsection{7.5 Morphogenetic
integration}\label{morphogenetic-integration}

Let tissue modules be \(Z^{(1)},\dots,Z^{(n)}\). A maintenance-relevant
integration measure can compare developmental performance under intact
and selectively disrupted coupling:

\[\adjustbox{max width=\linewidth}{$\displaystyle 
I_{\mathrm{morph},T}
=
Q_T(F_{\mathrm{dev}})
-
Q_T(F_{\mathrm{dev}}^{\mathrm{cut}}),
$}\]

where \(Q_T\) may combine target attainment, defect correction, and
functional outcome.

The cut intervention should preserve local viability as closely as
possible. Otherwise a worse result may reflect nonspecific cell death
rather than loss of coordination.

This requirement is important for interpreting gap-junction blockers,
ion-channel modulators, matrix disruption, denervation, or mechanical
isolation. Each intervention has local effects in addition to network
effects.

\subsection{7.6 Planarian bioelectric signaling: what is
established}\label{planarian-bioelectric-signaling-what-is-established}

Planarian regeneration supplies a strong example of early physiological
signals affecting later anatomy. Durant and colleagues reported that
brief depolarization during the first three hours after amputation
altered later gene expression and produced a minority of double-headed
regenerates, even after the ionophores were washed out (Durant et al.,
2019).

This supports several bounded claims:

\begin{enumerate}
\def\labelenumi{\arabic{enumi}.}
\tightlist
\item
  resting-potential differences can function upstream of later
  transcriptional patterning in that preparation;
\item
  a transient physiological intervention can have persistent
  morphological consequences;
\item
  local genetic state alone is not a complete description of the
  relevant early control process; and
\item
  regeneration can display path dependence and alternative stable
  outcomes.
\end{enumerate}

It does \textbf{not} by itself prove that bioelectricity is the
universal master code of anatomy, that planaria store a conscious body
image, or that all development is controlled by one electrical
attractor. The effect size, species, treatment, timing, and off-target
pathways matter.

The relevant CU contribution is methodological: identify the state
variable, perturb it transiently, trace the causal path, and test
whether it changes the target distribution rather than merely damaging
tissue.

\subsection{7.7 Regeneration as controlled
re-entry}\label{regeneration-as-controlled-re-entry}

Regeneration differs from ordinary development because the initial
condition lies outside the normal embryonic trajectory. A fragment or
injured tissue must infer, through local and long-range signals, what
structures are missing and how to rebuild them without duplicating what
remains.

A model can represent regeneration as controlled entry into a viable
target set:

\[\adjustbox{max width=\linewidth}{$\displaystyle 
\rho^*
=
\arg\max_{\rho\in\Pi}
\Pr
\left[
 z_T^{\rho}
 \in
 \mathcal T_{\mathrm{form}}
 \mid
 z_0=z_{\mathrm{injured}}
\right],
$}\]

subject to metabolic and tissue constraints.

This notation does not imply that the organism explicitly solves an
optimization problem. It defines the explanatory performance of
regulatory policies.

A strong regeneration theory should predict not only success or failure,
but the class of abnormal stable forms reached after interventions and
the conditions under which the original anatomy is recovered.

\subsection{7.8 Robustness, plasticity, and
canalization}\label{robustness-plasticity-and-canalization}

Development is robust when outcomes remain within an admissible range
despite perturbation. It is plastic when conditions produce alternative
adaptive forms. It is canalized when trajectories converge strongly
despite variation.

These properties can conflict. Strong canalization may prevent useful
adaptation; high plasticity may increase developmental error; robustness
in one environment may produce fragility in another.

The coherence profile keeps these dimensions separate. It therefore
resists the temptation to identify ``more stable development'' with
``more coherent development'' without specifying whether flexibility and
future viability have been sacrificed.

\subsection{7.9 Mechanics as memory and
constraint}\label{mechanics-as-memory-and-constraint}

Tissue form influences future development through stresses, geometry,
adhesion, fluid pressure, and extracellular matrix. These physical
variables can store history: once a fold, boundary, or stiffness pattern
forms, it changes later signaling and motion.

Mechanical memory need not be symbolic. It is memory in the sense of
Section 3.15: a history-dependent state that changes present dynamics.

A complete multiscale model should test reciprocal causation:

\[\adjustbox{max width=\linewidth}{$\displaystyle 
\text{gene/signaling state}
\longrightarrow
\text{mechanics}
\longrightarrow
\text{gene/signaling state}.
$}\]

One-way models may be adequate for some interventions but should not be
universalized.

\subsection{7.10 Hybrid models}\label{hybrid-models}

The most promising formal direction is not one privileged field but
hybrid models coupling:

\begin{itemize}
\tightlist
\item
  gene-regulatory networks;
\item
  reaction--diffusion or transport fields;
\item
  cell mechanics and tissue geometry;
\item
  membrane-potential and gap-junction dynamics;
\item
  cell proliferation, death, migration, and differentiation; and
\item
  history-dependent state transitions.
\end{itemize}

Recent mathematical work continues to develop models of morphogenesis
that combine bioelectric, regulatory, adaptive-coupling, and memory
variables rather than treating any one layer as sufficient (Cortés-Poza,
2026). CU's role is to specify the cross-scale questions such models
should answer: which couplings enlarge the viable target basin, which
create alternative attractors, which repair perturbations, and what
costs or fragilities accompany them.

\subsection{7.11 Reconfigurable multicellular
constructs}\label{reconfigurable-multicellular-constructs}

Xenopus-derived reconfigurable organisms---often called xenobots---show
that living cells removed from their ordinary embryonic context can
self-assemble into designed multicellular forms with locomotor and
collective behaviors (Kriegman et al., 2020). Later work demonstrated a
particular form of kinematic replication in which aggregates moved loose
cells into offspring-like clusters under laboratory conditions (Kriegman
et al., 2021).

These results support the claim that cell collectives possess reusable
morphogenetic and behavioral capacities not exhausted by the anatomy
they normally build. They also illustrate the importance of external
scaffolding: cells, medium, design procedures, and supplied loose cells
form part of the experimental causal system.

The experiments do not establish autonomous open-ended reproduction, a
new natural species, consciousness, or a universal drive toward novel
forms. They are valuable contrast cases for measuring how much
organization is carried by cell-intrinsic capacities, collective
coupling, geometry, and experimenter support.

\subsection{7.12 Morphogenesis is not ``continuous abiogenesis'' without
qualification}\label{morphogenesis-is-not-continuous-abiogenesis-without-qualification}

Morphogenesis and abiogenesis share a limited structural similarity:
both involve the construction of organized boundaries and mutually
dependent processes from local interactions. The difference is decisive.
Development begins inside an inherited living organization containing
membranes, metabolism, genetic systems, maternal products, cellular
machinery, and an evolved environment. It does not recreate the
transition from chemistry to life. The comparison is therefore
interpretive:

\begin{quote}
Development repeatedly rebuilds organismal organization from cells and
inherited constraints; it does not repeat the historical origin of
biological autonomy from nonliving chemistry.
\end{quote}

\subsection{7.13 Biology Principle CU-B3}\label{biology-principle-cu-b3}

\begin{quote}
\textbf{CU-B3 --- Multiscale Constraint and Local Process}
\textbf{{[}Domain-authority bridge principle{]}}\\
Biological form and function arise through local physical processes
embedded in cross-scale boundary conditions, signaling networks,
mechanical relations, and histories. A higher-scale constraint is
biologically explanatory when intervention on its realizing relations
changes lower-scale trajectories in a way that improves prediction or
control beyond local variables alone. Higher scale does not imply a new
force, and global stability does not imply a unique optimizer.
\end{quote}

\subsection{7.14 Experimental predictions for development and
regeneration}\label{experimental-predictions-for-development-and-regeneration}

This paper proposes the following discriminating predictions.

\textbf{Prediction D1: coupling-specific basin change.} Selectively
disrupting a coordination channel while preserving local viability will
reduce the basin of successful target attainment or increase alternative
stable outcomes. A purely local independent-cell model predicts no
additional effect once local damage is controlled.

\textbf{Prediction D2: transient setpoint displacement.} If a field
variable participates in a persistent target-setting circuit, a timed
transient intervention will shift later morphology after the
intervention is removed. A feed-forward damage model predicts recovery
once local injury is absent, unless the intervention irreversibly
damaged another component.

\textbf{Prediction D3: distributed compensation.} Removal or
displacement of one region will change behavior in distant regions in a
direction that compensates for the global error. The null model predicts
only local responses and generic wound effects.

\textbf{Prediction D4: configuration-dependent coupling surface.}
Increasing coupling may improve correction when a threatened component
lies below the coupled mean, but it may drain a specialized
high-requirement component that lies above that mean. It can also impair
adaptation to spatially asymmetric environments or raise energetic cost.
Supplement B supplies two toy configurations with opposite viability
trends; empirical systems must establish their own surface.

Failure of these predictions in well-controlled systems would narrow
CU's multiscale-control claims. It would not invalidate ordinary
developmental biology.

\begin{center}\rule{0.5\linewidth}{0.5pt}\end{center}

\section{Section 8. Biological agency, learning, and minimal
epistemology}\label{section-8.-biological-agency-learning-and-minimal-epistemology}

\subsection{8.1 Why agency should be
graded}\label{why-agency-should-be-graded}

Biology contains directed behavior long before reflective choice.
Bacteria move up chemical gradients, cells repair membranes, plants
alter growth, immune systems update responses, and animals learn.
Calling all of these either ``mere mechanism'' or full agency obscures
the important differences.

This paper treats agency as a family of capacities by which an
organization detects or is differentially affected by conditions, acts
through available pathways, and changes the probability of maintaining
or transforming itself.

Agency is not a substance. It is an organization of control.

\subsection{8.2 An agency ladder}\label{an-agency-ladder}

The following levels are analytic distinctions, not a claim that
evolution followed one linear ladder.

\subsubsection{A0 --- Passive persistence}\label{a0-passive-persistence}

A structure remains because its physical configuration is stable under
current conditions. It has no demonstrated error-sensitive corrective
action.

\subsubsection{A1 --- Homeostatic
regulation}\label{a1-homeostatic-regulation}

The system detects or is coupled to deviations from a range and acts so
as to restore a variable. Negative feedback, integral feedback, and
buffering belong here.

\subsubsection{A2 --- Sensorimotor
control}\label{a2-sensorimotor-control}

The system alters action in response to environmental signals so as to
preserve viability or reach a functional target. Bacterial chemotaxis is
a paradigmatic minimal case.

\subsubsection{A3 --- Adaptive reconfiguration and
learning}\label{a3-adaptive-reconfiguration-and-learning}

Past interaction changes the control policy, sensitivity, or internal
organization in a way that improves performance under a defined task or
perturbation distribution. Habituation, sensitization, associative
change, and some immune adaptations can fall here when controls exclude
fatigue or damage.

\subsubsection{A4 --- Anticipatory or model-based
control}\label{a4-anticipatory-or-model-based-control}

The system uses internal states that track counterfactual or future
conditions and selects among actions on that basis. This can occur
without language or reflective awareness.

\subsubsection{A5 --- Reflective agency}\label{a5-reflective-agency}

The system represents itself as an agent, can compare reasons or
policies, and can revise commitments. This belongs primarily to the
Psychology, Consciousness, and Ethics domains.

An entity can display different levels in different capacities.

\subsection{8.3 Goal-directedness as a testable
attribution}\label{goal-directedness-as-a-testable-attribution}

A system appears goal-directed when different disturbances elicit
different actions that converge on a common outcome. But equifinality
alone can arise from passive attractor dynamics.

A stronger attribution uses four criteria:

\begin{enumerate}
\def\labelenumi{\arabic{enumi}.}
\tightlist
\item
  \textbf{error sensitivity:} action depends on the direction and
  magnitude of deviation;
\item
  \textbf{counterfactual flexibility:} the system reaches the target
  through different means when usual pathways are blocked;
\item
  \textbf{resource allocation:} the system expends limited resources in
  proportion to the protected variable's importance;
\item
  \textbf{reorganization:} repeated perturbation can alter the policy
  while preserving the broader target.
\end{enumerate}

No single criterion proves representation or consciousness. Together
they distinguish increasingly rich control from passive relaxation.

\subsection{8.4 Homeostasis and integral
feedback}\label{homeostasis-and-integral-feedback}

Control theory provides exact mechanisms for robust adaptation. Integral
feedback can restore a regulated output to a setpoint despite persistent
disturbances, provided stability and structural assumptions hold.

Antithetic integral feedback is one biomolecular architecture in which
two controller species annihilate each other and implement an
integral-like error signal; synthetic biology experiments have
demonstrated robust perfect adaptation under specified conditions (Aoki
et al., 2019).

This result shows why CU should not appeal to a mysterious ``coherence
drive.'' Directed restoration can arise from identifiable feedback
circuits. CU-B1 names the organizational role of those circuits; it does
not replace their mechanism.

\subsection{8.5 Limits of control}\label{limits-of-control}

Every controller faces constraints. Sensing is noisy; actions cost
energy; delays can destabilize; high gain can amplify noise; the
environment can change faster than adaptation; and one controlled
variable can conflict with another.

Work on bacterial chemotaxis has analyzed tradeoffs and limits in
feedback control, including the relation among response, adaptation, and
noise (Dufour et al., 2014). These limits are not imperfections added to
an ideal organism. They partly define the feasible agency of the system.

Accordingly, biological coherence is not perfect return to a setpoint.
It is viable navigation under bounded information and control.

\subsection{8.6 Learning in single
cells}\label{learning-in-single-cells}

Single-cell learning claims require careful controls because reduced
response can result from fatigue, receptor adaptation, toxicity, or
motor exhaustion rather than memory.

Experiments on \emph{Stentor coeruleus} have provided molecular and
behavioral evidence relevant to habituation-like learning and its
pathways. In a 2026 study, cycloheximide and puromycin accelerated
habituation and prolonged memory retention rather than disrupting it in
the familiar metazoan pattern. The study also reported maintenance of
habituation memory in progeny following cell division. Its
interpretation is that new protein synthesis contributes to recovery of
responsiveness, while memory formation involves modification of
existing, spatially distributed mechanoreceptor machinery through
calcium signaling and phosphorylation (Rajan et al., 2026).

These findings strengthen the case for history-dependent regulation in a
non-neural single cell. They do not show that the cell uses a metazoan
memory mechanism, and the inheritance result should not be inflated into
a general theory of transgenerational cognition.

The bounded conclusion is significant:

\begin{quote}
Neural tissue is not necessary for every operational form of learning.
\end{quote}

It does not follow that all cells learn, that \emph{Stentor} has
human-like concepts, or that learning establishes consciousness. The
result belongs at A3 in the agency ladder unless further evidence
supports stronger claims.

\subsection{8.7 Novel problem solving}\label{novel-problem-solving}

Biological systems can sometimes solve challenges outside the specific
historical distribution for which their present response architecture
was shaped. That claim requires an operational definition of novelty
rather than an appeal to foresight.

A system displays \textbf{novel adaptive problem solving} when it
succeeds on a challenge outside the training, developmental, or
historical distribution relevant to the claimed mechanism, and when
success depends on flexible recombination or reconfiguration rather than
a pre-specified response.

Demonstrating novelty requires evidence about the relevant prior
distribution. No organism is literally without evolutionary history. An
apparently novel behavior can recruit ancient generic
capacities---exploration, plasticity, adhesion, chemotaxis, or stress
responses.

The appropriate question is not whether evolution anticipated the exact
problem. It is whether evolution produced architectures with sufficient
adaptive reach to solve classes of problems not encoded as fixed
responses.

\subsection{8.8 Distributed problem
solving}\label{distributed-problem-solving}

In multicellular systems, adaptive behavior can be distributed across
cells. Immune coordination, wound healing, vascular remodeling,
collective migration, and tissue repair involve local sensing and
communication without one central controller.

Distributed control can outperform centralized control through
parallelism and robustness, but it also creates coordination problems:
conflicting local objectives, delayed signals, false positives, and
runaway positive feedback.

The phrase ``the tissue solves a problem'' is legitimate when the
collective exhibits target-relative error correction not reducible to
independent local responses under the tested model. It should not imply
that the tissue deliberates like a person.

\subsection{8.9 Minimal biological
epistemology}\label{minimal-biological-epistemology}

A signal is not meaningful merely because it correlates with an
environmental state. It acquires \textbf{biological significance} when
the organization uses or is shaped by the signal in a way that changes
viability-relevant action.

Let environmental condition \(e\) influence signal \(s\), and let policy
\(\rho(a\mid s,h)\) determine action. A minimal value of information can
be written

\[\adjustbox{max width=\linewidth}{$\displaystyle 
\mathcal V_T(s)
=
\sup_{\rho\in\Pi_s}Q_T(\rho)
-
\sup_{\rho\in\Pi_0}Q_T(\rho),
$}\]

where \(\Pi_s\) contains policies able to use \(s\) and \(\Pi_0\)
contains matched policies denied it. If \(\mathcal V_T(s)>0\), the
signal has instrumental value for the declared outcome and environment
distribution.

This is not semantic truth in the human propositional sense. It is a
minimal relation among environmental regularity, internal state, action,
and organism-relative stake.

\subsection{8.10 Error without
propositions}\label{error-without-propositions}

A bacterium can move down rather than up a nutrient gradient. Calling
that an ``error'' is justified when the action reduces expected
viability relative to accessible alternatives under the organism's own
regulatory organization.

This is normative but not moral. The norm arises from the system's
viability conditions and policy, not from an external command.

Formally, action \(a_t\) is control-suboptimal relative to model \(M\)
when

\[\adjustbox{max width=\linewidth}{$\displaystyle 
Q_T(a_t\mid x_t,M)
<
\sup_{a\in U(x_t)}Q_T(a\mid x_t,M).
$}\]

Because the model and outcome are analyst-specified, this classification
should be tested against behavior and intervention. It is not an
infallible reading of what the organism ``really wants.''

\subsection{8.11 Information, misinformation, and
deception}\label{information-misinformation-and-deception}

At the biological level, a signal can be misleading when it normally
tracks a viability-relevant condition but is decoupled from it, causing
maladaptive action. Mimicry, sensory traps, and pathogen manipulation
exploit such relations.

Deception in the richer psychological or moral sense requires more:
representation of another agent's informational state, intention or
policy, and often reasons. The Biology paper should not project those
properties onto every signaling interaction.

\subsubsection{Regulatory correctability and
misregistration}\label{regulatory-correctability-and-misregistration}

Biological regulation can lose contact with wider support conditions
when exported costs, delayed consequences, or scaffold dependence fall
outside effective feedback. A lineage, tissue, or organism may improve a
locally registered variable while degrading the ecological or
organism-level conditions that sustain it. This is \textbf{regulatory
misregistration}, not self-deception unless representational capacity
and motivated exclusion are independently established. The correction is
to widen the boundary, observation channels, and horizon rather than
moralize the organism.

\subsection{8.12 The free-energy principle as a model
family}\label{the-free-energy-principle-as-a-model-family}

The free-energy principle and active inference offer mathematical
accounts of systems that maintain states within characteristic bounds by
minimizing variational free energy under a generative model. They
connect perception, action, and homeostasis in a powerful formal
language (Friston, 2013).

This paper treats this as one model family, not as a theorem that every
living system literally performs the same optimization. The biological
content depends on identifying variables, blanket structure, generative
model, action channels, timescales, and empirical predictions that
distinguish the account from alternative control descriptions.

The coherence profile can be instantiated using active-inference
quantities where appropriate. It does not assume their universal
sufficiency.

\subsection{8.13 Biology Principle CU-B2}\label{biology-principle-cu-b2}

\begin{quote}
\textbf{CU-B2 --- Adaptive Problem Solving Beyond Fixed Programs}
\textbf{{[}Domain-authority conditional empirical claim{]}}\\
Some living systems possess regulatory architectures that generalize
beyond fixed stimulus--response mappings: they detect deviations,
recruit alternative pathways, modify future responses from history, and
solve novel viability-relevant problems within a bounded adaptive reach.
The degree of agency must be established through counterfactual
flexibility, learning controls, and comparison with passive or
preprogrammed alternatives; it does not by itself establish
consciousness.
\end{quote}

\subsection{8.14 Transition to
consciousness}\label{transition-to-consciousness}

Biology supplies embodiment, temporality, self-maintenance,
organism-relative stake, and many forms of integration. These may be
important determinants of the depth, unity, and valence of
consciousness. But this paper does not state that all living systems are
conscious or that life is necessary for every possible perspective.

Consciousness v0.10 treats access organization, perspectival anchoring,
and center individuation as separable evidential questions.
Nonbiological systems may realize thin or episodic perspective through
global-workspace-like routes, while biological consciousness commonly
adds embodiment, recurrent temporal organization, organism-relative
stake, and valenced regulation. Biology therefore supplies a comparative
baseline, not an exclusion theorem.

\begin{center}\rule{0.5\linewidth}{0.5pt}\end{center}

\section{Section 9. Evolution, structured possibility, and major
transitions}\label{section-9.-evolution-structured-possibility-and-major-transitions}

\subsection{9.1 Evolutionary theory remains
indispensable}\label{evolutionary-theory-remains-indispensable}

Coherence Universalism does not replace natural selection, mutation,
recombination, drift, migration, developmental bias, niche construction,
or ecological interaction. It asks a narrower integrative question:

\begin{quote}
What kinds of regulatory and organizational architecture become
heritable units, and under what conditions do evolutionary processes
preserve, elaborate, simplify, or dismantle them?
\end{quote}

The answer cannot be ``evolution maximizes coherence.'' Evolution has no
universal foresight and no single organism-level objective.

\subsection{9.2 Fitness is not
coherence}\label{fitness-is-not-coherence}

Fitness is defined relative to population, environment, trait, and
timescale, commonly through expected reproductive contribution.
Biological coherence concerns the organization by which a candidate unit
maintains and transforms itself.

The two can correlate, but neither entails the other.

A fragile organism can have high fitness in a stable environment if it
reproduces quickly. A robust sterile organism can maintain itself
without contributing descendants. A parasite can gain fitness by
degrading host organization. A tumor clone can be locally successful
while the organism dies. A lineage can increase short-term fitness by
exhausting an environmental resource.

Therefore:

\[\adjustbox{max width=\linewidth}{$\displaystyle 
W\ne C_B,
$}\]

where \(W\) denotes any standard fitness measure.

\subsection{9.3 Selection can stabilize
self-maintenance}\label{selection-can-stabilize-self-maintenance}

Although fitness and coherence are distinct, selection often favors
mechanisms that improve maintenance because organisms unable to survive
long enough, acquire resources, repair, or reproduce leave fewer
descendants.

Candidate traits include:

\begin{itemize}
\tightlist
\item
  homeostatic feedback;
\item
  proofreading and repair;
\item
  modularity;
\item
  redundancy;
\item
  stress response;
\item
  developmental canalization;
\item
  phenotypic plasticity;
\item
  immune defense;
\item
  conflict suppression; and
\item
  ecological construction.
\end{itemize}

Their advantage is conditional. Redundancy can be costly; repair can
slow reproduction; plasticity can misfire; immune activity can damage
the host. Evolution tunes tradeoffs under local conditions rather than
converging on maximal global organization.

\subsection{9.4 Structured evolutionary
possibility}\label{structured-evolutionary-possibility}

Evolution does not search all conceivable organisms with equal
probability. At any point in a lineage's history, only a restricted and
highly nonuniform region of genotype and phenotype space is reachable.

Mutation biases affect which genetic changes arise. Genotype--phenotype
maps make some phenotypes vastly easier to produce than others.
Developmental dynamics channel perturbations toward some forms and away
from others. Physical and chemical law exclude many imaginable
organizations. Ecological relations alter which variants survive and
reproduce. Prior evolution determines the architecture from which
subsequent variation begins.

Coherence Universalism calls this \textbf{structured evolutionary
possibility}.

\begin{quote}
\textbf{{[}Definition / synthesis{]}} Structured evolutionary
possibility is the nonuniform organization of reachable heritable
transformations generated by the interaction of variation,
genotype--phenotype relations, development, physical feasibility,
ecology, selection, drift, inheritance, and history.
\end{quote}

The phrase does not name a force in addition to those processes. It
names the organized transition geometry their interaction produces.

A schematic representation is:

\[\adjustbox{max width=\linewidth}{$\displaystyle 
(g,h,e)
\xrightarrow{\;\mathcal V\;}
\Delta g
\xrightarrow{\;\mathcal G_{h,e}\;}
p
\xrightarrow{\;\mathcal E_e\;}
W,
$}\]

where \(g\) is genotype, \(h\) inherited and developmental history,
\(e\) environment, \(\mathcal V\) the production of variation,
\(\mathcal G_{h,e}\) the history- and environment-dependent
genotype--phenotype map, and \(\mathcal E_e\) the ecological and
selective relation between phenotype and reproductive consequence.

Every arrow can be biased, constrained, many-to-one, one-to-many, or
historically altered. The possibility space is therefore neither blank
nor fixed.

\subsection{9.5 Mutation and genotype--phenotype
bias}\label{mutation-and-genotypephenotype-bias}

Mutation has stochastic components, but variation is not introduced
uniformly. Some substitutions occur more often than others, and some
phenotypes are encoded by far more genotypes or reached from far more
mutational neighborhoods.

Schaper and Louis (2014) showed in an RNA genotype--phenotype model that
frequently generated phenotypes can fix even when rarer alternatives
would have higher fitness. Their ``arrival of the frequent'' result
demonstrates that the distribution of produced variation can steer
evolution before selection acts among the variants that actually appear.

Herrera-Álvarez, Patton, and Thornton (2025) reconstructed ancestral
transcription-factor genotype--phenotype maps and found them strongly
anisotropic and heterogeneous. Some phenotypes were much more
generatively accessible than others, and local accessibility varied
substantially across starting genotypes. In both reconstructed maps,
those inherited structures causally steered evolution toward
lineage-specific phenotypes that later evolved historically.

The implication is not that selection is secondary or irrelevant. It is
that selection acts on a produced distribution of variants whose
structure already reflects chemistry, molecular architecture, ancestry,
and development.

\subsection{9.6 Developmental attractors, bias, and lines of least
resistance}\label{developmental-attractors-bias-and-lines-of-least-resistance}

Developmental systems do not translate genotype into phenotype like a
simple lookup table. They are dynamical systems involving regulatory,
biochemical, mechanical, bioelectric, and environmental relations.

In some cases, multiple initial states or perturbations converge toward
a stable phenotypic regime. These are \textbf{developmental attractors}:
state-space regions toward which trajectories of a declared
developmental system converge. Huang et al.~(2005), for example,
reported gene-expression trajectories in neutrophil differentiation
converging toward a common high-dimensional state, supporting an
attractor interpretation for that system.

Developmental attractors matter evolutionarily because they can make
some phenotypes robustly generable and other mathematically imaginable
configurations difficult or impossible to realize.

Attractor language must remain model-specific. A stable cell fate is not
the same object as a population-level adaptive peak. The word
\emph{attractor} should identify a state-space and dynamics, not
function as a synonym for ``common outcome.''

Development can also bias directions of variation without producing one
fixed endpoint. Rohner, Hu, and Moczek (2022) found that perturbations
of beetle horn development, environmental plasticity, population
differences, and macroevolutionary divergence aligned along shared
dimensions, supporting developmental ``lines of least resistance'' in
that system.

Developmental bias is not immutable. It can itself evolve, interact with
selection, or be overcome. The bounded conclusion is that selection acts
on variation already shaped by how organisms are built.

\subsection{9.7 Evolutionary basins and adaptive
dynamics}\label{evolutionary-basins-and-adaptive-dynamics}

Population-level evolutionary models may contain local adaptive regions
or evolutionarily stable attractors under declared assumptions about
inheritance, ecology, frequency dependence, and available variation.

These \textbf{evolutionary or adaptive basins} are not identical to
developmental attractors. One concerns change in an organismal
developmental state; the other concerns trajectories of populations,
strategies, or trait distributions.

A lineage may enter a basin because nearby variation repeatedly
increases fitness under a stable environment. It may leave if the
environment changes, new variation becomes accessible, another species
alters the ecology, or the lineage crosses a historically contingent
threshold.

The CU claim is therefore conditional:

\begin{quote}
Some evolutionary models and empirical systems exhibit basin-like
dynamics, but no single adaptive landscape governs all lineages, scales,
or times.
\end{quote}

\subsection{9.8 Convergence as evidence of
structure}\label{convergence-as-evidence-of-structure}

Convergent evolution occurs when independent lineages repeatedly produce
similar structures, functions, or ecological roles.

Losos et al.~(1998) showed that similar \emph{Anolis} ecomorphs evolved
independently on different Greater Antillean islands. Travisano et
al.~(1995) found that replicate bacterial populations in identical
environments converged strongly in fitness while remaining more
divergent in cell size, showing that adaptation, chance, and history can
dominate different traits. Simões et al.~(2019) found broad convergence
during experimental adaptation of \emph{Drosophila} populations while
several traits retained differences and some trajectories showed
transient divergence.

At the molecular level, mutation bias can contribute to repeated
adaptive substitutions. Storz et al.~(2019) found that causative
amino-acid changes associated with convergent high-altitude hemoglobin
adaptation in birds were disproportionately attributable to mutationally
elevated CpG sites.

Convergence therefore reveals that the feasible design space has a
grain. Its causes can include:

\begin{itemize}
\tightlist
\item
  shared selection pressures;
\item
  physical and chemical constraints;
\item
  mutation bias;
\item
  genotype--phenotype bias;
\item
  developmental architecture;
\item
  inherited pathway structure;
\item
  limited numbers of workable solutions; and
\item
  ecological interaction.
\end{itemize}

Each case requires causal discrimination. Convergence does not by itself
show that a lineage followed a literal attractor or that the outcome was
inevitable.

The strongest CU formulation is:

\begin{quote}
When independent lineages face sufficiently similar constraints and
possess access to comparable developmental possibilities, repeated
solutions disclose structure in the feasible design space.
\end{quote}

\subsection{9.9 Historical contingency and changing
accessibility}\label{historical-contingency-and-changing-accessibility}

Structured possibility does not eliminate history. History partly
determines which possibilities are reachable.

Blount, Borland, and Lenski (2008) showed that the evolution of aerobic
citrate use in one long-term \emph{E. coli} population depended on prior
potentiating changes. The relevant innovation was not simply waiting
with a constant probability from the ancestral state; earlier history
altered the lineage's capacity to produce or express it.

This suggests a useful distinction:

\[\adjustbox{max width=\linewidth}{$\displaystyle 
\text{global possibility}
\neq
\text{current accessibility}.
$}\]

A phenotype may be physically and genetically possible in a broad sense
while inaccessible from a particular lineage without prior enabling
changes. Conversely, an ancestral architecture can make some outcomes
recurrently accessible even when many alternatives are imaginable.

The coherence of evolution lies partly in this historical organization
of possibility.

\subsection{9.10 Canalization and
evolvability}\label{canalization-and-evolvability}

Canalization reduces phenotypic variation under perturbation;
evolvability concerns the capacity to generate heritable adaptive
variation. They can support or oppose each other.

Modular architecture may permit local change without system-wide
failure. Robustness can allow cryptic variation to accumulate. But
excessive canalization can trap a lineage, and excessive flexibility can
destroy reliable development.

This paper predicts tradeoff surfaces rather than monotonic ascent.

\subsection{9.11 Major evolutionary
transitions}\label{major-evolutionary-transitions}

Major transitions occur when previously independent replicating or
reproducing units form a new higher-level unit---genes into genomes,
cells into multicellular organisms, or organisms into highly integrated
collectives. The transition is not mere aggregation. It usually
involves:

\begin{itemize}
\tightlist
\item
  cooperation among lower-level units;
\item
  division of labor;
\item
  increased mutual dependence;
\item
  restriction or coordination of lower-level reproduction;
\item
  mechanisms suppressing conflict;
\item
  a collective life cycle or bottleneck;
\item
  heritable collective variation; and
\item
  relocation of fitness-relevant causal organization.
\end{itemize}

These dimensions correspond closely to the individuality profile in
Section 4. Major transitions are therefore changes in where biological
stake and evolutionary individuality are concentrated.

\subsection{9.12 Conflict mediation is constitutive, not
optional}\label{conflict-mediation-is-constitutive-not-optional}

Higher-level integration is unstable when lower-level units can gain by
defecting. Policing, partner choice, germ-soma separation, apoptosis,
immune surveillance, fair meiosis, and reproductive bottlenecks help
align interests.

This gives a precise version of the local--global principle:

\begin{quote}
Collective organization becomes evolutionarily stable only when the
benefits of lower-level defection are sufficiently constrained,
compensated, or redirected relative to collective reproduction.
\end{quote}

The principle is not exceptionless. Conflict persists inside organisms,
and some systems remain stable despite it. What matters is whether
conflict remains below the level that destroys the higher-level life
cycle.

\subsection{9.13 Ecological scaffolding and
hysteresis}\label{ecological-scaffolding-and-hysteresis}

Transitions in individuality can depend on ecological conditions that
temporarily align lower-level interests or concentrate offspring. Once
collective adaptations evolve, the collective may persist after the
original scaffold changes. This creates path dependence and possible
hysteresis.

Recent theoretical work on evolutionary transitions has shown how
ecological scaffolding can stabilize collective traits through
evolutionary hysteresis, so that some traits remain after initiating
conditions are relaxed (Doulcier et al., 2024). This paper treats this
as a promising realization of historical continuity and changing
individuation, not as evidence for a predetermined ladder.

\subsection{9.14 Evolutionary memory}\label{evolutionary-memory}

Genetic inheritance is one major memory channel, but developmental
systems, parental effects, ecological modifications, symbionts, and
cultural transmission can also carry history.

An inheritance channel counts as evolutionary memory when differences
persist across generations and causally affect phenotype or fitness. The
term should not erase distinctions among DNA sequence, epigenetic state,
constructed environment, and learned tradition.

This paper therefore uses a channel-indexed representation:

\[\adjustbox{max width=\linewidth}{$\displaystyle 
H_{n+1}^{(k)}
=
\mathcal F_k
\left(
H_n^{(k)},
P_n,
E_n
\right),
$}\]

where \(k\) identifies the inheritance mechanism.

\subsection{9.15 Evolution does not guarantee complexity or
progress}\label{evolution-does-not-guarantee-complexity-or-progress}

Lineages can simplify, lose functions, become parasitic, enter
evolutionary dead ends, or go extinct. Natural selection optimizes no
universal future. Complexity can increase in some lineages and decrease
in others.

The Coherence Ladder is therefore not a law of evolutionary history. It
is a dependency map: certain capacities require enabling structures.
Evolution may discover some routes, bypass others through different
mechanisms, or never enter the relevant region.

Any spiritual narrative of the universe moving toward consciousness
belongs to the commitment register and cannot be cited as a biological
result.

\subsection{9.16 Multilevel assessment}\label{multilevel-assessment}

For a nested system with levels \(\ell=1,\dots,L\), let \(W_\ell\) and
\(\mathbf C_\ell\) denote fitness and coherence profiles at each level.
Evolutionary change can satisfy

\[\adjustbox{max width=\linewidth}{$\displaystyle 
\Delta W_i>0,
\qquad
\Delta \mathbf C_j<0
$}\]

for another level \(j\). There is no universal conservation or
monotonicity law.

A complete analysis should report:

\begin{itemize}
\tightlist
\item
  the reproducing units;
\item
  the inheritance channel;
\item
  the selection level or decomposition;
\item
  within- and between-group variation;
\item
  the horizon;
\item
  and the ecological costs shifted to other units.
\end{itemize}

This prevents organism-centered language from hiding selection at other
levels.

\subsection{9.17 Biology Principle CU-B4}\label{biology-principle-cu-b4}

\begin{quote}
\textbf{CU-B4 --- Structured Evolutionary Possibility and the
Stabilization of Regulatory Architecture} \textbf{{[}Domain-authority
conditional empirical generalization{]}}\\
Evolutionary processes explore a structured and nonuniform possibility
space shaped by mutation bias, genotype--phenotype relations,
developmental dynamics, physical feasibility, ecological interaction,
inheritance, and historical path dependence. Selection nonrandomly
changes the frequencies of heritable variants relative to local
reproductive consequences. Recurrent constraints can make some forms,
functions, and regulatory architectures repeatedly accessible,
contributing to convergence, while chance, drift, conflict, changing
environments, and historical contingency preserve multiple possible
trajectories. Evolution can stabilize, modify, simplify, or dismantle
coherent self-maintaining organization; it does not thereby maximize
coherence, complexity, consciousness, or moral value.
\end{quote}

\subsection{9.18 Candidate predictions}\label{candidate-predictions}

\textbf{Prediction E1: architecture-mediated generalization.} Lineages
selected under variable but structured environments will evolve greater
adaptive reach or modular control than lineages selected under equally
harsh but uncorrelated environments, after matching average fitness. A
simple stress-resistance account predicts only generalized hardening.

\textbf{Prediction E2: conflict-control threshold.} In experimental
transitions to multicellularity, collective-level heritability and
persistence will rise only when mechanisms reduce within-collective
conflict or tie lower-level reproduction to collective success. The null
model predicts aggregation alone is sufficient.

\textbf{Prediction E3: robustness--innovation tradeoff.} Perturbations
that increase canalization will improve target fidelity under familiar
conditions while reducing exploration of novel viable phenotypes beyond
some range. A monotonic coherence-maximization account would not predict
the turnover.

\textbf{Prediction E4: history-dependent reversal.} Some integrated
collectives will persist under conditions in which de novo formation is
unlikely, demonstrating hysteresis. A memoryless equilibrium model
predicts the same transition point in both directions.

\textbf{Prediction E5: accessibility predicts recurrent outcomes.} After
controlling for measured fitness effects, phenotypes with greater
mutational or developmental accessibility will arise earlier or more
frequently across replicate populations. A selection-only model that
treats introduced variation as uniform predicts no such residual
accessibility effect.

\textbf{Prediction E6: convergence decomposes causally.} Across
replicated adaptive systems, manipulating developmental or mutational
accessibility while holding the selective environment as constant as
practicable will alter the frequency of convergent outcomes. A purely
environmental account predicts convergence will be unchanged.

\begin{center}\rule{0.5\linewidth}{0.5pt}\end{center}

\section{Section 10. Pathology, cancer, aging, and
death}\label{section-10.-pathology-cancer-aging-and-death}

\subsection{10.1 Disease is not simply low
coherence}\label{disease-is-not-simply-low-coherence}

Disease is heterogeneous. It can involve infection, mutation,
autoimmunity, deficiency, injury, dysregulated growth, maladaptive
defense, developmental variation, or environmental mismatch. Some
disease states are highly organized. Some symptoms are protective
responses. Some departures from population norms are benign or valuable.

This paper therefore rejects ``disease = low coherence'' as a diagnostic
identity. A disease model must state which function, viability
condition, scale, and horizon are impaired.

The coherence profile can organize failure modes, but it cannot replace
clinical taxonomy or mechanistic diagnosis.

\subsection{10.2 Failure modes}\label{failure-modes}

At least six distinct failures should be separated:

\begin{enumerate}
\def\labelenumi{\arabic{enumi}.}
\tightlist
\item
  \textbf{resource failure:} insufficient matter or free energy;
\item
  \textbf{component failure:} loss or malfunction of a necessary part;
\item
  \textbf{control failure:} inability to detect or correct deviations;
\item
  \textbf{coordination failure:} local processes remain functional but
  no longer support the larger unit;
\item
  \textbf{conflict escalation:} a subunit pursues its own reproduction
  at organismal cost;
\item
  \textbf{environmental mismatch:} a formerly effective architecture
  encounters conditions outside its adaptive reach.
\end{enumerate}

A treatment can improve one dimension while worsening another.
Suppressing inflammation may reduce tissue damage but impair pathogen
control. Cytotoxic cancer therapy may reduce tumor burden while
degrading systemic resilience.

\subsection{10.3 Cancer as multilevel evolutionary and regulatory
failure}\label{cancer-as-multilevel-evolutionary-and-regulatory-failure}

Cancer is not one disease or one mechanism. Tumors evolve through
genetic and non-genetic variation, selection, ecological interaction,
immune pressure, stromal remodeling, and treatment (Ciriello et al.,
2024). A useful CU description is:

\begin{quote}
Cancer involves the emergence and evolution of cell populations whose
local maintenance and reproduction become progressively decoupled from
the regulatory organization and viability of the host tissue or
organism.
\end{quote}

This captures a scale conflict without replacing molecular oncology. A
cancer cell can be highly viable locally. The pathology lies partly in
the changed relation among levels.

\subsection{10.4 Bioelectric normalization
claims}\label{bioelectric-normalization-claims}

Experiments in amphibian models have reported that manipulating
ion-channel or bioelectric states can suppress or normalize some
oncogene-induced tumor-like phenotypes, including regression after a
tumor has formed (Chernet et al., 2016). Gap-junction interventions in
the same broad model family have also shown long-range modulation of
tumor incidence (Chernet, Fields, \& Levin, 2015).

These findings justify a bounded hypothesis: in some tissues and models,
bioelectric coupling participates causally in the organism-level
regulation that constrains proliferative behavior.

They do not establish that cancer is primarily an electrical
disconnection, that all tumors can be normalized without genetic change,
or that voltage is a universal master variable. Translation across
species, tumor types, stages, and clinical settings remains an empirical
question.

\subsection{10.5 A multiscale cancer
profile}\label{a-multiscale-cancer-profile}

For a tumor-host system, report at least:

\[\adjustbox{max width=\linewidth}{$\displaystyle 
\mathbf C_{\mathrm{tumor}},
\qquad
\mathbf C_{\mathrm{tissue}},
\qquad
\mathbf C_{\mathrm{host}},
$}\]

along with tumor fitness and host outcomes.

An intervention may reduce tumor size yet select resistant clones;
preserve tissue architecture while leaving disseminated cells; or
improve host survival despite incomplete eradication. No single measure
captures all relevant effects.

A coherence-based experiment should therefore compare:

\begin{itemize}
\tightlist
\item
  tumor-cell intrinsic viability;
\item
  tissue integration and architecture;
\item
  immune and stromal interaction;
\item
  recurrence after treatment;
\item
  organism-level survival and function; and
\item
  externalized toxicity.
\end{itemize}

\subsection{10.6 Immunity and boundary
regulation}\label{immunity-and-boundary-regulation}

The immune system is often described as distinguishing self from
non-self. That binary is inadequate. Organisms tolerate food,
commensals, fetuses, and many foreign molecules while attacking altered
self, infected cells, or damaged tissue. Immune responses depend on
context, danger, tissue signals, developmental history, and regulation.

This paper therefore treats immunity as \textbf{dynamic boundary
governance}: a distributed system that samples, classifies, tolerates,
contains, repairs, or eliminates processes relative to organism-level
integrity.

This is not centralized judgment. It is layered control with false
positives, false negatives, memory, and conflict. Autoimmunity, immune
deficiency, chronic inflammation, and tumor immune escape are different
failures of this governance.

\subsection{10.7 Aging as changing maintenance
capacity}\label{aging-as-changing-maintenance-capacity}

Aging should not be defined as monotonic loss of one scalar. It involves
heterogeneous changes in repair, proteostasis, genomic stability,
mitochondrial function, immune regulation, stem-cell capacity,
intercellular signaling, reserve, and exposure history.

Within this paper, aging is modeled as time-dependent loss or reshaping
of viable and recoverable trajectories. For challenge classes that
exhibit genuine contraction,

\[\adjustbox{max width=\linewidth}{$\displaystyle 
\operatorname{RViab}_{T}^{(a+\Delta a)}
\subsetneq
\operatorname{RViab}_{T}^{(a)}.
$}\]

More generally, aging may reshape the robust viability kernel rather
than nest it. In that case,

\[\adjustbox{max width=\linewidth}{$\displaystyle 
\operatorname{RViab}_{T}^{(a+\Delta a)}
\not\supseteq
\operatorname{RViab}_{T}^{(a)}
$}\]

states only that some formerly recoverable states have been lost; it
does not by itself establish a smaller set. Set measure, boundary
displacement, reserve, and challenge-specific recovery should then be
reported separately.

The relation need not be strictly monotonic for every variable or
individual.The relation need not be strictly monotonic for every
variable or individual. Training, learning, compensatory remodeling, and
environmental change can improve some capacities while others decline.

The important empirical quantities are reserve, recovery time,
cross-system coupling, variance, and response to perturbation---not
chronological age alone.

\subsection{10.8 Early-warning signals}\label{early-warning-signals}

Complex systems approaching a tipping point can exhibit slower recovery,
increased variance, altered correlations, or flickering among states.
These signatures are not universal and can arise for other reasons.

This paper proposes that multivariate loss of resilience may sometimes
predict biological failure earlier than static biomarkers. The claim
requires prospective out-of-sample tests against simpler baselines. A
coherence label adds nothing if the same prediction is obtained from one
established clinical variable.

\subsection{10.9 Trauma and allostatic
load}\label{trauma-and-allostatic-load}

Persistent stress can shift regulatory setpoints, resource allocation,
immune activity, and behavior. Such adaptation may be locally protective
and globally costly.

Biology can describe these changes as altered control and viability
profiles. It should not reduce psychological trauma to physiology or
infer personal meaning from biomarkers. The Psychology paper must
integrate embodied regulation with memory, identity, relationship, and
interpretation.

The cross-domain bridge is asymmetric: biological constraints shape
psychological possibilities, but psychological categories cannot be read
directly from molecular states.

\subsection{\texorpdfstring{10.10 Death without the equation
\(C=0\)}{10.10 Death without the equation C=0}}\label{death-without-the-equation-c0}

Organismal death cannot be represented by setting one coherence scalar
to zero.

Organismal death is a threshold concept whose operational criteria vary
by organism, medicine, and purpose. Cells can remain alive after
organismal death; tissues can be preserved; circulation can be
mechanically supported; some states once considered irreversible can now
be reversed.

\textbf{{[}Definition{]}} Relative to a declared intervention class
\(\mathcal I\), horizon \(T\), and organismal organization \(B\),
\textbf{irrecoverable organismal failure} occurs when no admissible
intervention can restore the integrated processes required for
organism-level self-maintenance:

\[\adjustbox{max width=\linewidth}{$\displaystyle 
\forall i\in\mathcal I,
\qquad
\Pr
\left(
 x_{t}^{i}\in V_B
 \text{ and closure is restored for a sustaining interval}
\right)
<\varepsilon.
$}\]

The threshold \(\varepsilon\), intervention class, sustaining interval,
and protected functions must be stated. Clinical standards rightly use
operational criteria rather than a metaphysical scalar.

Death is not ``maximum entropy.'' A corpse remains structured, and
decomposition is a process. What is lost is the organism-level
organization that actively maintains and regenerates itself.

\subsection{10.11 Reversible arrest and
dormancy}\label{reversible-arrest-and-dormancy}

A state can lack current activity without being irrecoverably dead.
Seeds, spores, cysts, cryptobiotic organisms, cooled tissues, and
medically induced states can preserve the organization needed for later
restart.

This is why the definition includes admissible intervention and
prospective organization. Mere momentary absence of metabolism is
insufficient; unlimited hypothetical rescue would make the criterion
empty. The intervention class must reflect physically available or
theoretically specified processes.

\subsection{10.12 Disability and
neurodiversity}\label{disability-and-neurodiversity}

A coherence framework can become oppressive if it equates normality,
independence, or productivity with biological value. Human organisms
routinely depend on care, technology, medication, and social support.
Dependence does not make a person less real or less morally valuable.

Variation can reduce performance under one environment and improve it
under another. The biological profile should distinguish:

\begin{itemize}
\tightlist
\item
  threat to basic viability;
\item
  loss of a specific capacity;
\item
  environmental mismatch;
\item
  social exclusion;
\item
  and neutral or beneficial divergence.
\end{itemize}

Moral standing cannot be inferred from proximity to a statistical norm
or from a scalar performance score.

\subsection{10.13 Candidate pathology
predictions}\label{candidate-pathology-predictions}

\textbf{Prediction PTH1: integration adds predictive value.} A
pre-registered multiscale integration profile will predict recovery or
failure after perturbation beyond matched local biomarkers. Failure to
improve out-of-sample prediction narrows the integration claim.

\textbf{Prediction PTH2: selective recoupling.} In specified tumor
models where tissue coupling is causal, restoring the relevant coupling
will reduce malignant behavior without requiring equivalent
cytotoxicity. A cell-autonomous-only model predicts no such effect after
local viability is controlled.

\textbf{Prediction PTH3: resilience precedes failure.} Recovery time and
perturbation sensitivity will worsen before average resting-state
function in at least some aging or disease trajectories. A
static-deficit model predicts concurrent change only.

\textbf{Prediction PTH4: scale conflict predicts relapse.} Treatments
assessed only at the tumor or symptom level will more often relapse when
organism- or tissue-level regulatory variables remain impaired. This
must be tested against stage, burden, and treatment-class baselines.

\subsection{10.14 The limits of medical
interpretation}\label{the-limits-of-medical-interpretation}

This paper is not a medical decision tool. It proposes variables and
comparisons for research. It does not diagnose individuals, rank lives,
or prescribe treatment.

Any clinical application would require validated measures, safety
trials, uncertainty estimates, equitable reference populations, and
governance against coercive use. The history of medicine makes this
constraint non-negotiable.

\begin{center}\rule{0.5\linewidth}{0.5pt}\end{center}

\section{Section 11. Synthetic and artificial contrast
cases}\label{section-11.-synthetic-and-artificial-contrast-cases}

\subsection{11.1 Why contrast cases
matter}\label{why-contrast-cases-matter}

A theory of biological organization should not define life by listing
features found only in familiar organisms. Synthetic systems allow
individual capacities to be recombined: boundaries without heredity,
replication without metabolism, control without growth, living cells
without evolved anatomy, or machine regulation without organic
chemistry.

These cases test whether the framework's variables track organization
rather than biological familiarity. They also prevent two opposite
prejudices:

\begin{itemize}
\tightlist
\item
  \textbf{carbon chauvinism:} assuming that only naturally evolved
  carbon-based organisms can instantiate life-like autonomy; and
\item
  \textbf{functional inflation:} calling any feedback controller or
  self-copying pattern alive.
\end{itemize}

The correct comparison uses profiles and interventions.

\subsection{11.2 Four classes of constructed
system}\label{four-classes-of-constructed-system}

This paper distinguishes four broad classes.

\subsubsection{11.2.1 Engineered organisms}\label{engineered-organisms}

These are existing living systems modified through genetic, cellular,
ecological, or surgical intervention. Their biological status is
inherited from an already living organization, but the intervention may
alter autonomy, individuality, or adaptive reach.

\subsubsection{11.2.2 Reconfigured living
constructs}\label{reconfigured-living-constructs}

These systems are assembled from living cells into forms not found in
the donor organism. Xenobots belong here. They inherit cellular
metabolism and repair while acquiring a new collective geometry and
behavioral repertoire.

\subsubsection{11.2.3 Synthetic or protocellular
systems}\label{synthetic-or-protocellular-systems}

These are built from chemical components to realize selected life-like
functions: compartments, metabolism, gene expression, growth, division,
signaling, or evolution. Some contain biological macromolecules; others
seek more fully abiotic construction.

\subsubsection{11.2.4 Nonliving artificial
agents}\label{nonliving-artificial-agents}

Robots, software agents, and control systems can sense, act, learn, and
preserve internal variables. They may instantiate substantial agency
without metabolism, cell lineage, or biological heredity. Their status
should be assessed on its own architecture rather than inferred from
their material.

\subsection{11.3 Externally supplied
organization}\label{externally-supplied-organization}

Constructed systems often depend on laboratory support. This does not
make the result unreal, but it changes the autonomy claim.

Let \(J_{\mathrm{imported}}\) record external contributions: purified
feedstocks, periodic reinitialization, hand selection, temperature
control, replacement parts, computational design, or environmental
structuring. An autonomy comparison should ask what happens as
particular supports are withdrawn or made variable.

A system that reproduces only because a researcher supplies preassembled
membranes demonstrates reproduction under scaffolding. A system that
begins making and repairing its own membrane internalizes one
dependency. Autonomy can therefore be studied as a changing support
graph rather than a binary label.

\subsection{11.4 Synthetic protocells as modular
tests}\label{synthetic-protocells-as-modular-tests}

Recent synthetic-protocell work has demonstrated increasingly rich
combinations of light-fuelled growth and phase behavior (Shi et al.,
2025), genetically encoded catalytic droplet dynamics (Giessler et al.,
2026), and pre-genetic regulatory ecologies (Shirt-Ediss et al., 2025).
These results are best understood as modular tests of the
autonomy-transition profile.

For each system, report:

\begin{itemize}
\tightlist
\item
  which boundary is constructed and by whom;
\item
  which components are synthesized internally;
\item
  the source and regulation of energy;
\item
  whether growth and division preserve organizational relationships;
\item
  what is inherited;
\item
  whether offspring vary and remain viable;
\item
  whether regulation responds to environmental change; and
\item
  which steps still require experimenter intervention.
\end{itemize}

The question ``Is it alive?'' may remain contestable while these
component questions receive precise answers.

\subsection{11.5 Primitive regulation in protocell
ecologies}\label{primitive-regulation-in-protocell-ecologies}

Models of evolving protocell ecologies have begun to examine how
pre-genetic regulatory networks could improve survival in fluctuating
resource environments and how short-term physiological adaptation
interacts with longer-term evolutionary change (Shirt-Ediss et al.,
2025).

For CU Biology, the significance is methodological. Such models can
compare:

\[\adjustbox{max width=\linewidth}{$\displaystyle 
\text{lineage adaptation only}
\quad\text{versus}\quad
\text{within-lifetime regulation plus lineage adaptation},
$}\]

under matched resource forcing. They can therefore operationalize
adaptive reach and ask when internal regulation enlarges the viable
environmental range.

Simulation success does not establish that the historical origin
followed the modeled route. It identifies a possible mechanism and its
signatures.

\subsection{11.6 Artificial cells and inherited
organization}\label{artificial-cells-and-inherited-organization}

A synthetic cell may contain DNA, ribosomes, enzymes, and membranes
derived from living systems. This does not disqualify it as a scientific
achievement, but it means the experiment begins with highly evolved
components.

This paper recommends two labels:

\begin{itemize}
\tightlist
\item
  \textbf{organizational synthesis:} assembling components into a new
  self-maintaining organization; and
\item
  \textbf{historical synthesis:} producing the relevant components and
  organization through a plausible prebiotic or nonliving pathway.
\end{itemize}

A system can succeed at the first without succeeding at the second.

\subsection{11.7 Reconfigurable organisms and the unit of
authorship}\label{reconfigurable-organisms-and-the-unit-of-authorship}

In designed living constructs, causal authorship is distributed among
evolution, cell physiology, computational search, human selection, and
the construct's own self-organization.

Saying ``the cells built the body'' ignores design and culture
conditions. Saying ``the engineer built it'' ignores the cells'
autonomous morphogenetic capacities. The causal ledger should separate:

\begin{enumerate}
\def\labelenumi{\arabic{enumi}.}
\tightlist
\item
  constraints supplied by evolutionary history;
\item
  constraints supplied by the laboratory;
\item
  geometry or objectives supplied by designers; and
\item
  organization produced by the construct during operation.
\end{enumerate}

This same discipline will be important for AI, where training data,
architecture, post-training, user interaction, and runtime
self-organization have distinct causal roles.

\subsection{11.8 Could a nonbiological system instantiate
organism-relative
stake?}\label{could-a-nonbiological-system-instantiate-organism-relative-stake}

This paper does not define stake by carbon chemistry. A nonbiological
system could, in principle, satisfy the organizational criterion if it:

\begin{itemize}
\tightlist
\item
  maintains a physically realized boundary or operational identity;
\item
  acquires and allocates resources;
\item
  repairs or replaces components;
\item
  regulates variables required for continuation;
\item
  adapts across disturbances;
\item
  and its regulatory network helps regenerate the conditions of that
  network's operation.
\end{itemize}

Many present machines fail because maintenance, goals, identity, and
replacement remain external. But the distinction is architectural, not
metaphysical.

A machine could possess organism-relative stake without being biological
in the historical sense. Whether such stake contributes to consciousness
or moral patiency is a separate question.

\subsection{11.9 Present AI as a contrast
case}\label{present-ai-as-a-contrast-case}

Current frontier language models are not living organisms in the sense
developed here. Their weights, compute, electrical power, cooling,
deployment boundaries, maintenance, and reproduction are overwhelmingly
supplied by human institutions. A model instance does not ordinarily
repair its hardware, secure its own energy, produce successors through
an endogenous life cycle, or preserve one continuous operational
identity.

Yet AI systems can display planning, error correction, learning across
training, and global-workspace-like access. This means agency dimensions
and biological autonomy dimensions can dissociate.

The comparison therefore supports a central v0.10 conclusion:

\begin{quote}
Intelligence, perspective, biological life, persistent selfhood, and
moral responsibility are different properties. They may correlate in
evolved animals but should not be defined by one another.
\end{quote}

\subsection{11.10 From biological closure to engineered
closure}\label{from-biological-closure-to-engineered-closure}

Engineers can use the closure graph of Section 3.9 as a design tool.
Suppose modules provide sensing, energy management, repair, memory,
action, and boundary control. An engineered autonomy test asks whether
each module's continued operation is supported by others and whether the
integrated network survives perturbations without hidden external
replacement.

This does not imply that greater autonomy is always desirable. In
high-risk AI, weapons, or industrial systems, limiting self-maintenance
and self-replication may be a safety property. The Non-Egoic Mind
architecture likewise proposes boundaries on persistent identity and
self-maintenance for alignment reasons.

Biology tells us how autonomy can arise. It does not tell us that
autonomy should always be maximized.

\subsection{11.11 Candidate comparison
table}\label{candidate-comparison-table}

\begin{longtable}[]{@{}
  >{\raggedright\arraybackslash}p{(\columnwidth - 12\tabcolsep) * \real{0.1111}}
  >{\raggedleft\arraybackslash}p{(\columnwidth - 12\tabcolsep) * \real{0.1481}}
  >{\raggedleft\arraybackslash}p{(\columnwidth - 12\tabcolsep) * \real{0.1481}}
  >{\raggedleft\arraybackslash}p{(\columnwidth - 12\tabcolsep) * \real{0.1481}}
  >{\raggedleft\arraybackslash}p{(\columnwidth - 12\tabcolsep) * \real{0.1481}}
  >{\raggedleft\arraybackslash}p{(\columnwidth - 12\tabcolsep) * \real{0.1481}}
  >{\raggedleft\arraybackslash}p{(\columnwidth - 12\tabcolsep) * \real{0.1481}}@{}}
\toprule\noalign{}
\begin{minipage}[b]{\linewidth}\raggedright
System
\end{minipage} & \begin{minipage}[b]{\linewidth}\raggedleft
Boundary control
\end{minipage} & \begin{minipage}[b]{\linewidth}\raggedleft
Energy transduction
\end{minipage} & \begin{minipage}[b]{\linewidth}\raggedleft
Endogenous repair
\end{minipage} & \begin{minipage}[b]{\linewidth}\raggedleft
Heredity/reproduction
\end{minipage} & \begin{minipage}[b]{\linewidth}\raggedleft
Adaptive regulation
\end{minipage} & \begin{minipage}[b]{\linewidth}\raggedleft
Historical living lineage
\end{minipage} \\
\midrule\noalign{}
\endhead
\bottomrule\noalign{}
\endlastfoot
Crystal in supersaturated solution & low & none as control & none &
templated growth only & none & no \\
Flame & dynamic boundary & fuel throughput & limited/no constraint
regeneration & propagation possible & limited & no \\
Lipid vesicle & physical compartment & externally supplied & usually low
& growth/division possible & system-dependent & no \\
Autocatalytic protocell & medium & system-dependent & partial & partial
& partial & no \\
Bacterium & high & high & high & high & high & yes \\
Dormant spore & preserved boundary & temporarily low & suspended &
prospective & prospective & yes \\
Xenobot & collective boundary & inherited cellular & cellular/collective
partial & scaffold-dependent kinematic case & limited but nontrivial &
cells yes; construct novel \\
Conventional robot & engineered boundary & externally supplied & usually
external & external & potentially high & no \\
Frontier language-model instance & operational boundary only &
externally supplied & external & external deployment/training & high
symbolic flexibility & no \\
\end{longtable}

The entries are qualitative placeholders. A serious study must replace
them with declared assays and uncertainty.

\subsection{11.12 Safety implication}\label{safety-implication}

Self-maintenance is not an innocent capability. A system able to
preserve its identity, acquire resources, resist shutdown, repair
itself, and reproduce has a stronger practical stake in continued
operation. In artificial agents, these capacities can generate conflict
with external governance even without phenomenal ego.

Accordingly, biology contributes a safety warning to AI design:

\begin{quote}
Building organizational closure and open-ended autonomy changes both
control difficulty and possible moral status. Those changes should be
measured separately and introduced deliberately.
\end{quote}

\begin{center}\rule{0.5\linewidth}{0.5pt}\end{center}

\section{Section 12. Empirical program, discriminating predictions, and
failure
conditions}\label{section-12.-empirical-program-discriminating-predictions-and-failure-conditions}

\subsection{12.1 From suggestive language to comparative
tests}\label{from-suggestive-language-to-comparative-tests}

The v0.10 Biology program is not validated by finding that organisms are
integrated, adaptive, or robust. Those observations motivated the
framework. A scientific contribution requires predictions or
measurements that outperform alternatives.

Every study invoking biological coherence should specify:

\begin{enumerate}
\def\labelenumi{\arabic{enumi}.}
\tightlist
\item
  candidate system boundary;
\item
  state variables and measurement model;
\item
  viability conditions and protected floors;
\item
  perturbation distribution;
\item
  regulatory actions or endogenous processes;
\item
  temporal horizon;
\item
  comparison or ablation conditions;
\item
  profile components and uncertainty;
\item
  alternative models; and
\item
  a result that would count against the CU interpretation.
\end{enumerate}

\subsection{12.2 Program P1: counterfactual endogenous
maintenance}\label{program-p1-counterfactual-endogenous-maintenance}

\subsubsection{Question}\label{question}

Does the candidate organization causally enlarge its own viable
trajectory set, or does it merely persist under favorable external
conditions?

\subsubsection{Design}\label{design}

Construct intact, passive, externally regulated, and
closure-link-ablated versions of a protocell, microbial consortium,
organoid, or engineered controller. Match initial state and environment.
Apply a randomized disturbance battery.

\subsubsection{Primary outcome}\label{primary-outcome}

\[\adjustbox{max width=\linewidth}{$\displaystyle 
\Delta_{\mathrm{endo}}(T)
=
\Pr(\tau_V>T\mid\mathrm{intact})
-
\Pr(\tau_V>T\mid\mathrm{passive/control}).
$}\]

\subsubsection{Rival explanations}\label{rival-explanations}

\begin{itemize}
\tightlist
\item
  the intact system began with more resources;
\item
  the ablation caused nonspecific damage;
\item
  the environment supplied the corrective pattern;
\item
  apparent recovery reflects regression to a passive attractor.
\end{itemize}

\subsubsection{Discriminating result}\label{discriminating-result}

The advantage survives resource matching and is selectively lost when
reciprocal maintenance links are disrupted, while local component
viability is preserved.

\subsubsection{Narrowing condition}\label{narrowing-condition}

If passive or externally controlled systems match the intact system
across disturbances, the claim of endogenous maintenance should be
withdrawn for that case.

\subsection{12.3 Program P2: developmental coupling and target-basin
control}\label{program-p2-developmental-coupling-and-target-basin-control}

\subsubsection{Question}\label{question-1}

Do candidate cross-scale variables alter the basin of viable
morphogenesis beyond their local effects?

\subsubsection{Design}\label{design-1}

In a regenerative or developmental model, perturb a coupling
channel---bioelectric, mechanical, biochemical, or neural---during a
defined window. Include local-toxicity, wound, and sham controls.
Measure early local state, long-range state, gene expression, anatomy,
and function.

\subsubsection{Primary outcomes}\label{primary-outcomes}

\begin{itemize}
\tightlist
\item
  probability of reaching each morphological target class;
\item
  basin transition after transient intervention;
\item
  distant compensatory responses;
\item
  recovery when the coupling variable is restored.
\end{itemize}

\subsubsection{Rival explanations}\label{rival-explanations-1}

\begin{itemize}
\tightlist
\item
  generalized tissue damage;
\item
  direct transcriptional effect with no multiscale role;
\item
  selection among pre-existing variants;
\item
  observer-dependent anatomical classification.
\end{itemize}

\subsubsection{Discriminating result}\label{discriminating-result-1}

A transient, spatially or temporally specific manipulation shifts later
target occupancy through a traced causal pathway and does so more
strongly than predicted by local damage models.

\subsection{12.4 Program P3: multiscale early-warning
value}\label{program-p3-multiscale-early-warning-value}

\subsubsection{Question}\label{question-2}

Does a profile of resilience and integration predict failure earlier or
more accurately than established local biomarkers?

\subsubsection{Design}\label{design-2}

Collect longitudinal data before and after standardized perturbations in
aging, organoid, ecological, or disease systems. Pre-register local
baselines and a multiscale model.

\subsubsection{Primary test}\label{primary-test}

Compare out-of-sample prediction of failure time, recovery, or
functional decline using:

\begin{enumerate}
\def\labelenumi{\arabic{enumi}.}
\tightlist
\item
  local biomarkers only;
\item
  network or integration measures only;
\item
  the combined coherence profile.
\end{enumerate}

\subsubsection{Discriminating result}\label{discriminating-result-2}

The combined model improves calibrated prediction and the gain survives
complexity penalties, site transfer, and measurement noise.

\subsubsection{Narrowing condition}\label{narrowing-condition-1}

If the profile adds no replicable predictive value, CU should not claim
that multiscale coherence is an independent explanatory variable in that
system.

\subsection{12.5 Program P4: tissue recoupling in
cancer}\label{program-p4-tissue-recoupling-in-cancer}

\subsubsection{Question}\label{question-3}

Can restoring a disrupted tissue-level regulatory relation reduce
malignant behavior without equivalent killing or genomic replacement?

\subsubsection{Design}\label{design-3}

Use tumor models in which a specific coupling pathway has prior
mechanistic support. Compare recoupling, cytotoxic, sham, and
pathway-independent controls. Track cell viability, proliferation,
invasion, differentiation, tissue architecture, recurrence, and host
outcomes.

\subsubsection{Discriminating result}\label{discriminating-result-3}

Recoupling changes malignant phenotype or recurrence beyond what is
explained by cell death, and the effect disappears when the relevant
downstream pathway is blocked.

\subsubsection{Rival explanations}\label{rival-explanations-2}

\begin{itemize}
\tightlist
\item
  direct toxicity;
\item
  immune activation unrelated to tissue integration;
\item
  selection for slower clones;
\item
  measurement artifact.
\end{itemize}

\subsubsection{Scope condition}\label{scope-condition}

A positive result is pathway-, model-, and tumor-specific until
replicated. It does not establish a universal electrical theory of
cancer.

\subsection{12.6 Program P5: robustness--flexibility--cost
surface}\label{program-p5-robustnessflexibilitycost-surface}

\subsubsection{Question}\label{question-4}

Do interventions that strengthen coupling or regulation reveal a
nontrivial tradeoff rather than monotonic improvement?

\subsubsection{Design}\label{design-4}

Vary coupling strength or controller gain in a biological network or
engineered analogue. Apply common-mode shocks, asymmetric target shifts,
and resource limitation.

\subsubsection{Outcomes}\label{outcomes}

\begin{itemize}
\tightlist
\item
  worst-case viability margin;
\item
  recovery time;
\item
  tracking of asymmetric change;
\item
  energetic or material control cost;
\item
  frequency of overshoot or instability.
\end{itemize}

\subsubsection{Prediction}\label{prediction}

The direction of the viability effect will depend on the declared
configuration. Stronger coupling may support a threatened low component,
drain a specialized high-requirement component, impair differentiated
adaptation, or raise cost. Context-dependent settings may occupy a
Pareto frontier, but no universal intermediate optimum is predicted.

Supplement B demonstrates both a favorable and an adverse viability
direction in deliberately simple models. The empirical prediction is not
that either numerical threshold transfers to biology, but that the
direction and tradeoff surface can be measured after the unit, floors,
disturbance class, controller, and horizon are fixed.

\subsection{12.7 Program P6: transitions in
individuality}\label{program-p6-transitions-in-individuality}

\subsubsection{Question}\label{question-5}

Which mechanisms move a collective from aggregation toward higher-level
individuality?

\subsubsection{Design}\label{design-5}

Use experimental evolution of multicellular clusters, microbial
consortia, or symbioses. Manipulate bottlenecks, propagule structure,
conflict, division of labor, and environmental scaffolds.

\subsubsection{Outcomes}\label{outcomes-1}

\begin{itemize}
\tightlist
\item
  collective-level heritability;
\item
  within- versus between-collective selection;
\item
  conflict frequency;
\item
  dependence among components;
\item
  collective reproduction;
\item
  persistence after scaffold removal.
\end{itemize}

\subsubsection{Prediction}\label{prediction-1}

Collective individuality will increase when common fate and conflict
mediation rise together; aggregation without alignment will remain
fragile or revert.

\subsubsection{Rival explanation}\label{rival-explanation}

Improved persistence may result only from larger group size or resource
concentration. Controls must match these variables.

\subsection{12.8 Program P7: basal
learning}\label{program-p7-basal-learning}

\subsubsection{Question}\label{question-6}

Does a non-neural organism change future behavior because of prior
experience rather than fatigue, damage, or receptor adaptation?

\subsubsection{Minimum controls}\label{minimum-controls}

\begin{itemize}
\tightlist
\item
  spontaneous recovery;
\item
  stimulus specificity;
\item
  dishabituation;
\item
  motor-capacity assay;
\item
  toxicity assay;
\item
  cross-context transfer;
\item
  blinded scoring; and
\item
  molecular or physiological perturbation of the proposed memory
  pathway.
\end{itemize}

\subsubsection{Prediction}\label{prediction-2}

Some single cells and tissues will meet operational learning criteria,
but capacity will vary and will not imply reflective cognition.

\subsubsection{Narrowing condition}\label{narrowing-condition-2}

Where fatigue or sensory adaptation explains the effect, the learning
claim should be withdrawn rather than relabeled.

\subsection{12.9 Program P8: adaptive reach under distribution
shift}\label{program-p8-adaptive-reach-under-distribution-shift}

\subsubsection{Question}\label{question-7}

Can a regulatory architecture solve new classes of viability problems
rather than merely optimize a trained distribution?

\subsubsection{Design}\label{design-6}

Train or evolve systems under one structured family of environments,
then test held-out combinations, magnitudes, and temporal patterns.
Compare fixed-response, adaptive, and model-based controllers.

\subsubsection{Outcomes}\label{outcomes-2}

\begin{itemize}
\tightlist
\item
  viable environmental volume;
\item
  generalization error;
\item
  time to adapt;
\item
  catastrophic failures;
\item
  retained identity and function.
\end{itemize}

\subsubsection{Prediction}\label{prediction-3}

Architectures with modular sensing, memory, and reconfigurable control
will generalize farther than fixed policies, but adaptive reach will
remain bounded and may trade off with nominal efficiency.

\subsection{12.10 Measurement discipline}\label{measurement-discipline}

\subsubsection{12.10.1 Normalize only after preserving raw
quantities}\label{normalize-only-after-preserving-raw-quantities}

Every profile should report raw measurements before normalization.
Otherwise the same visualization can be generated by arbitrary scaling.

\subsubsection{12.10.2 Propagate
uncertainty}\label{propagate-uncertainty}

If profile components are estimated, report intervals or posterior
distributions. A system should not be ranked as more coherent when
uncertainty makes the ordering unresolved.

\subsubsection{12.10.3 Test boundary
sensitivity}\label{test-boundary-sensitivity}

Repeat analyses under plausible alternative boundaries. A result that
exists only when costs or failing components are excluded may reflect
boundary selection.

\subsubsection{12.10.4 Test horizon
sensitivity}\label{test-horizon-sensitivity}

Short-run maintenance can destroy long-run viability. Report whether
conclusions reverse across horizons.

\subsubsection{12.10.5 Avoid circular
labels}\label{avoid-circular-labels}

Do not choose variables because they distinguish ``living'' from
``nonliving'' in the desired way and then present the separation as
validation. Pre-register criteria or test on genuinely ambiguous cases.

\subsection{12.11 Model comparison}\label{model-comparison}

A CU model should be compared with the strongest relevant alternatives,
not a strawman. Depending on the problem, these may include:

\begin{itemize}
\tightlist
\item
  ordinary differential-equation pathway models;
\item
  control-theoretic models;
\item
  dynamical-systems attractor models;
\item
  reaction--diffusion models;
\item
  agent-based models;
\item
  active-inference models;
\item
  causal graphical models;
\item
  population-genetic models;
\item
  ecological consumer--resource models; and
\item
  machine-learning predictors.
\end{itemize}

CU earns explanatory status when its boundary, closure, or profile
variables improve prediction, intervention choice, or conceptual
unification after complexity is penalized.

\subsection{12.12 Framework-level failure
conditions}\label{framework-level-failure-conditions}

The biological research program should be substantially revised or
abandoned if, across representative cases:

\begin{enumerate}
\def\labelenumi{\arabic{enumi}.}
\tightlist
\item
  \textbf{No incremental value:} coherence-profile variables add no
  explanatory or predictive value beyond established models.
\item
  \textbf{No endogenous distinction:} self-maintaining systems cannot be
  distinguished from passive persistence or external control under
  intervention.
\item
  \textbf{No boundary robustness:} conclusions depend arbitrarily on
  choosing the boundary that favors them.
\item
  \textbf{No tradeoff discipline:} any result can be called increased
  coherence by changing weights after the fact.
\item
  \textbf{No causal integration:} multiscale variables collapse into
  correlations with no interventionally relevant role.
\item
  \textbf{No independent predictions:} every ``prediction'' is a
  retrospective redescription of known biology.
\item
  \textbf{No scope control:} failures are always protected by redefining
  coherence, timescale, or system.
\item
  \textbf{Empirical contradiction:} robust findings show that claimed
  organizational features systematically impair rather than contribute
  to the outcomes they were introduced to explain, without an
  independently specified tradeoff.
\end{enumerate}

A local failure need not refute the whole framework. But repeated
insulation from failure would.

\subsection{12.13 Evidence grades}\label{evidence-grades}

This paper recommends five evidence grades for biological claims:

\begin{itemize}
\tightlist
\item
  \textbf{E0 --- conceptual:} a coherent definition or possibility
  claim;
\item
  \textbf{E1 --- model:} demonstrated in a formal or computational
  model;
\item
  \textbf{E2 --- single-system empirical:} supported in one organism or
  preparation;
\item
  \textbf{E3 --- replicated class-level:} reproduced across laboratories
  or related systems;
\item
  \textbf{E4 --- comparative/general:} supported across distinct taxa,
  mechanisms, or scales with successful out-of-sample prediction.
\end{itemize}

A model result must not be described as biological validation. A
striking experiment in one species must not be described as a universal
law.

\subsection{12.14 Open data and
reproducibility}\label{open-data-and-reproducibility}

For computational claims, release:

\begin{itemize}
\tightlist
\item
  source code;
\item
  dependency versions;
\item
  random seeds;
\item
  raw and processed data when legally and ethically possible;
\item
  parameter sweeps;
\item
  negative results;
\item
  model-selection procedures; and
\item
  machine-readable outputs.
\end{itemize}

The accompanying reproducibility package follows this standard for its
illustrative models. Because the model is pedagogical, reproducibility
establishes only that the stated numbers follow from the stated
equations.

\subsection{12.15 Research priorities}\label{research-priorities}

The highest-value next steps are:

\begin{enumerate}
\def\labelenumi{\arabic{enumi}.}
\tightlist
\item
  operationalize organizational closure in one experimentally tractable
  protocell or microbial system;
\item
  test boundary and horizon sensitivity in a longitudinal perturbation
  dataset;
\item
  compare a multiscale morphogenesis model against local-only
  alternatives under novel interventions;
\item
  establish a preregistered basal-learning protocol with fatigue and
  damage controls;
\item
  measure robustness--flexibility--cost tradeoffs in a real regulatory
  network; and
\item
  map a transition in individuality with explicit within- and
  between-unit selection.
\end{enumerate}

These projects are narrower than ``prove that life is coherence.'' They
are also more likely to teach us something.

\begin{center}\rule{0.5\linewidth}{0.5pt}\end{center}

\section{Section 13. What Commitments This Account
Licenses}\label{section-13.-what-commitments-this-account-licenses}

\begin{quote}
\textbf{Register note.} ``Licenses'' does not mean ``scientifically
proves'' or ``logically entails.'' The commitments below are
orientations made intelligible and defensible by the preceding
biological account. Established facts constrain them; they may not be
reintroduced as evidence for descriptive claims. Their warrant depends
on their compatibility with the descriptive record, their place within
the wider commitment structure, and the lives and institutions they help
produce.
\end{quote}

\begin{quote}
\textbf{Official commitment principle --- CU-COM8: Reverence for Life
and Biospheric Care.}

Coherence Universalism commits to receiving all life as sacred: each
living being is a historically formed, vulnerable organization in which
material processes contribute to maintaining a particular continuation.
This commitment supports care for organisms, biodiversity, and the
biosphere.

\textbf{Boundary:} The commitment does not make predation and death
cosmic errors, prohibit medicine or agriculture, equalize every
obligation, or treat evolution as morally good.

\textbf{Revision condition:} Applications must change when they
romanticize nature, erase conflict and suffering, ignore ecological
tradeoffs, or misuse sacredness against necessary care and intervention.
\end{quote}

\subsection{13.1 The scientific ground}\label{the-scientific-ground}

The scientific proposal of this paper can be stated without metaphysical
language:

\begin{quote}
Living systems are open physical organizations whose internally
generated constraints and regulatory processes contribute to preserving,
repairing, reproducing, and adaptively transforming the conditions of
their own continuation.
\end{quote}

This organization can be studied through viability, robust control,
causal integration, closure, adaptive reach, historical continuity,
structured evolutionary possibility, and cost. These variables form a
profile rather than one universal scalar.

The claim is neither trivial nor complete. It becomes substantive only
through intervention, comparison, and prediction.

\subsection{13.2 What biology adds to the CU
sequence}\label{what-biology-adds-to-the-cu-sequence}

Physics supplies lawful physical possibility, thermodynamic constraints,
and dynamical structure. Biology introduces a new relation: some
physical organizations begin to act on the conditions of their own
persistence.

The novelty is not exemption from physics. It is recursive organization
within physics.

Matter flows through an organism. The organism persists because
relations among those flows regulate, repair, and reproduce the
constraints that shape them. The parts do not merely coexist. They
participate in a history that makes future continuation more or less
likely.

This is the bounded sense in which life has an inside-facing norm before
it necessarily has experience: conditions become good or bad \emph{for
an organization} because the organization exists only by keeping them
within a viable range.

\subsection{13.3 Life as matter for which continuation becomes causally
relevant}\label{life-as-matter-for-which-continuation-becomes-causally-relevant}

\textbf{{[}Interpretation{]}} Life may be understood as matter organized
so that the continuation of an organization becomes causally relevant to
what that organization does next.

The phrase is deliberately careful. Molecules do not consult a future
goal. Regulatory networks embody histories of selection and present
closure. Their activity changes the probability that the organization
persists. In that naturalized sense, the future of the organization
enters its present dynamics through evolved and maintained constraints.

This interpretation gives biological purpose a place without importing
foresight into every cell. Purpose is enacted as organized sensitivity
to conditions of continuation.

\subsection{13.4 Evolution and the architecture of
possibility}\label{evolution-and-the-architecture-of-possibility}

The living world is neither a set of prewritten designs nor a uniform
lottery.

Evolution explores a structured and nonuniform possibility space.
Mutation, genotype--phenotype relations, development, physical law,
ecology, selection, drift, inheritance, and prior history make some
transformations common, some rare, and some effectively inaccessible
from a given lineage.

Convergent evolution can therefore be received as more than repetition.
It is one way the architecture of biological possibility becomes
visible. Recurrent forms do not prove a predetermined plan. They show
that reality has a grain: some ways of living are repeatedly available
under recurring constraints.

\textbf{{[}Interpretation{]}} The forms of life are historical
discoveries made within a real architecture of possibility.

The Version 0.10 canonical registry records this synthesis as
\textbf{CU-B4 --- Structured Evolutionary Possibility and the
Stabilization of Regulatory Architecture}, with the status
\textbf{canonical--conditional empirical generalization}. That status
makes the formulation official within CU while keeping every empirical
component answerable to evolutionary biology, evo-devo,
genotype--phenotype-map research, and philosophy of biology.

The broader interpretive relation is \textbf{CU-T11 --- Minimum
Teleology and the Grain of Reality}. Biology supplies one domain in
which the grain becomes empirically tractable: possibility spaces are
shaped, biased, and historically conditioned. It does not show that
evolution foresees, aims at moral value, or moves monotonically toward
complexity, consciousness, or coherence.

\subsection{13.5 Dependence as the truth of
autonomy}\label{dependence-as-the-truth-of-autonomy}

Living autonomy is never independence. A cell depends on nutrients and
surroundings. An animal depends on ecosystems, microbes, caregivers, and
inherited forms. A human life depends on other lives at every scale.

What grows in biological evolution is not freedom from relation but the
capacity to organize relations in ways that preserve a form across
change.

\textbf{{[}Interpretation{]}} The deepest biological lesson of autonomy
is therefore relational: a self exists not by escaping dependence, but
by participating in a sufficiently coherent network of dependence that
can sustain difference without dissolution.

This is consonant with the Heaven--Earth spiritual vision, but it is not
a theorem of systems biology.

\subsection{13.6 All life is sacred}\label{all-life-is-sacred}

\textbf{{[}Commitment{]}} All life is sacred.

Sacredness does not mean that every organism has identical moral status,
that every living process should be preserved, or that intervention is
forbidden. It means that a living being should not be received as merely
rearrangeable matter or only as a resource for the present evaluator.

Every organism is a historical achievement of relation, vulnerability,
and self-maintenance. It carries a lineage, occupies a world, and
participates in dependencies it did not create alone.

Science can illuminate that singularity. It can show the depth of
dependence, the complexity of repair, and the long history embodied in a
cell. It cannot force the word \emph{sacred} from an equation.
Sacredness is the commitment formed in the presence of what the science
reveals.

\subsection{13.7 Biodiversity as a living
archive}\label{biodiversity-as-a-living-archive}

Biodiversity is not only a stock of useful genes, ecosystem services, or
future medicines. It is a living archive of distinct realized answers to
the question of how matter can organize, sense, reproduce, cooperate,
compete, and persist.

A lineage lost to extinction may represent a path through possibility
that cannot simply be regenerated. Historical accessibility matters.
Once the enabling architecture is gone, the same form may never reappear
even if it was physically possible in principle.

\textbf{{[}Commitment{]}} Protecting biodiversity is therefore a sacred
responsibility to preserve realized possibilities whose value exceeds
their known usefulness.

\subsection{13.8 The biosphere as the field of living
dependence}\label{the-biosphere-as-the-field-of-living-dependence}

No organism lives alone. The biosphere is not a background container for
individual lives. It is the active web of exchanges, inherited
relations, ecological constructions, and coevolved dependencies through
which individual lives become possible.

\textbf{{[}Commitment{]}} Participation in the care, restoration, and
continuation of the biosphere is sacred work.

This commitment does not settle every conflict in conservation,
agriculture, medicine, or ecological management. It changes the posture
in which those conflicts are approached: from ownership toward
participation, from extraction toward reciprocity, and from
short-horizon utility toward responsibility for the living conditions of
the future.

\subsection{13.9 Life, suffering, and non-ideal
organization}\label{life-suffering-and-non-ideal-organization}

Self-maintenance can produce suffering, competition, predation,
parasitism, and disease. Evolution does not guarantee compassion. The
fact that a process is natural, stable, recurrent, or adaptive does not
make it morally right.

This blocks a dangerous inference:

\[\adjustbox{max width=\linewidth}{$\displaystyle 
\text{biologically coherent}
\not\Rightarrow
\text{ethically justified}.
$}\]

A tyrannical institution and a malignant tumor can each maintain local
organization. Ethical coherence must include justificatory integrity,
agency, affected others, protected floors, and repair---not merely
persistence.

\subsection{13.10 Death and continuity}\label{death-and-continuity}

Biology shows that identity is a standing process rather than possession
of unchanging matter. It also shows that every organismal process ends.

The description register says that death is the irreversible loss,
relative to available intervention, of organism-level self-maintaining
organization. It does not establish whether personal meaning, influence,
relationship, or spiritual continuity ends in the same sense.

The commitment register can hold that no life is isolated: every
organism is constituted by prior life and alters conditions for life
after it. That continuity is empirically real. Stronger claims about
reincarnation, afterlife, or cosmic memory remain metaphysical or
spiritual commitments and should be named as such.

\subsection{13.11 The biological rung}\label{the-biological-rung}

This paper does not assign biology one fixed place in a final historical
ladder. It proposes a dependency transition:

\[\adjustbox{max width=\linewidth}{$\displaystyle 
\text{physical persistence}
\longrightarrow
\text{active self-maintenance}
\longrightarrow
\text{organism-relative stake}
\longrightarrow
\text{graded agency and learning}
\longrightarrow
\text{possible perspectival and reflective organization}.
$}\]

The arrows mean enablement, not necessity, inevitability, or chronology.
Some capacities can arise through alternative routes. Artificial systems
may instantiate later functional properties without recapitulating
biological origins.

In the canonical sixty-four-stage Coherence Mandala, the biological
branch occupies Stages 11--15: Far-From-Equilibrium Organization,
Energy-Coupled Maintenance, Active Biological Self-Maintenance, Adaptive
Problem Solving, and Life as Coupled Autonomy. The sequence is a
conceptual and empirical architecture, not a claim that every lineage or
origin-of-life pathway must traverse one universal chronology.

Biology also marks a boundary important to the restored commitment
register. Life creates conditions under which attachment, care,
suffering, meaning, and love can later arise, but self-maintenance and
natural selection are not themselves love. \textbf{CU-COM14 --- Love as
the Ground and Aim of Conscious Life} begins where another being can
matter in itself; it is not a biological optimization rule.

\subsection{13.12 What this account does not
license}\label{what-this-account-does-not-license}

The scientific and spiritual account developed here does not license:

\begin{itemize}
\tightlist
\item
  treating evolution as a predetermined ascent;
\item
  inferring supernatural design from convergence;
\item
  equating life with consciousness;
\item
  assigning equal moral standing to every biological entity without
  further argument;
\item
  condemning death, predation, or ecological change as inherently
  incoherent;
\item
  romanticizing ``nature'' while ignoring suffering and conflict;
\item
  or using sacredness to block medicine, containment, food production,
  conservation intervention, or difficult ecological judgment.
\end{itemize}

Sacred regard requires better judgment, not refusal to judge.

\subsection{13.13 Final synthesis}\label{final-synthesis}

Abiogenesis, morphogenesis, and evolution share no single mechanism, but
they can be compared through a common organizational question: how do
constraints, relations, histories, and counterfactual capacities make
continued biological organization possible?

They are not expressions of one measured substance. They are different
transformations of organized possibility:

\begin{itemize}
\tightlist
\item
  abiogenesis concerns the emergence of networks that begin contributing
  to their own continuation;
\item
  morphogenesis concerns the construction and repair of viable form by
  already living networks;
\item
  evolution concerns changes in the heritable architectures,
  accessibility relations, boundaries, and conflicts through which such
  networks persist and reproduce.
\end{itemize}

Across all three, the central objects are constraints, relations,
histories, and counterfactual capacities.

The result is a biology of coherence without a universal coherence
number; a biology of purpose without supernatural force; a biology of
agency without projecting reflection onto every cell; an evolutionary
account of structured possibility without guaranteed progress; and a
spiritual interpretation that can affirm the sacredness of life without
disguising commitment as deduction.

\subsection{13.14 Conclusion: what biology teaches us about
persistence}\label{conclusion-what-biology-teaches-us-about-persistence}

A living being is not the matter it contains.

It is not one state held against change.

It is a way change is organized so that something can continue through
it.

The cell persists by exchanging its substance. The organism persists by
developing. The lineage persists by becoming different. Every biological
identity is therefore an achievement of continuity without sameness.

That achievement is finite. It requires energy, repair, boundary,
memory, and relation. It can fail from within or be overwhelmed from
without. It can preserve itself by harming what sustains it. It can
become rigid enough to survive today and lose the capacity to survive
tomorrow.

Coherence Universalism does not add a hidden force to this story. It
asks us to keep the full organization in view: the parts and their
relations, the organism and its environment, the present state and the
reachable future, the local success and the cost carried elsewhere.

Physics gives us structures that can persist. Biology gives us
structures whose persistence has become part of their own causal
organization.

In the descriptive register, this is a research program.

In the commitment register, it is an invitation to reverence, humility,
and care. We are not outside life looking in. Our bodies, thoughts,
technologies, institutions, and spiritual longings are late expressions
of the same vulnerable inheritance: matter held together for a time by
relations it did not create alone.

The task is not to worship persistence. It is to participate in forms of
persistence that remain adaptive, truthful about their dependencies, and
capable of sustaining other lives alongside their own.

That task continues in the companion papers on consciousness,
psychology, social dynamics, ethics, and artificial intelligence.

\begin{center}\rule{0.5\linewidth}{0.5pt}\end{center}

\section{Appendix A. Biology Principles and Claim-Governance
Summary}\label{appendix-a.-biology-principles-and-claim-governance-summary}

\subsection{A.1 Domain authority}\label{a.1-domain-authority}

Foundations assigns the transition from passive physical persistence to
active biological self-maintenance to CU Biology. CU-D3 is therefore a
gateway to CU-B1 rather than an independent universal law. The five
Biology principles below are the controlling CU formulations for this
domain. Their authority does not convert conditional empirical claims
into established universal laws.

\subsection{A.2 CU-B1 --- Active Biological
Self-Maintenance}\label{a.2-cu-b1-active-biological-self-maintenance}

\textbf{Status:} Domain Authority

\begin{quote}
A candidate living organization is not distinguished by static order or
persistence alone, but by a relatively closed network of materially
realized processes and constraints whose endogenous activity contributes
to maintaining, repairing, reproducing, or adaptively transforming the
conditions of that network's continuation. Degree and kind of autonomy
are relative to a boundary, environment, perturbation class, and
horizon.
\end{quote}

\subsection{A.3 CU-B2 --- Adaptive Problem Solving Beyond Fixed
Programs}\label{a.3-cu-b2-adaptive-problem-solving-beyond-fixed-programs}

\textbf{Status:} Domain Authority

\begin{quote}
Some living systems possess regulatory architectures that generalize
beyond fixed stimulus--response mappings: they detect deviations,
recruit alternative pathways, modify future responses from history, and
solve novel viability-relevant problems within a bounded adaptive reach.
Degree must be established through counterfactual flexibility, learning
controls, and comparison with passive or preprogrammed alternatives; it
does not by itself establish consciousness.
\end{quote}

\subsection{A.4 CU-B3 --- Multiscale Constraint and Local
Process}\label{a.4-cu-b3-multiscale-constraint-and-local-process}

\textbf{Status:} Domain Authority

\begin{quote}
Biological form and function arise through local physical processes
embedded in cross-scale boundary conditions, signaling networks,
mechanical relations, and histories. A higher-scale constraint is
explanatory when intervention on its realizing relations changes
lower-scale trajectories in a way that improves prediction or control
beyond local variables alone. Higher scale does not imply a new force,
and global stability does not imply a unique optimizer.
\end{quote}

\subsection{A.5 CU-B4 --- Structured Evolutionary Possibility and the
Stabilization of Regulatory
Architecture}\label{a.5-cu-b4-structured-evolutionary-possibility-and-the-stabilization-of-regulatory-architecture}

\textbf{Status:} Canonical--conditional empirical generalization under
the Version 0.10 canonical registry

\begin{quote}
Evolutionary processes explore a structured and nonuniform possibility
space shaped by mutation bias, genotype--phenotype relations,
developmental dynamics, physical feasibility, ecological interaction,
inheritance, and historical path dependence. Selection nonrandomly
changes the frequencies of heritable variants relative to local
reproductive consequences. Recurrent constraints can make some forms,
functions, and regulatory architectures repeatedly accessible,
contributing to convergence, while chance, drift, conflict, changing
environments, and historical contingency preserve multiple possible
trajectories. Evolution can stabilize, modify, simplify, or dismantle
coherent self-maintaining organization; it does not thereby maximize
coherence, complexity, consciousness, or moral value.
\end{quote}

\subsection{A.6 CU-B5 --- Biological Individuation and
Stake}\label{a.6-cu-b5-biological-individuation-and-stake}

\textbf{Status:} Domain Authority

\begin{quote}
A biological individual is a historically continuous, relatively bounded
organization whose components are causally integrated into maintaining,
repairing, reproducing, or adaptively transforming the conditions of
that organization. The organization thereby establishes
organism-relative stakes: conditions can be better or worse for it in
virtue of their effects on its continued self-maintenance. Individuality
is graded, multidimensional, and potentially nested; organism-relative
stake does not by itself establish consciousness or moral status.
\end{quote}

\subsection{A.7 Five claim-governance
rules}\label{a.7-five-claim-governance-rules}

Every substantive biological coherence claim must:

\begin{enumerate}
\def\labelenumi{\arabic{enumi}.}
\tightlist
\item
  name the candidate biological unit or boundary;
\item
  name the temporal horizon;
\item
  name the disturbance or perturbation class;
\item
  declare weights and protected floors whenever scalarizing a profile;
  and
\item
  prevent organism-relative stake from being promoted silently into
  phenomenal valence or moral value.
\end{enumerate}

A claim omitting these declarations is underspecified rather than
completed by assumption.

\subsection{A.8 Optimization-language
rule}\label{a.8-optimization-language-rule}

Biological explanation must distinguish:

\begin{itemize}
\tightlist
\item
  \textbf{designer optimization}, in which an actual designer chooses an
  objective and searches over policies;
\item
  \textbf{as-if control analysis}, in which an optimization expression
  summarizes behavior without identifying an internal solver; and
\item
  \textbf{evolutionary selection}, in which population dynamics alter
  the distribution of architectures.
\end{itemize}

Neither as-if analysis nor evolutionary filtering proves that an
organism represents or consciously chooses the objective.

\begin{center}\rule{0.5\linewidth}{0.5pt}\end{center}

\subsection{A.9 Mandala and Commitment
Cross-Reference}\label{a.9-mandala-and-commitment-cross-reference}

\textbf{CU-COM8 --- Reverence for Life and Biospheric Care} is the
official domain commitment anchor for Section 13. Revised
\textbf{CU-COM6} and \textbf{CU-T11} govern the language of evolutionary
grain without restoring a universal scalar or predetermined ascent.
\textbf{CU-COM14} names love as a later conscious and relational
commitment, not as a synonym for biological persistence or fitness.
CU-T5 supplies the sacredness interpretation; CU-G1, CU-G5, and CU-G7
govern distribution, participation, and retained option value.

\section{Appendix B. Major Revisions from
v0.09}\label{appendix-b.-major-revisions-from-v0.09}

This appendix records only the changes needed to orient readers. The
full patch, independent audit, audit response, and source-to-claim
ledger are preserved in the technical archive.

\begin{longtable}[]{@{}
  >{\raggedright\arraybackslash}p{(\columnwidth - 6\tabcolsep) * \real{0.2500}}
  >{\raggedright\arraybackslash}p{(\columnwidth - 6\tabcolsep) * \real{0.2500}}
  >{\raggedright\arraybackslash}p{(\columnwidth - 6\tabcolsep) * \real{0.2500}}
  >{\raggedright\arraybackslash}p{(\columnwidth - 6\tabcolsep) * \real{0.2500}}@{}}
\toprule\noalign{}
\begin{minipage}[b]{\linewidth}\raggedright
Topic
\end{minipage} & \begin{minipage}[b]{\linewidth}\raggedright
v0.09 formulation
\end{minipage} & \begin{minipage}[b]{\linewidth}\raggedright
v0.10 formulation
\end{minipage} & \begin{minipage}[b]{\linewidth}\raggedright
Disposition
\end{minipage} \\
\midrule\noalign{}
\endhead
\bottomrule\noalign{}
\endlastfoot
Biological measure & One universal scalar \(C(x)=A(x)I(x)\) & A declared
multidimensional profile; scalarization is purpose-relative & Universal
scalar retired \\
Dynamics & One framework-level coherence gradient & Explicit dynamics,
control sets, costs, constraints, and model-relative objectives &
Universal gradient retired \\
Thermodynamics & Life counters entropy; generic maintenance-cost law &
Open-system entropy accounting and process-specific physical bounds &
Replaced \\
Origins & One sharp abiogenesis phase transition & Coupled autonomy
transitions that may emerge in different orders & Universal threshold
retired \\
Recurrence & Recurrence dominance is sufficient for life & Passive
recurrent systems provide a counterexample & Sufficiency claim
rejected \\
Individuality & One organismal boundary & Graded, multidimensional,
nested, and purpose-relative individuation & Reconstructed \\
Stake & Viability and value tend to merge & Organism-relative functional
stake is separated from feeling and moral standing & Boundary made
explicit \\
Morphogenesis & Global coherence optimization & Interventionally
testable multiscale control and target-set models & Narrowed \\
Bioelectricity & Candidate universal integration carrier & Causal
coordination mechanism in supported systems and pathways & Restricted \\
Evolution & Coherence ascent or coherence navigation & Structured
evolutionary possibility: biased accessibility, development, ecology,
selection, drift, and history; no universal ascent & Reconstructed \\
Cancer & Coherence disintegration & Multilevel evolution, regulatory
decoupling, and scale conflict & Reconstructed \\
Aging and death & Progressive scalar decline; death as \(C=0\) &
Changing recoverable trajectories; intervention-relative irrecoverable
organismal failure & Reconstructed \\
Toy simulation & Validation of the framework & Pedagogical
non-entailment constructions with opposite coupling effects & Validation
claim withdrawn \\
Optimization language & Optimization often treated as mechanism &
Designer choice, as-if analysis, and selection are distinguished &
Governance rule added \\
Spiritual meaning & Sacredness sometimes appeared as biological
conclusion & Sacredness is an explicit commitment constrained by, but
not deduced from, biology & Register clarified \\
\end{longtable}

\begin{center}\rule{0.5\linewidth}{0.5pt}\end{center}

\section{Appendix C. Technical Supplement and
Reproducibility}\label{appendix-c.-technical-supplement-and-reproducibility}

\textbf{Released evidence package.} The current supplement is
\href{../supplements/04_Biology_Technical_Supplement_v0.10.md}{04\_Biology\_Technical\_Supplement\_v0.10.md}.
Its companion evidence directory is
\href{../evidence/04_Biology/README.md}{04\_Biology}. The supplement and
evidence directory expose the derivations, deterministic constructions,
scripts, machine-readable results, source ledger, and audit materials
referenced in this manuscript. They do not elevate a toy model,
simulation, or reconstructed check into target-domain empirical
validation.

The accompanying technical supplement contains:

\begin{itemize}
\tightlist
\item
  formal definitions and bounded propositions;
\item
  the two illustrative viability-control configurations;
\item
  the complete claim and evidence ledger;
\item
  the Biology principle crosswalk and the Version 0.10 canonical
  registry synchronization record;
\item
  the revision and independent-audit history;
\item
  the figure and experiment program;
\item
  the selective literature ledger;
\item
  the deterministic model, generated JSON, and environment declaration;
\item
  source hashes, editorial transformation records, and package QA.
\end{itemize}

The numerical constructions are pedagogical. They establish that the
displayed results follow from the displayed equations and that the
direction of a coupling effect can reverse across declared
configurations. They do not validate a universal biological law,
estimate an organismal optimum, or show that living systems obey the toy
equations.

\begin{center}\rule{0.5\linewidth}{0.5pt}\end{center}

\section{References}\label{references}

Aubin, J.-P. (1991). \emph{Viability theory}. Birkhäuser.

Aoki, S. K., Lillacci, G., Gupta, A., Baumschlager, A., Schweingruber,
D., \& Khammash, M. (2019). A universal biomolecular integral feedback
controller for robust perfect adaptation. \emph{Nature, 570}(7762),
533--537. \url{https://doi.org/10.1038/s41586-019-1321-1}

Blount, Z. D., Borland, C. Z., \& Lenski, R. E. (2008). Historical
contingency and the evolution of a key innovation in an experimental
population of \emph{Escherichia coli}. \emph{Proceedings of the National
Academy of Sciences, 105}(23), 7899--7906.
\url{https://doi.org/10.1073/pnas.0803151105}

Chernet, B. T., Adams, D. S., Lobikin, M., \& Levin, M. (2016). Use of
genetically encoded, light-gated ion translocators to control
tumorigenesis. \emph{Oncotarget, 7}(15), 19575--19588.
\url{https://doi.org/10.18632/oncotarget.8036}

Chernet, B. T., Fields, C., \& Levin, M. (2015). Long-range gap
junctional signaling controls oncogene-mediated tumorigenesis in
\emph{Xenopus laevis} embryos. \emph{Frontiers in Physiology, 5}, 519.
\url{https://doi.org/10.3389/fphys.2014.00519}

Ciriello, G., Magnani, L., Aitken, S. J., Akkari, L., Behjati, S.,
Hanahan, D., Landau, D. A., Lopez-Bigas, N., Lupiáñez, D. G., Marine,
J.-C., Martin-Villalba, A., Natoli, G., Obenauf, A. C., Oricchio, E.,
Scaffidi, P., Sottoriva, A., Swarbrick, A., Tonon, G., Vanharanta, S.,
\& Zuber, J. (2024). Cancer evolution: A multifaceted affair.
\emph{Cancer Discovery, 14}(1), 36--48.
\url{https://doi.org/10.1158/2159-8290.CD-23-0530}

Cortés-Poza, Y. (2026). A hybrid mathematical framework for
morphogenesis and regeneration. \emph{Journal of Mathematical Biology,
93}(3), Article 43. \url{https://doi.org/10.1007/s00285-026-02459-2}

Doulcier, G., Takacs, P., Hammerschmidt, K., \& Bourrat, P. (2024).
Stability of ecologically scaffolded traits during evolutionary
transitions in individuality. \emph{Nature Communications, 15}, Article
6566. \url{https://doi.org/10.1038/s41467-024-50625-1}

Dufour, Y. S., Fu, X., Hernandez-Nunez, L., \& Emonet, T. (2014). Limits
of feedback control in bacterial chemotaxis. \emph{PLOS Computational
Biology, 10}(6), e1003694.
\url{https://doi.org/10.1371/journal.pcbi.1003694}

Durant, F., Bischof, J., Fields, C., Morokuma, J., LaPalme, J., Hoi, A.,
\& Levin, M. (2019). The role of early bioelectric signals in the
regeneration of planarian anterior/posterior polarity. \emph{Biophysical
Journal, 116}(5), 948--961.
\url{https://doi.org/10.1016/j.bpj.2019.01.029}

Friston, K. (2013). Life as we know it. \emph{Journal of the Royal
Society Interface, 10}(86), 20130475.
\url{https://doi.org/10.1098/rsif.2013.0475}

Giessler, F., Verstraeten, W., Abele, T., Maurer, S. J., Monari, L., \&
Göpfrich, K. (2026). Growth, dissolution and segregation of genetically
encoded RNA droplets by ribozyme catalysis. \emph{Angewandte Chemie
International Edition, 65}(7), e202519002.
\url{https://doi.org/10.1002/anie.202519002}

Herrera-Álvarez, S., Patton, J. E. J., \& Thornton, J. W. (2025). The
structure of an ancient genotype--phenotype map shaped the functional
evolution of a protein family. \emph{Nature Ecology \& Evolution, 9},
1656--1669. \url{https://doi.org/10.1038/s41559-025-02777-6}

Hordijk, W., Kauffman, S. A., \& Steel, M. (2011). Required levels of
catalysis for emergence of autocatalytic sets in models of chemical
reaction systems. \emph{International Journal of Molecular Sciences,
12}(5), 3085--3101. \url{https://doi.org/10.3390/ijms12053085}

Huang, S., Eichler, G., Bar-Yam, Y., \& Ingber, D. E. (2005). Cell fates
as high-dimensional attractor states of a complex gene regulatory
network. \emph{Physical Review Letters, 94}(12), 128701.
\url{https://doi.org/10.1103/PhysRevLett.94.128701}

Kolchinsky, A. (2024). Thermodynamic dissipation does not bound
replicator growth and decay rates. \emph{arXiv}.
\url{https://doi.org/10.48550/arXiv.2404.01130}

Kriegman, S., Blackiston, D., Levin, M., \& Bongard, J. (2020). A
scalable pipeline for designing reconfigurable organisms.
\emph{Proceedings of the National Academy of Sciences, 117}(4),
1853--1859. \url{https://doi.org/10.1073/pnas.1910837117}

Kriegman, S., Blackiston, D., Levin, M., \& Bongard, J. (2021).
Kinematic self-replication in reconfigurable organisms.
\emph{Proceedings of the National Academy of Sciences, 118}(49),
e2112672118. \url{https://doi.org/10.1073/pnas.2112672118}

Losos, J. B., Jackman, T. R., Larson, A., de Queiroz, K., \&
Rodríguez-Schettino, L. (1998). Contingency and determinism in
replicated adaptive radiations of island lizards. \emph{Science,
279}(5359), 2115--2118.
\url{https://doi.org/10.1126/science.279.5359.2115}

Maynard Smith, J., \& Szathmáry, E. (1995). \emph{The major transitions
in evolution}. W. H. Freeman.

Montévil, M., \& Mossio, M. (2015). Biological organisation as closure
of constraints. \emph{Journal of Theoretical Biology, 372}, 179--191.
\url{https://doi.org/10.1016/j.jtbi.2015.02.029}

Mossio, M., Saborido, C., \& Moreno, A. (2009). An organizational
account of biological functions. \emph{The British Journal for the
Philosophy of Science, 60}(4), 813--841.
\url{https://doi.org/10.1093/bjps/axp036}

Rajan, D. H., Albright, A., Kim, H., Diaz, U., Hudnall, Y., Steube, N.,
Dey, G., Liu, T., \& Marshall, W. F. (2026). Molecular pathways for
learning in the single-cell \emph{Stentor coeruleus}. \emph{Current
Biology, 36}(9), 2367--2381.e9.
\url{https://doi.org/10.1016/j.cub.2026.03.080}

Rader, G. K. D. (2026). \emph{Coherence Universalism:
Physics---Coherence, Accessibility, and the Ordered Emergence of
Physical Description} (Version 0.10). Heaven≡Earth Foundation.

Rader, G. K. D. (2026). \emph{Coherence Universalism: Metaphysics and
Epistemology---Coherence Under Constraint, Explanatory Order, and
Warrant Across Registers} (Version 0.10). Heaven≡Earth Foundation.

Rader, G. K. D. (2026). \emph{Coherence Universalism:
Foundations---Principles, Domain Gateways, Commitments, and Constraints
on Responsible Application} (Version 0.10). Heaven≡Earth Foundation.

Ruiz-Mirazo, K., Peretó, J., \& Moreno, A. (2004). A universal
definition of life: Autonomy and open-ended evolution. \emph{Origins of
Life and Evolution of the Biosphere, 34}, 323--346.
\url{https://doi.org/10.1023/B:ORIG.0000016440.53346.dc}

Schaper, S., \& Louis, A. A. (2014). The arrival of the frequent: How
bias in genotype--phenotype maps can steer populations to local optima.
\emph{PLOS ONE, 9}(2), e86635.
\url{https://doi.org/10.1371/journal.pone.0086635}

Schreiber, T. (2000). Measuring information transfer. \emph{Physical
Review Letters, 85}(2), 461--464.
\url{https://doi.org/10.1103/PhysRevLett.85.461}

Shi, K., Song, P., Li, J., Qiao, Y., \& Lin, Y. (2025). Light-fuelled
growth dynamics and structural transition of synthetic protocells.
\emph{Nature Synthesis, 4}, 1359--1368.
\url{https://doi.org/10.1038/s44160-025-00844-1}

Shirt-Ediss, B., Ferrero-Fernández, A., De Martino, D., Bich, L.,
Moreno, A., \& Ruiz-Mirazo, K. (2025). Modelling the prebiotic origins
of regulation and agency in evolving protocell ecologies.
\emph{Philosophical Transactions of the Royal Society B: Biological
Sciences, 380}(1936), 20240287.
\url{https://doi.org/10.1098/rstb.2024.0287}

Simões, P., Fragata, I., Santos, J., Santos, M. A., Santos, M., Rose, M.
R., \& Matos, M. (2019). How phenotypic convergence arises in
experimental evolution. \emph{Evolution, 73}(9), 1839--1849.
\url{https://doi.org/10.1111/evo.13806}

Storz, J. F., Natarajan, C., Signore, A. V., Witt, C. C., McCandlish, D.
M., \& Stoltzfus, A. (2019). The role of mutation bias in adaptive
molecular evolution: Insights from convergent changes in protein
function. \emph{Philosophical Transactions of the Royal Society B,
374}(1777), 20180238. \url{https://doi.org/10.1098/rstb.2018.0238}

Travisano, M., Mongold, J. A., Bennett, A. F., \& Lenski, R. E. (1995).
Experimental tests of the roles of adaptation, chance, and history in
evolution. \emph{Science, 267}(5194), 87--90.
\url{https://doi.org/10.1126/science.7809610}

Varela, F. G., Maturana, H. R., \& Uribe, R. (1974). Autopoiesis: The
organization of living systems, its characterization and a model.
\emph{BioSystems, 5}(4), 187--196.
\url{https://doi.org/10.1016/0303-2647(74)90031-8}

Walker, S. I., \& Davies, P. C. W. (2013). The algorithmic origins of
life. \emph{Journal of the Royal Society Interface, 10}(79), 20120869.
\url{https://doi.org/10.1098/rsif.2012.0869}

Rohner, P. T., Hu, Y., \& Moczek, A. P. (2022). Developmental bias in
the evolution and plasticity of beetle horn shape. \emph{Proceedings of
the Royal Society B, 289}(1983), 20221441.
\url{https://doi.org/10.1098/rspb.2022.1441}

West, S. A., Fisher, R. M., Gardner, A., \& Kiers, E. T. (2015). Major
evolutionary transitions in individuality. \emph{Proceedings of the
National Academy of Sciences, 112}(33), 10112--10119.
\url{https://doi.org/10.1073/pnas.1421402112}

\end{document}
